\documentclass[11pt,a4paper,oneside]{book}

\usepackage{iftex}
\ifPDFTeX
  \usepackage[utf8]{inputenc}
  \usepackage[T1]{fontenc}
\else
  \usepackage{fontspec}
\fi

\usepackage[margin=2.0cm, headheight=24pt, footskip=26pt]{geometry}
\usepackage{xcolor}
\usepackage{tcolorbox}
\tcbuselibrary{breakable,skins}
\usepackage{booktabs}
\usepackage{longtable}
\usepackage{array}
\usepackage{tabularx}
\usepackage{listings}
\usepackage{fancyhdr}
\usepackage{titlesec}
\usepackage{enumitem}
\usepackage{parskip}
\usepackage{microtype}
\usepackage{amsmath}
\PassOptionsToPackage{obeyspaces,hyphens}{url}
\usepackage{xurl}
\usepackage{hyperref}

% ==============================================================================
% PALET WARNA & IDENTITAS VISUAL CSIP
% ==============================================================================
\definecolor{csipDark}{HTML}{092B2C}
\definecolor{csipTeal}{HTML}{0F3D3E}
\definecolor{csipAccent}{HTML}{17A394}
\definecolor{csipAmber}{HTML}{D97706}
\definecolor{csipBlue}{HTML}{0369A1}
\definecolor{csipCrimson}{HTML}{B91C1C}
\definecolor{csipSlate}{HTML}{1E293B}
\definecolor{csipMuted}{HTML}{475569}
\definecolor{csipLightBg}{HTML}{F8FAFC}
\definecolor{csipCodeBg}{HTML}{F1F5F9}

\hypersetup{
  colorlinks=true,
  linkcolor=csipTeal,
  urlcolor=csipBlue,
  pdftitle={Buku Induk Rekayasa, Sains Data \& Skrip Presentasi Master — CSIP v1.1.0},
  pdfauthor={Tim Rekayasa \& Sains Data CSIP}
}

% Makro aman untuk penulisan identifier kode (ukuran nyaman 10pt & putus baris otomatis)
\newcommand{\code}[1]{{\small\nolinkurl{#1}}}
\setlength{\tabcolsep}{4.5pt}
\tolerance=1200
\emergencystretch=3em
\AtBeginEnvironment{tcolorbox}{\sloppy}

% Spasi nyaman untuk buku referensi panjang (>50 halaman) & naskah bicara
\linespread{1.22}
\setlength{\parindent}{0pt}
\setlength{\parskip}{6pt}
\setlist[itemize]{topsep=3pt, itemsep=3pt, parsep=2pt, leftmargin=1.4em}
\setlist[enumerate]{topsep=3pt, itemsep=3pt, parsep=2pt, leftmargin=1.6em}

% Gaya blok kode
\lstset{
  basicstyle=\ttfamily\small,
  backgroundcolor=\color{csipCodeBg},
  frame=leftline,
  framerule=2.5pt,
  rulecolor=\color{csipTeal},
  breaklines=true,
  breakatwhitespace=true,
  columns=flexible,
  keepspaces=true,
  showstringspaces=false,
  xleftmargin=8pt,
  xrightmargin=4pt,
  aboveskip=8pt,
  belowskip=8pt
}

% Header & footer profesional (konsisten 10pt, bebas tabrakan)
\pagestyle{fancy}
\fancyhf{}
\fancyhead[L]{\small\color{csipMuted}\textbf{CSIP v1.1.0}}
\fancyhead[R]{\small\color{csipMuted}\nouppercase{\leftmark}}
\fancyfoot[L]{\small\color{csipMuted}Core Strategy \& Intelligence Platform (\code{csip})}
\fancyfoot[R]{\small\color{csipMuted}Halaman \thepage}
\renewcommand{\headrulewidth}{0.4pt}
\renewcommand{\footrulewidth}{0.4pt}

% Gaya judul Bagian & Bab yang tegas
\newcommand{\csippartbox}[1]{\colorbox{csipTeal}{\parbox{0.92\textwidth}{\centering\vspace{12pt}\color{white}\Large #1\vspace{12pt}}}}
\titleformat{\part}[display]
  {\centering\bfseries\Huge\color{csipDark}}
  {\vspace*{1.5cm}\Large\color{csipAccent}BAGIAN \thepart}
  {14pt}
  {\csippartbox}

\titleformat{\chapter}[display]
  {\bfseries\color{csipTeal}}
  {\large\color{csipAccent}BAB \thechapter}
  {4pt}
  {\LARGE}
  [\vspace{4pt}{\color{csipTeal}\titlerule[1.5pt]}]
\titlespacing*{\chapter}{0pt}{-10pt}{16pt}

\titleformat{\section}{\large\bfseries\color{csipDark}}{\thesection}{0.7em}{}
\titleformat{\subsection}{\normalsize\bfseries\color{csipBlue}}{\thesubsection}{0.7em}{}

% Penamaan daftar isi berbahasa Indonesia
\renewcommand{\contentsname}{Daftar Isi Lengkap}
\renewcommand{\listtablename}{Daftar Tabel Utama}
\renewcommand{\partname}{Bagian}
\renewcommand{\chaptername}{Bab}

% ==============================================================================
% KOTAK PANGGILAN (CALLOUT BOXES) UNTUK NON-TEKNIS, TEKNIS & SKRIP PRESENTASI
% ==============================================================================
\newtcolorbox{bisnisbox}[1][Analogi Bisnis \& Penjelasan Mudah]{
  breakable, enhanced, before upper={\sloppy},
  colback=blue!3!white, colframe=csipBlue, colbacktitle=csipBlue,
  fonttitle=\bfseries\small\color{white}, title={#1},
  left=8pt, right=8pt, top=6pt, bottom=6pt, boxrule=0.8pt
}

\newtcolorbox{teknisbox}[1][Bedah Teknis \& Logika Rekayasa Data]{
  breakable, enhanced, before upper={\sloppy},
  colback=csipLightBg, colframe=csipSlate, colbacktitle=csipSlate,
  fonttitle=\bfseries\small\color{white}, title={#1},
  left=8pt, right=8pt, top=6pt, bottom=6pt, boxrule=0.8pt
}

\newtcolorbox{kejujuranbox}[1][Disiplin Ilmiah \& Batasan Klaim]{
  breakable, enhanced, before upper={\sloppy},
  colback=csipAmber!7!white, colframe=csipAmber, colbacktitle=csipAmber,
  fonttitle=\bfseries\small\color{white}, title={#1},
  left=8pt, right=8pt, top=6pt, bottom=6pt, boxrule=0.8pt
}

\newtcolorbox{stagebox}[1][Panduan Layar / Slide Saat Presentasi]{
  breakable, enhanced, before upper={\sloppy},
  colback=cyan!4!white, colframe=csipBlue, colbacktitle=csipBlue,
  fonttitle=\bfseries\small\color{white}, title={#1},
  left=8pt, right=8pt, top=5pt, bottom=5pt, boxrule=0.8pt
}

\newtcolorbox{scriptbox}[1][Naskah Bicara Lengkap (Siap Baca Kata-demi-Kata)]{
  breakable, enhanced, before upper={\sloppy},
  colback=csipTeal!4!white, colframe=csipTeal, colbacktitle=csipTeal,
  fonttitle=\bfseries\small\color{white}, title={#1},
  left=9pt, right=9pt, top=7pt, bottom=7pt, boxrule=1pt
}

\newtcolorbox{tipbox}[1][Tips Penyampaian \& Hal yang Jangan Salah Ucap]{
  breakable, enhanced, before upper={\sloppy},
  colback=red!2!white, colframe=csipCrimson, colbacktitle=csipCrimson,
  fonttitle=\bfseries\small\color{white}, title={#1},
  left=8pt, right=8pt, top=5pt, bottom=5pt, boxrule=0.8pt
}

% ==============================================================================
% MULAI DOKUMEN
% ==============================================================================
\begin{document}

\frontmatter

% HALAMAN SAMPUL
\begin{titlepage}
\thispagestyle{empty}
\begin{tcolorbox}[colback=csipDark, colframe=csipDark, arc=0pt, left=18pt, right=18pt, top=24pt, bottom=24pt]
{\color{csipAccent}\bfseries\large DOKUMEN ACUAN INDUK PROYEK \& PANDUAN PRESENTASI EKSEKUTIF-TEKNIS}\\[10pt]
{\color{white}\bfseries\fontsize{22}{28}\selectfont BUKU INDUK REKAYASA,\\[3pt] SAINS DATA \& SKRIP PRESENTASI}\\[10pt]
{\color{csipAccent!35!white}\Large Core Strategy \& Intelligence Platform (CSIP v1.1.0)}
\end{tcolorbox}

\vspace{1.0cm}
\large
\textbf{Cakupan Lengkap Dokumen:}
\begin{itemize}
  \item \textbf{Bagian I --- Fondasi Bisnis, Data \& Rekayasa Fitur (Bab 1--4):} Masalah operasional e-commerce dari titik nol, arsitektur 5-lapis, pembersihan 9 tabel relasional Olist (\code{99.441 -> 95.824} baris), pencegahan \emph{Cartesian Fan-Out}, kontrak data 4-tingkat, isolasi waktu keputusan ($T_0, T_1, T_2$), dan embargo 90 hari.
  \item \textbf{Bagian II --- Sains Data, Pemodelan, Hipotesis \& Kejujuran Empiris (Bab 5--8):} Tangga baseline B0--B4, cara kerja seluruh algoritma (LightGBM, Logistic/Ridge, Bi-LSTM CPU, K-Means, Kaplan-Meier, VAR Ridge), papan peringkat empiris lengkap beserta penolakan model gagal, uji hipotesis H1--H6, ketahanan wilayah, kalibrasi Platt Scaling, atribusi SHAP, dan simulasi kebijakan S0--S5.
  \item \textbf{Bagian III --- Decision Engine, CSV Studio, MLOps \& Pengujian (Bab 9--12):} Mesin rekomendasi 15 field tervalidasi, kalibrasi kuota kapasitas CS, rancangan uji A/B holdout 80/20, arsitektur Universal CSV Studio (\code{ml_engine.py} \& \code{semantic_profiler.py}), detektor drift PSI/KS/Wasserstein/JSD, siklus Champion-Challenger, 7 rute API FastAPI, kokpit Streamlit CSIP Studio, 216 pengujian \code{pytest} (39 berkas), 12 gerbang kualitas, dan pemetaan SKKNI Level 7.
  \item \textbf{Bagian IV --- Paket Presentasi Master Siap Manggung (Bab 13--16):} Strategi multi-audiens, manajemen waktu (5 menit, 15 menit, dan 45 menit), \textbf{skrip bicara kata-demi-kata sangat rinci}, panduan klik \emph{Live Demo} langkah-demi-langkah, serta \textbf{Ensiklopedia 25 Pertanyaan Tanya-Jawab} tingkat manajemen dan penguji Data Science.
\end{itemize}

\vfill
\begin{tcolorbox}[colback=csipLightBg, colframe=csipTeal, left=12pt, right=12pt, top=10pt, bottom=10pt]
\normalsize
\begin{tabular}{@{}l p{10.6cm}@{}}
\textbf{Versi Platform:} & \code{1.1.0} (\code{pyproject.toml} \& \code{FINAL_RELEASE_MANIFEST.json}) \\
\textbf{Kompatibilitas Perangkat:} & Laptop Standar RAM $\le$ 8 GB, CPU-Only (Tanpa GPU Wajib) \\
\textbf{Keamanan Data:} & Pemrosesan Lokal 100\% (Tanpa API LLM Cloud Eksternal) \\
\textbf{Sasaran Pembaca:} & Pimpinan Non-Teknis, Manajer CS/Logistik, Data Scientist, ML/Data Engineer, serta Asesor Kompetensi SKKNI No. 299 Tahun 2020 (KKNI Level 7) \\
\end{tabular}
\end{tcolorbox}
\end{titlepage}

% KATA PENGANTAR & PETUNJUK MEMBACA
\chapter*{Kata Pengantar \& Cara Menggunakan Buku Ini}
\addcontentsline{toc}{chapter}{Kata Pengantar \& Cara Menggunakan Buku Ini}

Buku ini ditulis dengan satu prinsip utama: \textbf{tidak seorang pun boleh tersesat saat memahami atau mempresentasikan proyek \emph{Core Strategy \& Intelligence Platform} (CSIP v1.1.0)}, baik Anda seorang pimpinan bisnis non-teknis yang berfokus pada efisiensi biaya dan kepuasan pelanggan, maupun seorang praktisi atau penguji Data Science yang ingin membedah logika kode, alur pemodelan, hingga bukti pengujian produksinya.

Banyak dokumentasi teknis gagal karena memakai bahasa yang terlalu kaku untuk manajemen atau membuang rumus matematika tingkat lanjut yang jauh melampaui ruang lingkup kerja Data Scientist sehari-hari. Sebaliknya, banyak skrip presentasi gagal karena penulisnya menghafal kalimat pemasaran tanpa benar-benar mengerti angka apa yang ada di balik layar saat penguji bertanya mendadak.

Oleh karena itu, seluruh isi buku ini menyatukan \textbf{Buku Acuan Teknis-Bisnis} dan \textbf{Skrip Presentasi Lengkap} ke dalam satu sistem belajar berlapis:

\begin{itemize}
  \item \textbf{Kotak Biru (Analogi Bisnis \& Penjelasan Mudah):} Menjelaskan konsep rumit dengan bahasa sehari-hari yang dapat langsung dipahami orang awam dalam 20 detik.
  \item \textbf{Kotak Abu-abu (Bedah Teknis \& Logika Rekayasa):} Membahas struktur modul Python, parameter algoritma, kontrak skema data, dan aturan validasi sesuai isi aktual repositori.
  \item \textbf{Kotak Kuning (Disiplin Ilmiah \& Batasan Klaim):} Mengingatkan batas kejujuran model agar pembicara tidak pernah mengeluarkan klaim berlebihan di depan penguji.
  \item \textbf{Bagian IV (Skrip Presentasi Kata-demi-Kata \& Tanya-Jawab):} Berisi naskah siap baca yang lengkap dengan petunjuk slide, gerakan kursor saat \emph{Live Demo}, serta jawaban teruji atas 25 pertanyaan kritis.
\end{itemize}

Seluruh angka metrik empiris, nama modul, nama kelas, batas uji, dan status kualitas yang tercantum di dalam buku ini diambil langsung dari berkas aktual repositori CSIP v1.1.0 pada saat kompilasi, antara lain \code{locked_test_reevaluation.json}, \code{ml_leaderboard.csv}, \code{model_explainability.json}, \code{hypothesis_evaluation.json}, \code{clustering_evaluation.json}, \code{model_slice_analysis.json}, \code{scenario_simulation_results.json}, \code{world_model_evaluation.json}, serta \code{resource_benchmark.json}.

\tableofcontents
\listoftables

\mainmatter

% ==============================================================================
% BAGIAN I — FONDASI BISNIS, DATA & REKAYASA FITUR
% ==============================================================================
\part{Fondasi Bisnis, Data \& Rekayasa Fitur (Dari Titik Nol)}

\chapter{Gambaran Besar Proyek dari Titik Nol: Masalah Bisnis, Visi Produk \& Strategi Dua Lapis}
\chaptermark{Gambaran Besar \& Strategi Dua Lapis}

\section{Titik Awal: Mengapa Platform Seperti CSIP Dibutuhkan?}

Setiap proyek Data Science yang bernilai tinggi selalu berangkat dari \textbf{masalah operasional nyata}, bukan dari keinginan pamer algoritma. Untuk memahami mengapa CSIP dibangun dari titik nol, mari kita tempatkan diri kita sebagai pengelola sebuah ekosistem \emph{marketplace e-commerce} skala nasional.

Di dalam ekosistem transaksi daring yang memproses hampir \textbf{100.000 pesanan} melintasi wilayah geografis raksasa, ribuan paket berpindah setiap minggu dari ratusan penjual menuju pembeli melalui jaringan kurir berlapis. Pada kondisi normal, sebagian besar barang tiba tepat waktu dan pembeli merasa puas. Namun di balik angka kelancaran tersebut, selalu muncul segelintir transaksi yang bermasalah: paket tersangkut di gudang transit, estimasi tanggal kirim meleset beberapa hari, atau pembeli merasa biaya kirim yang mereka bayar terlalu mahal dibandingkan harga produknya.

Akibatnya berantai bagi perusahaan:
\begin{enumerate}
  \item Pembeli yang kecewa memberikan ulasan bintang 1 atau bintang 2 (\emph{negative review}).
  \item Reputasi penjual dan kepercayaan terhadap platform menurun.
  \item Tim \emph{Customer Service} (CS) baru mengetahui adanya krisis setelah ulasan buruk telanjur terbit di halaman publik --- ketika kesempatan menyelamatkan pengalaman pelanggan sudah hampir tertutup.
  \item Jika perusahaan ingin mengirim voucher permintaan maaf atau menghubungi pembeli satu per satu, jumlah tenaga manusia tidak mencukupi. Menghubungi ribuan pesanan secara acak jelas memboroskan biaya, sedangkan tidak bertindak sama sekali berisiko kehilangan pelanggan.
\end{enumerate}

\begin{bisnisbox}[Analogi Sederhana: Dokter Triase IGD untuk Pesanan E-Commerce]
Bayangkan tim CS dan tim Logistik seperti unit IGD rumah sakit dengan jumlah dokter yang terbatas. Mereka tidak bisa memeriksa kondisi ribuan pasien sehat satu per satu sepanjang hari. Mereka membutuhkan \textbf{sistem triase cerdas} yang mampu memberi tahu secara terukur:
\begin{itemize}
  \item Pasien mana yang kemungkinan besar sedang mengalami gejala kritis paling parah (\textbf{Probabilitas Risiko}).
  \item Mengapa ia berisiko (\textbf{Faktor Pemicu Utama}, misalnya telat berapa hari dan nilai belanjanya berapa).
  \item Siapa yang harus menangani terlebih dahulu sesuai kapasitas jumlah dokter yang bertugas hari ini (\textbf{Daftar Prioritas Terukur Kuota}).
\end{itemize}
\end{bisnisbox}

\section{Dua Jebakan Klasik Proyek Data Science yang Dihindari CSIP}

Selama pengembangan proyek analitik di industri maupun lingkungan akademik, sering ditemui dua kutub kegagalan yang sama-sama ingin dipatahkan oleh CSIP:

\begin{enumerate}
  \item \textbf{Jebakan Pertama --- Proyek Notebook Akademik yang Bocor Data:} Banyak eksperimen berhenti sebagai berkas notebook statis yang hanya memamerkan skor ROC-AUC 0,95 atau akurasi 98\%. Ketika diaudit lebih teliti, angka tinggi tersebut muncul akibat pemecahan acak (\emph{random train-test split}) atau karena memasukkan informasi masa depan (misalnya sudah tahu paket telat 6 hari padahal klaimnya memprediksi saat pembeli baru membayar di kasir). Begitu diuji di data masa depan nyata, model tersebut langsung lumpuh.
  \item \textbf{Jebakan Kedua --- Aplikasi Dashboard Kosmetik / Wrapper Chatbot:} Aplikasi terlihat penuh warna, namun angka metrik di dalamnya diketik manual, tidak memakai kontrak validasi data, tidak mengukur apakah probabilitas model sudah terkalibrasi, serta tidak menghasilkan keluaran aksi apa yang harus dilakukan tim bisnis.
\end{enumerate}

\textbf{CSIP v1.1.0} didesain sejak baris kode pertama sebagai sistem kecerdasan terpadu yang menjawab \textbf{lima tingkatan pertanyaan bisnis} dalam satu alur berkesinambungan:

\begin{longtable}{@{} >{\raggedright\arraybackslash}p{2.6cm} >{\raggedright\arraybackslash}p{4.8cm} >{\raggedright\arraybackslash}p{8.4cm} @{}}
\caption{Strategi Dua Lapis Arsitektur CSIP (Kedalaman Vertikal vs Keluasan Universal)} \\
\toprule
\textbf{Tingkatan Analitik} & \textbf{Pertanyaan Bisnis Utama} & \textbf{Jawaban Nyata di dalam Sistem CSIP} \\
\midrule
\endfirsthead
\toprule
\textbf{Tingkatan Analitik} & \textbf{Pertanyaan Bisnis Utama} & \textbf{Jawaban Nyata di dalam Sistem CSIP} \\
\midrule
\endhead
\bottomrule
\endfoot
\textbf{1. Deskriptif} & \emph{Apa yang sedang terjadi di operasi kita?} & Profil kualitas 9 tabel Olist, pemetaan funnel populasi 99.441 transaksi, lama tunggu rata-rata, rasio ongkir, dan distribusi ulasan bintang 1--5. \\
\textbf{2. Diagnostik} & \emph{Mengapa ketidakpuasan tersebut cenderung terjadi?} & Pengujian 6 hipotesis preregistered (\code{H1}--\code{H6}), atribusi fitur SHAP pada 800 sampel, serta segmentasi perilaku pelanggan K-Means. \\
\textbf{3. Prediktif} & \emph{Pesanan mana yang berisiko terlambat atau menerima ulasan buruk?} & Model LightGBM terpisah per titik waktu ($T_0$ checkout, $T_1$ kurir, $T_2$ pasca-kirim), analisis sentimen teks Bi-LSTM, dan model tren makro mingguan VAR. \\
\textbf{4. Preskriptif} & \emph{Tindakan apa yang harus diambil tim hari ini dengan kuota terbatas?} & \code{BusinessDecisionEngine} yang mengubah skor probabilitas + nilai keranjang menjadi prioritas (\code{CRITICAL/HIGH/MEDIUM/LOW}) terpotong kuota kapasitas harian CS. \\
\textbf{5. Evaluatif} & \emph{Apakah sistem masih akurat bulan depan dan bagaimana membuktikan dampak intervensinya?} & Rancangan uji A/B holdout 80/20, observatorium pergeseran distribusi (PSI, KS, Wasserstein, JSD), dan gerbang promosi Champion-Challenger. \\
\end{longtable}

\section{Strategi Arsitektur Dua Lapis: Kedalaman Vertikal + Keluasan Lintas Industri}

Salah satu pertanyaan strategis terpenting yang sering diajukan pengambil keputusan saat melihat sistem berbasis dataset Olist adalah: \emph{``Apakah platform ini hanya berguna untuk satu dataset toko online Brasil saja? Bagaimana jika tim kami ingin memeriksa file CSV data medis, SDM, atau kendali mutu pabrik kami sendiri?''}

Memaksa file CSV bidang lain masuk ke aturan logistik Olist tentu akan merusak kevalidan ilmiah. Sebaliknya, hanya membuat aplikasi pembaca CSV generik tanpa studi kasus mendalam membuat proyek terlihat dangkal. Oleh sebab itu, CSIP mengambil keputusan desain \textbf{Strategi Arsitektur Dua Lapis (\emph{Depth vs. Breadth})}:

\begin{itemize}
  \item \textbf{Lapis 1 --- Pembuktian Mendalam Vertikal (\emph{Vertical Domain Proof}):} Seluruh paket inti di \code{src/csip/}, model terlatih di \code{artifacts/models/}, protokol preregistered di \code{config/analysis_plan.yaml}, layanan API FastAPI, serta 14 notebook riset membedah tuntas skema relasional 9 tabel Olist dengan disiplin metode Data Science tingkat produksi.
  \item \textbf{Lapis 2 --- Mesin Eksplorasi Lintas Industri (\emph{Horizontal Studio Proof}):} Antarmuka interaktif utama \textbf{CSIP Studio} (\code{app/streamlit_app.py}) yang ditenagai \code{ml_engine.py} dan \code{semantic_profiler.py} memungkinkan pengguna mengunggah file \code{.csv} atau \code{.xlsx}/\code{.xls} dari domain mana saja --- maupun menguji 3 kasus bawaan (E-Commerce Olist, Kesehatan Medis, dan Mutu Produk) --- tanpa mengotori artefak produksi Olist. Mesin secara otomatis memetakan kosakata 11 sektor industri, mendeteksi risiko kebocoran kolom tanggal/ID, memilih jalur regresi vs klasifikasi biner, melakukan \emph{Cross-Validation} bebas kebocoran, dan menyajikan simulator interaktif.
\end{itemize}

\section{Peta Alur Kecerdasan Terpadu dari Hulu ke Hilir}

Berikut aliran kerja lengkap tujuh langkah yang menjadi tulang punggung sistem CSIP:

\begin{lstlisting}
[Langkah 1: Pertanyaan Bisnis & Protokol Preregistered (config/analysis_plan.yaml)]
        |
        v
[Langkah 2: Kontrak Data 4-Tingkat & Pencatatan Hash SHA-256 (contracts.py & lineage.py)]
        |-----(Jika Cacat Fatal: CRITICAL_FAILURE)----> [Pipeline Berhenti Aman]
        v
[Langkah 3: Pembersihan Deterministik & Konsolidasi ke 1 Baris per order_id (cleaner.py)]
        |
        v
[Langkah 4: Isolasi Fitur Waktu (Pos T0 Checkout | Pos T1 Kurir | Pos T2 Pasca-Kirim)]
        |
        v
[Langkah 5: Tangga Model B0-B4 + Kalibrasi Platt Scaling + Atribusi SHAP + Ketahanan Wilayah]
        |
        v
[Langkah 6: Business Decision Engine (Prioritas + Kuota Kapasitas + Desain Uji A/B 80/20)]
        |
        v
[Langkah 7: Serving Terpisah (FastAPI + CSIP Studio) + Observatorium Drift + Retraining Gate]
\end{lstlisting}

% ------------------------------------------------------------------------------
\chapter{Arsitektur Sistem 5-Lapis, Pilihan Teknologi \& Jejak Sumber Daya}
\chaptermark{Arsitektur 5-Lapis \& Sumber Daya}

\section{Keputusan Batasan Lingkungan Eksekusi (Target Laptop RAM $\le$ 8 GB)}

Sejak tahap perancangan, CSIP menetapkan kendala rekayasa yang realistis: seluruh pipeline pembersihan data, pelatihan model utama, evaluasi SHAP, layanan API FastAPI, dan kokpit Streamlit harus mampu dijalankan di \textbf{laptop standar dengan RAM $\le$ 8 GB tanpa kartu grafis (GPU)} serta menjaga \textbf{privasi data lokal penuh} (nol panggilan ke API LLM cloud eksternal atau pengiriman data pelanggan ke internet).

Batasan ini secara langsung membentuk lima keputusan teknologi utama CSIP:

\begin{longtable}{@{} >{\raggedright\arraybackslash}p{3.4cm} >{\raggedright\arraybackslash}p{3.8cm} >{\raggedright\arraybackslash}p{8.6cm} @{}}
\caption{Pustaka Teknologi Inti CSIP v1.1.0 dan Alasan Penguncian Lingkungan} \\
\toprule
\textbf{Keputusan Teknis} & \textbf{Pilihan yang Diambil} & \textbf{Mengapa Dipilih \& Apa yang Ditolak} \\
\midrule
\endfirsthead
\toprule
\textbf{Keputusan Teknis} & \textbf{Pilihan yang Diambil} & \textbf{Mengapa Dipilih \& Apa yang Ditolak} \\
\midrule
\endhead
\bottomrule
\endfoot
\textbf{1. Format Data Internal} & \textbf{Apache Parquet} (\code{float32} / \code{int32}) & File \code{order_level.parquet} hanya berukuran $\sim$18,43 MB dibanding CSV mentah gabungan $>$120 MB. Memuat jauh lebih cepat dan tipe kolom tetap terjaga presisi, alih-alih membaca CSV mentah berulang yang boros RAM. \\
\textbf{2. Algoritma Tabular Utama} & \textbf{LightGBM} (\code{LGBMClassifier} / \code{LGBMRegressor}) & Pada data transaksi 10--15 fitur berdistribusi miring, LightGBM cepat di CPU, ukuran model hanya 433--520 KB, inferensi 1.000 baris $<$15 ms, dan mudah dijelaskan lewat SHAP. Deep learning tabular (\code{TabularEmbeddingMLP}) diuji di tangga B4 tetapi kalah efisien untuk produksi CPU. \\
\textbf{3. Pemrosesan Teks Ulasan} & \textbf{PyTorch Bi-LSTM} lokal ringan ($\sim$1,03 MB) & Menggunakan model bahasa besar (LLM) atau memanggil API cloud melanggar target RAM 8 GB dan privasi lokal. Bi-LSTM berkapasitas 5.000 kosakata berjalan mulus di CPU dengan \code{Macro-F1} 0,8420. \\
\textbf{4. Pemodelan Dinamika Makro} & \textbf{VAR(1) Ridge Regression} (\code{state_transition.py}) & Panel makro mingguan Olist hanya berjumlah 87 minggu. Regresi vektor autoregresif ber-regularisasi L2 ($\alpha = 0,1$) menjamin solusi matematis tertutup dan kestabilan jangka panjang tanpa overfitting jaringan saraf. \\
\textbf{5. Pemisahan Backend \& UI} & \textbf{FastAPI} (port 8000) terpisah dari \textbf{Streamlit} (port 8501) & UI analitik tidak boleh menjadi pintu integrasi antar-sistem. Pemisahan API memungkinkan sistem checkout atau CRM eksternal memanggil endpoint khusus T0/T2 dengan validasi skema ketat Pydantic. \\
\end{longtable}

\section{Arsitektur Modul Perangkat Lunak 5-Lapis (\texttt{src/csip/} \& \texttt{app/})}

Seluruh kode sumber dipisahkan menurut tanggung jawab tunggal (\emph{separation of concerns}) menjadi lima lapisan arsitektur:

\begin{longtable}{@{} >{\raggedright\arraybackslash}p{3.3cm} >{\raggedright\arraybackslash}p{4.6cm} >{\raggedright\arraybackslash}p{7.9cm} @{}}
\caption{Arsitektur Lima Lapis Perangkat Lunak CSIP} \\
\toprule
\textbf{Lapisan Arsitektur} & \textbf{Direktori \& File Kunci} & \textbf{Tanggung Jawab Rekayasa Utama} \\
\midrule
\endfirsthead
\toprule
\textbf{Lapisan Arsitektur} & \textbf{Direktori \& File Kunci} & \textbf{Tanggung Jawab Rekayasa Utama} \\
\midrule
\endhead
\bottomrule
\endfoot
\textbf{1. Data \& Cleaning Layer} & \code{src/csip/data/}, \code{src/csip/cleaning/} & Pemuatan hemat memori (\code{loader.py}), validasi kontrak 4-tingkat (\code{contracts.py}), silsilah SHA-256 (\code{lineage.py}), konsolidasi relasional (\code{cleaner.py}), dan pembuat kamus bahasa alami (\code{semantic_profiler.py}). \\
\textbf{2. Feature \& Hypothesis Layer} & \code{src/csip/features/}, \code{src/csip/hypotheses/} & Pembentukan 4 objek analitik (\code{order_level}, \code{customer_cohort}, \code{review_corpus}, \code{weekly_macro}), fitur interaksi, isolasi waktu, serta uji hipotesis H1--H6 dengan koreksi FDR (\code{runner.py}). \\
\textbf{3. Modeling, Calibration \& World Layer} & \code{src/csip/models/}, \code{evaluation/}, \code{discovery/}, \code{world/}, \code{uncertainty/} & Tangga baseline B0--B4 (\code{classic.py}, \code{trainer.py}, \code{deep_tabular.py}, \code{deep_seq.py}), CV embargo 90 hari (\code{temporal_cv.py}), kalibrasi Platt \& SHAP (\code{explainer.py}), ketahanan wilayah (\code{robustness.py}), survival Kaplan-Meier (\code{survival_analysis.py}), klastering (\code{clusterer.py}), VAR World Model (\code{state_transition.py}), bootstrap CI, dan mesin universal (\code{ml_engine.py}). \\
\textbf{4. Decision \& MLOps Layer} & \code{src/csip/decision_engine/}, \code{monitoring/}, \code{evidence/} & Penilaian aturan bisnis, kalibrasi kuota CS, desain A/B holdout (\code{rules.py}, \code{engine.py}, \code{schemas.py}), pengawas drift (\code{drift_engine.py}), pelatih ulang (\code{retrainer.py}), registri promosi (\code{registry.py}), dan registri bukti kompetensi (\code{evidence/registry.py}). \\
\textbf{5. Serving \& UI Layer} & \code{src/csip/api/}, \code{app/} & Peladen REST API FastAPI 7 endpoint HTTP (\code{server.py}, \code{client.py}) beserta antarmuka kokpit CSIP Studio (\code{streamlit_app.py}, \code{data_provider.py}, \code{ui_components.py}, \code{components/charts.py}). \\
\end{longtable}

\section{Pengukuran Jejak Sumber Daya Aktual (\texttt{resource\_benchmark.json})}

CSIP tidak menebak penggunaan memori dari perkiraan kertas. Pengukuran aktual dijalankan langsung pada sampel 1.000 baris \code{order_level.parquet} sambil memuat ketiga model LightGBM inti di CPU (\code{artifacts/reports/resource_benchmark.json}):

{\small
\begin{longtable}{@{} >{\raggedright\arraybackslash}p{7.4cm} r r r r @{}}
\caption{Benchmark Penggunaan Memori \& Waktu Inferensi Aktual Model LightGBM} \\
\toprule
\textbf{Model Artefak (.joblib)} & \textbf{Muat (ms)} & \textbf{Batch (ms)} & \textbf{/Baris (ms)} & \textbf{Tunggal (ms)} \\
\midrule
\endfirsthead
\toprule
\textbf{Model Artefak (.joblib)} & \textbf{Muat (ms)} & \textbf{Batch (ms)} & \textbf{/Baris (ms)} & \textbf{Tunggal (ms)} \\
\midrule
\endhead
\bottomrule
\endfoot
\code{best_negative_review_lightgbm.joblib} & 4.910,88* & 13,02 & 0,0130 & 3,78 \\
\code{best_early_warning_review_lightgbm.joblib} & 5,61 & 4,54 & 0,0045 & 2,94 \\
\code{best_delivery_delay_lightgbm.joblib} & 4,40 & 4,33 & 0,0043 & 2,62 \\
\end{longtable}
}

*Catatan: Waktu muat model pertama (4.910,88 ms) mencakup inisialisasi awal pustaka LightGBM dan scikit-learn dari disk ke memori Python (\emph{cold-start}). Setelah siap di memori, total pemakaian RAM proses seluruh aplikasi tercatat \textbf{263,68 MB} (\emph{Resident Set Size / RSS}), dengan alokasi memori puncak Python (\code{tracemalloc peak}) sebesar \textbf{51,57 MB} --- kurang dari 4\% kapasitas laptop RAM 8 GB.

% ------------------------------------------------------------------------------
\chapter{Bedah Ekosistem Data Olist (9 Tabel), Pembersihan Deterministik \& Kontrak Data 4-Tingkat}
\chaptermark{Data Olist 9 Tabel \& Kontrak Data}

\section{Anatomi 9 Tabel Mentah Olist: Mengapa Ini Ekosistem Kompleks?}

Dataset publik \emph{Brazilian E-Commerce Public Dataset by Olist} (rentang transaksi September 2016 hingga Oktober 2018) kerap diremehkan sebagai ``satu dataset tunggal''. Padahal dari sudut pandang Data Engineering dan Data Science, ia adalah \textbf{basis data relasional 9 tabel penuh} yang menampung lebih dari \textbf{1,5 juta rekaman mentah} dengan dinamika bisnis geografis skala benua across 27 negara bagian di Brasil.

{\small
\begin{longtable}{@{} >{\raggedright\arraybackslash}p{6.2cm} r >{\raggedright\arraybackslash}p{2.2cm} >{\raggedright\arraybackslash}p{5.5cm} @{}}
\caption{Anatomi Sembilan Tabel Mentah Olist dan Jumlah Rekaman Hasil Pembersihan} \\
\toprule
\textbf{Nama Tabel Mentah (CSV)} & \textbf{Baris} & \textbf{Baris Bersih} & \textbf{Peran Bisnis \& Kualitas Data Awal} \\
\midrule
\endfirsthead
\toprule
\textbf{Nama Tabel Mentah (CSV)} & \textbf{Baris} & \textbf{Baris Bersih} & \textbf{Peran Bisnis \& Kualitas Data Awal} \\
\midrule
\endhead
\bottomrule
\endfoot
\code{olist_orders_dataset.csv} & 99.441 & 96.470 & Tabel fakta utama: status pemenuhan dan 5 stempel waktu (checkout, approval, serah kurir, tiba aktual, janji estimasi). \\
\code{olist_order_items_dataset.csv} & 112.650 & Diagregasi & Rincian barang, harga, ongkir, dan ID penjual. Satu pesanan dapat memuat banyak item dan multi-penjual. \\
\code{olist_order_payments_dataset.csv} & 103.886 & Diagregasi & Pembayaran terpecah (kartu kredit, voucher, boleto) dan jumlah tenor cicilan. \\
\code{olist_order_reviews_dataset.csv} & 104.719 & 1 / order & Rating 1--5 dan komentar bahasa Portugis. Mengandung ulasan ganda/revisi serta pesanan tanpa ulasan. \\
\code{olist_geolocation_dataset.csv} & 1.000.163 & 19.015 unik & Jutaan titik GPS dengan duplikasi berat dan sejumlah koordinat nyasar di luar teritori Brasil. \\
\code{olist_customers_dataset.csv} & 99.441 & 75.318 akun & Penghubung \code{customer_id} (unik per pesanan) ke \code{customer_unique_id} (identitas asli pelanggan). \\
\code{olist_products_dataset.csv} & 32.951 & 32.951 & Dimensi fisik barang (berat gram, panjang, tinggi, lebar), kategori, dan jumlah foto. \\
\code{olist_sellers_dataset.csv} & 3.095 & 3.095 & Lokasi kota, negara bagian, dan kode pos penjual. \\
\code{product_category_name_translation.csv} & 71 & 71 & Kamus penerjemah 71 nama kategori Portugis ke bahasa Inggris. \\
\end{longtable}
}

\section{Jebakan Rekayasa Pertama: \emph{Cartesian Fan-Out} Saat Menggabungkan Tabel}

\begin{bisnisbox}[Analogi Sederhana: Mengapa Menggabung Tabel Sembarangan Merusak Hitungan]
Bayangkan Ibu Sari membeli \textbf{3 buah baju} dalam 1 nomor pesanan. Ia membayarnya dengan \textbf{2 metode pembayaran} (separuh pakai voucher, separuh pakai kartu kredit), lalu sempat mengirim \textbf{2 kali revisi ulasan}. Jika analis langsung menggabungkan (\code{LEFT JOIN}) tabel Pesanan $\rightarrow$ Barang $\rightarrow$ Pembayaran $\rightarrow$ Ulasan tanpa meringkasnya lebih dulu, komputer akan menghasilkan $3 \times 2 \times 2 = \mathbf{12\text{ baris}}$ untuk satu nomor pesanan yang sama! Total omzet tercatat seolah membengkak 4--12 kali lipat dan model akan mengulang hafalan pesanan yang sama berkali-kali.
\end{bisnisbox}

\begin{teknisbox}[Aturan Konsolidasi Ketat di OlistDataCleaner (\code{src/csip/cleaning/cleaner.py})]
Modul \code{OlistDataCleaner} melarang penggabungan langsung ber-grain campuran dan memaksa seluruh tabel anak diagregasi terlebih dahulu ke \textbf{tepat 1 baris per \code{order_id}}:
\begin{enumerate}
  \item \textbf{Deduplikasi Geospasial:} Dari 1.000.163 ping GPS, sistem menyaring titik di dalam kotak batas geografis teritori Brasil (Latitude $-34,0$ hingga $+5,5$; Longitude $-74,0$ hingga $-34,5$). Tiap prefiks kode pos diringkas menggunakan \textbf{median koordinat} (bukan rata-rata, agar kebal terhadap pencilan GPS ekstrem), menyusutkan tabel sebesar 98,1\% menjadi \textbf{19.015 prefiks kode pos unik} (\code{clean_geolocation_dedup.parquet}).
  \item \textbf{Konsolidasi Item Belanja:} Tabel item (112.650 baris) dijumlahkan per \code{order_id} menjadi \code{total_item_price}, \code{total_freight}, \code{item_count}, serta rata-rata berat produk (\code{mean_product_weight_g}).
  \item \textbf{Konsolidasi Pembayaran:} Tabel pembayaran (103.886 baris) diringkas per \code{order_id} menjadi total nilai bayar (\code{total_payment_value}), tenor cicilan maksimum (\code{payment_installments_max}), dan metode bayar dominan (\code{payment_type_primary}).
  \item \textbf{Deduplikasi Ulasan \& Pelestarian Label Hilang:} Tabel ulasan (104.719 baris) dikelompokkan per \code{order_id} dan hanya mempertahankan \textbf{ulasan paling akhir} berdasarkan \code{review_creation_date} / \code{review_answer_timestamp}. Jika pesanan tidak pernah diulas, nilai rating dibiarkan \code{NaN} (\emph{forbidden artificial imputation}) dan dicatat pada indikator biner \code{has_observed_review} ($1$ jika ada, $0$ jika kosong). Target \code{negative_review_flag} bernilai $1,0$ jika skor $\le 2$, bernilai $0,0$ jika skor $\ge 3$, dan tetap \code{NaN} bila belum diulas.
\end{enumerate}
\end{teknisbox}

\section{Transparansi Penyaringan Populasi (\emph{Population \& Selection Bias Funnel})}

Dalam praktik riset ilmiah, analis tidak boleh menyembunyikan baris data mana yang dibuang sebelum melatih model. Berkas \code{reports/methodology/POPULATION_AND_SELECTION_BIAS.md} serta gerbang kualitas \textbf{G7} mencatat alur penyusutan populasi secara utuh:

\begin{longtable}{@{} >{\raggedright\arraybackslash}p{4.7cm} r >{\raggedright\arraybackslash}p{8.2cm} @{}}
\caption{Corong Transparansi Penyaringan Populasi Pesanan Olist} \\
\toprule
\textbf{Tahap Penyaringan Populasi} & \textbf{Jumlah Pesanan} & \textbf{Penjelasan Ilmiah \& Batas Bias Seleksi} \\
\midrule
\endfirsthead
\toprule
\textbf{Tahap Penyaringan Populasi} & \textbf{Jumlah Pesanan} & \textbf{Penjelasan Ilmiah \& Batas Bias Seleksi} \\
\midrule
\endhead
\bottomrule
\endfoot
\textbf{1. Populasi Mentah Sistem} & 99.441 & Seluruh rekaman pesanan di \code{olist_orders_dataset.csv} (2016--2018). \\
\textbf{2. Disaring ke Status \code{delivered}} & 96.478 & Memisahkan \textbf{625 pesanan \code{canceled}} (tingkat ulasan negatif 87,19\%) dan \textbf{609 pesanan \code{unavailable}} (87,60\%) plus status tengah jalan karena tidak memiliki siklus pengiriman kurir selesai. \\
\textbf{3. Konsolidasi Bersih (\code{order_level})} & 96.470 & Grain tunggal \code{order_id} siap olah setelah penggabungan bersih antar-tabel. \\
\textbf{4. Pesanan dengan Ulasan Teramati} & 95.824 & Sebanyak \textbf{646 pesanan terkirim (0,67\%)} tidak memiliki ulasan. Analisis batas seleksi terburuk (\emph{Manski bounds}) menunjukkan ketidakpastian akibat ulasan hilang ini hanya selisih $\le 0,59\%$ (rentang 11,44\% hingga 12,10\%) sehingga tidak mengubah kesimpulan model. \\
\end{longtable}

\begin{kejujuranbox}[Pernyataan Estimand Kondisional yang Jujur]
Karena 1.234 pesanan yang dibatalkan atau stoknya kosong dipisahkan pada tahap ke-2 (dan 99,0\% pesanan \code{unavailable} tidak memiliki rincian produk sama sekali di database Olist menurut laporan \code{CHECKOUT_RISK_FEASIBILITY.md}), model risiko kepuasan CSIP selalu dinyatakan secara jujur sebagai prediksi \textbf{bersyarat pada pesanan yang statusnya selesai terkirim} (\emph{conditional on fulfillment}), bukan prediksi pembatalan pesanan di gudang.
\end{kejujuranbox}

\section{Optimasi Memori RAM 8 GB, Kontrak Data 4-Tingkat \& Sidik Jari SHA-256}

Untuk menjaga stabilitas pipeline pada laptop RAM 8 GB dan mencegah data cacat masuk ke model, modul data menjalankan tiga pengaman berurutan:
\begin{enumerate}
  \item \textbf{Penghematan Memori (\code{downcast_dataframe} di \code{src/csip/data/loader.py}):} Tipe \code{float64} bawaan pandas dikonversi aman ke \code{float32}, dan \code{int64} ke \code{int8}/\code{int16}/\code{int32} sesuai rentang angka tiap kolom. Ini memangkas jejak RAM hampir 50\% saat membaca tabel besar seperti 1 juta titik GPS.
  \item \textbf{Kontrak Data 4-Tingkat (\code{src/csip/data/contracts.py}):} Setiap tahap dilindungi oleh kelas \code{DataContract} dan \code{ValidationReport} dengan 4 level keparahan melalui enum \code{ContractSeverity}:
    \begin{itemize}
      \item \code{INFO}: Profil skema normal dan statistik jumlah baris.
      \item \code{WARNING}: Peringatan non-fatal (misalnya nilai di luar kewajaran kecil).
      \item \code{RECOVERABLE_ERROR}: Anomali minor yang dapat ditangani deterministik oleh cleaner.
      \item \code{CRITICAL_FAILURE}: Cacat fatal (kolom wajib hilang, primary key \code{order_id} duplikat, atau tipe rusak) yang memicu \code{DataContractViolationError} melalui \code{report.assert_pipeline_valid()} dan \textbf{menghentikan seluruh pipeline seketika}.
    \end{itemize}
  \item \textbf{Garis Silsilah Kriptografi SHA-256 (\code{src/csip/data/lineage.py}):} Setiap berkas mentah hingga tabel Parquet akhir direkam sidik jari SHA-256-nya ke \code{artifacts/metadata/DATA_LINEAGE.json} agar setiap eksperimen selalu bisa diaudit ke versi data asalnya.
  \item \textbf{Strategi Distribusi ZIP $<$ 100 MB:} File CSV mentah di \code{data/raw/} berukuran total $>120$ MB (khusus tabel GPS saja 58,44 MB). Oleh karena itu paket distribusi \code{csip_v1.1.0.zip} ($\sim$82 MB) menyertakan tabel olahan bersih \code{data/processed/} dan \code{data/interim/} siap pakai, sambil menyediakan skrip \code{scripts/download_raw_olist.py} bagi siapa saja yang ingin memverifikasi atau mengulang tahap pembersihan dari CSV mentah Kaggle.
\end{enumerate}

% ------------------------------------------------------------------------------
\chapter{Rekayasa Fitur, Isolasi Waktu ($T_0, T_1, T_2$) \& Pertahanan Anti-Leakage}
\chaptermark{Rekayasa Fitur \& Isolasi Waktu T0/T1/T2}

\section{Analogi Kebocoran Data (\emph{Data Leakage}): Kesalahan Paling Mematikan di Industri}

\begin{bisnisbox}[Analogi Prakiraan Cuaca yang Menipu]
Bayangkan sebuah aplikasi mengklaim bisa memprediksi hujan besok pagi dengan akurasi 99\%. Namun setelah dibongkar, salah satu kolom input yang dipakainya adalah \emph{``apakah trotoar besok pagi basah?''}. Tentu saja modelnya terlihat hebat di kertas karena masa depan sudah bocor ke soal ujian! Begitu diluncurkan malam ini, kolom besok pagi belum tersedia dan aplikasi gagal total.
\end{bisnisbox}

Di proyek e-commerce, kesalahan paling sering terjadi persis seperti itu: analis memasukkan fitur \code{delivery_delay_days} (apakah paket telat atau tidak) atau \code{delivery_lead_time_days} (berapa lama paket sampai), lalu mengklaim modelnya mampu ``memperingatkan risiko ulasan buruk saat pelanggan baru membayar di halaman checkout''.

CSIP memberantas kebocoran ini di level kode sumber dengan mendefinisikan \textbf{tiga pos waktu keputusan yang terisolasi ketat}.

\section{Tiga Pos Waktu Operasional ($T_0, T_1, T_2$)}

Modul \code{src/csip/features/engineer.py}, \code{src/csip/models/trainer.py}, dan \code{src/csip/api/server.py} mengunci daftar fitur secara eksplisit untuk masing-masing titik keputusan bisnis:

\subsection{Pos $T_0$ --- Titik Checkout Pembayaran (\texttt{E1\_checkout\_retrospective\_conditional}, 10 Fitur)}

Pada detik pembeli menekan tombol bayar, paket belum dikemas dan belum diserahkan ke kurir. Seluruh informasi perjalanan pengiriman \textbf{diblokir keras}. Hanya 10 fitur sah berikut yang diizinkan:

\begin{longtable}{@{} >{\raggedright\arraybackslash}p{4.8cm} >{\raggedright\arraybackslash}p{11.2cm} @{}}
\caption{Sepuluh Fitur Ekonomi \& Geografis pada Pos Waktu T0 (Checkout)} \\
\toprule
\textbf{Nama Fitur (Pos $T_0$)} & \textbf{Makna Bisnis \& Cara Pembuatannya di Kode} \\
\midrule
\endfirsthead
\toprule
\textbf{Nama Fitur (Pos $T_0$)} & \textbf{Makna Bisnis \& Cara Pembuatannya di Kode} \\
\midrule
\endhead
\bottomrule
\endfoot
\code{total_item_price} & Total harga murni seluruh barang dalam satu nomor pesanan. \\
\code{total_freight} & Total ongkos kirim yang dibebankan kepada pembeli. \\
\code{total_order_value} & Nilai transaksi keseluruhan (\code{total_item_price + total_freight}). \\
\code{freight_ratio} & Rasio ongkir terhadap harga barang (\code{total_freight / total_item_price}). Memotret rasa ketidakadilan biaya kirim (misalnya harga barang murah tapi ongkirnya hampir setengah harga barang). \\
\code{item_count} & Jumlah unit barang dalam satu keranjang belanja. \\
\code{payment_installments_max} & Tenor cicilan maksimum sebagai indikator komitmen finansial pembeli. \\
\code{seller_buyer_distance_km} & Jarak tempuh garis lurus antara centroid kode pos penjual dan pembeli dalam satuan kilometer (dihitung vektorisasi menggunakan \code{haversine_np} di \code{engineer.py}). \\
\code{estimated_lead_time_days} & Janji durasi kirim yang ditampilkan kepada pembeli saat checkout (selisih tanggal estimasi sistem minus tanggal beli). \\
\code{distance_x_weight} & Fitur interaksi antara jarak pengiriman dan rata-rata berat produk untuk memotret beban fisik ekspedisi lintas pulau/negara bagian. \\
\code{prior_order_count} & Jumlah pesanan masa lalu yang sudah pernah dilakukan \code{customer_unique_id} \textbf{sebelum stempel waktu transaksi saat ini} (\emph{strictly backward-looking as-of feature}). Transaksi masa depan pelanggan yang sama dijamin tidak ikut terhitung. \\
\end{longtable}

\subsection{Pos $T_1$ --- Persetujuan Pembayaran \& Disposisi Kurir (\texttt{DELAY\_FEATURES}, 5 Fitur)}

Digunakan oleh model regresi keterlambatan \code{best_delivery_delay_lightgbm.joblib} untuk memperkirakan potensi hari keterlambatan logistik sebelum paket menuntaskan perjalanannya:
\begin{itemize}
  \item \code{total_order_value}, \code{freight_ratio}, \code{estimated_lead_time_days}, \code{seller_buyer_distance_km}, serta \code{approval_lag_hours} (jeda waktu dalam jam antara pembelian dibuat hingga verifikasi pembayaran lunas oleh sistem bank/gerbang bayar).
\end{itemize}

\subsection{Pos $T_2$ --- Pasca-Pengiriman (\texttt{E3\_combined\_interactions}, 15 Fitur)}

Baru diizinkan setelah paket dinyatakan tiba di tangan pembeli. Tujuannya adalah \textbf{diagnosis akar masalah, pemeringkatan keluhan sebelum atau sesudah survei ditutup, dan alokasi tim CS}. Di atas fitur ekonomi dan geografis sebelumnya, Pos $T_2$ menambahkan lima sinyal pengalaman nyata kurir:

\begin{longtable}{@{} >{\raggedright\arraybackslash}p{4.8cm} >{\raggedright\arraybackslash}p{11.2cm} @{}}
\caption{Lima Fitur Sinyal Aktual pada Pos Waktu T2 (Pasca-Pengiriman)} \\
\toprule
\textbf{Fitur Tambahan (Pos $T_2$)} & \textbf{Makna Operasional \& Hipotesis Bisnis} \\
\midrule
\endfirsthead
\toprule
\textbf{Fitur Tambahan (Pos $T_2$)} & \textbf{Makna Operasional \& Hipotesis Bisnis} \\
\midrule
\endhead
\bottomrule
\endfoot
\code{delivery_lead_time_days} & Lama hari tunggu aktual yang benar-benar dirasakan pembeli sejak bayar hingga paket tiba. \\
\code{delivery_delay_days} & Selisih hari kedatangan aktual dikurangi tanggal janji estimasi. Nilai $>0$ artinya terlambat; nilai $\le 0$ artinya tiba lebih cepat atau tepat waktu. \\
\code{is_late} & Indikator biner ($1$ jika \code{delivery_delay_days > 0}, $0$ jika tidak). \\
\code{delay_x_freight_ratio} & Fitur interaksi \textbf{kekecewaan ganda} (\emph{double frustration}): pembeli sudah membayar ongkir yang mahal relatif terhadap harga barang, tetapi pengirimannya terlambat pula. \\
\code{high_value_late_flag} & Indikator biner khusus untuk pesanan bernilai tinggi yang mengalami keterlambatan --- segmen dengan risiko kerugian reputasi terbesar. \\
\end{longtable}

\section{Partisi Waktu Kronologis \& Validasi Silang dengan Embargo 90 Hari}

Selain memisahkan pos fitur, CSIP melarang keras penggunaan \emph{random K-fold split} pada evaluasi model Olist karena akan mencampur pesanan dari bulan yang sama di masa depan ke dalam data latih masa lalu.

Dalam \code{config/analysis_plan.yaml}, \code{src/csip/models/trainer.py}, dan \code{src/csip/evaluation/temporal_cv.py}, pemecahan data mengikuti urutan kronologis mutlak:
\begin{enumerate}
  \item \textbf{Partisi Latih (\emph{Train Set}):} Pesanan hingga \code{2018-03-31} ($N = 63.666$ baris pada pelatihan model ulasan negatif; $64.156$ pada regresi keterlambatan).
  \item \textbf{Partisi Validasi (\emph{Validation Set}):} Pesanan dari \code{2018-04-01} hingga \code{2018-06-30} ($N = 19.587$ pada ulasan negatif; $19.684$ pada keterlambatan). Digunakan secara eksklusif untuk pemilihan model, penyetelan, dan fitting kalibrasi Platt Scaling.
  \item \textbf{Partisi Uji Terkunci (\emph{Locked Holdout Test Set}):} Pesanan setelah \code{2018-06-30} ($N = 12.571$ untuk evaluasi artefak \code{.joblib} ulasan negatif; $N = 12.630$ untuk keterlambatan pengiriman). Partisi ini disegel hingga seluruh keputusan pemodelan selesai dan hanya dinilai satu kali sebagai bukti kemampuan generalisasi ke masa depan.
\end{enumerate}

Untuk mengecek stabilitas lintas musim, kelas \code{RollingOriginTemporalCV} menjalankan \textbf{4 lipatan jendela mengembang (\emph{4 expanding folds})} dengan aturan wajib \textbf{Embargo minimal 90 Hari} antara pesanan terakhir di partisi latih dan pesanan pertama di partisi validasi tiap lipatan. Mengapa harus 90 hari? Karena durasi kirim paket jarak jauh Brasil ditambah tenggat pengisian ulasan dan jendela pembelian ulang berada dalam rentang waktu berminggu-minggu; jeda 90 hari memastikan nol tumpang tindih label antara data latih dan data validasi. Pada kohort retensi (\code{repeat_purchase}), data di ujung kalender (\code{2018-07-19} hingga \code{2018-10-17}, $2.275$ pelanggan) juga disaring sebagai \emph{right-censoring buffer} serta mengeluarkan 962 pelanggan yang sudah belanja ulang sebelum paket pertamanya sempat datang agar definisi target 90 hari pasca-pengiriman tetap bersih.


% ==============================================================================
% BAGIAN II: SAINS DATA, PEMODELAN, HIPOTESIS & KEJUJURAN EMPIRIS
% ==============================================================================
\part{Sains Data, Pemodelan, Hipotesis \& Kejujuran Empiris}

\chapter{Tangga Baseline B0--B4 \& Cara Kerja Seluruh Algoritma}
\chaptermark{Tangga Baseline B0--B4 \& Algoritma}

\begin{bisnisbox}[Analogi Bisnis: Mengapa Tidak Langsung Pakai AI Tercanggih?]
Bayangkan Anda membuka cabang minimarket baru. Sebelum menghabiskan anggaran untuk memasang kamera sensor pelacakan wajah berbiaya mahal, manajer toko yang cerdas tentu bertanya lebih dulu: \emph{``Kalau kita hanya pakai catatan kasir manual atau rumus Excel sederhana, seberapa akurat tebakan kita?''}

Di dunia Sains Data, pemikiran sehat itulah yang disebut \textbf{Tangga Baseline (B0 hingga B4)}. Banyak proyek AI gagal meyakinkan direksi karena langsung memamerkan model berat dengan skor ``88\% akurasi'', padahal menebak semua transaksi sebagai \emph{``aman''} secara buta saja sudah menghasilkan akurasi 88\% karena mayoritas pembeli memang puas!

CSIP memaksa setiap algoritma baru memanjat anak tangga secara adil: mulai dari tebakan paling bodoh (B0), model statistik linear klasik (B1), algoritma pohon dan kedekatan jarak klasik (B2), rakitan pohon keputusan modern (B3), hingga eksperimen jaringan saraf tiruan dan rangkaian waktu (B4). Setiap langkah harus membuktikan apakah kerumitan tambahan benar-benar memberi kenaikan performa yang sepadan dengan beban memori dan biaya perawatan.
\end{bisnisbox}

\section{Filosofi Tangga Baseline B0 hingga B4}

Disiplin pemodelan di modul \code{src/csip/models/classic.py}, \code{src/csip/models/trainer.py}, dan \code{src/csip/models/deep_tabular.py} menetapkan lima tingkatan komparasi wajib:

\begin{enumerate}
  \item \textbf{Tingkat B0 --- Tolok Ukur Naif (\emph{Naive Baselines}):}
    \begin{itemize}
      \item \code{B0_majority_class} (\code{DummyClassifier(strategy="most_frequent")}): Untuk tugas klasifikasi, selalu menebak kelas mayoritas (tidak ada keluhan). Berguna sebagai lantai dasar metrik Precision-Recall.
      \item \code{B0_delay_heuristic}: Aturan bisnis deterministik murni tanpa pembelajaran mesin, yakni menandai risiko keluhan hanya berdasarkan fakta apakah paket datang terlambat (\code{delivery_delay_days > 0}). Pada regresi keterlambatan, pembanding B0 adalah \code{Naive Mean} yang sekadar menebak angka rata-rata historis untuk semua pengiriman.
    \end{itemize}
  \item \textbf{Tingkat B1 --- Model Linear Teregulasi (\emph{Regularized Linear Models}):}
    \begin{itemize}
      \item \code{B1_logistic} (\code{LogisticRegression(class_weight="balanced", max_iter=1000)} yang didahului \code{StandardScaler}): Menimbang kontribusi tiap fitur secara proporsional. Karena keluhan hanya mencakup sekitar 11\%--12\% pesanan, argumen \code{class_weight="balanced"} memberi bobot penalti kesalahan lebih tinggi pada kelas minoritas agar model tidak malas menebak semua aman.
      \item \code{B1_ridge} (\code{Ridge(alpha=1.0)} setelah penskalaan fitur): Digunakan untuk memprediksi angka kontinu keterlambatan hari. Penalti L2 (\code{alpha=1.0}) mencegah koefisien fitur saling tarik-menarik menjadi nilai ekstrem saat dua fitur jarak dan ongkir berkorelasi tinggi.
    \end{itemize}
  \item \textbf{Tingkat B2 --- Pemodelan Non-Linear Klasik (\emph{Single Interpretable \& Non-Parametric Models}):}
    \begin{itemize}
      \item \code{B2_decision_tree} (\code{DecisionTreeClassifier(max_depth=6, min_samples_split=50)}): Membangun aturan bercabang berbentuk \emph{Jika-Maka} yang mudah dibaca tim operasional. Kedalaman dikunci maksimal 6 lapis dan minimal 50 sampel sebelum membelah simpul agar pohon tidak menghafal kebisingan data latih (\emph{overfitting}).
      \item \code{B2_knn} (\code{KNeighborsClassifier(n_neighbors=15, weights="distance")}): Memprediksi nasib pesanan baru dengan mencari 15 pesanan paling mirip di masa lalu dan memberi suara lebih besar pada tetangga yang jarak fiturnya paling dekat.
      \item \code{B2_naive_bayes} (\code{GaussianNB()}): Pengklasifikasi probabilistik cepat dengan asumsi fitur saling independen satu sama lain. Sering kali kurang terkalibrasi secara probabilitas, tetapi berharga sebagai pembanding berkecepatan mikrodetik.
    \end{itemize}
  \item \textbf{Tingkat B3 --- Ensembel Pohon Gradien \& Hutan Acak (\emph{Tabular Ensemble Champions}):}
    \begin{itemize}
      \item \code{B3_random_forest} (\code{RandomForestClassifier(n_estimators=100, max_depth=8, min_samples_leaf=20, class_weight="balanced")}): Melatih 100 pohon keputusan independen di atas sub-sampel data acak lalu merata-ratakan voting mereka. Sangat kokoh meredam varians tanpa banyak penyetelan rumit.
      \item \code{B3_lightgbm} (\code{LGBMClassifier} / \code{LGBMRegressor} dengan \code{n_estimators=150}, \code{max_depth=6}, \code{num_leaves=31}, \code{learning_rate=0.05}, \code{class_weight="balanced"}): Berbeda dengan Random Forest yang pohonnya bekerja paralel secara terpisah, LightGBM membangun 150 pohon secara berurutan di mana setiap pohon baru khusus menambal kesalahan prediksi pohon sebelumnya. Teknik berbasis histogram ini hemat memori RAM dan cepat di CPU laptop standar.
    \end{itemize}
  \item \textbf{Tingkat B4 --- Pemodelan Dalam \& Runtun Waktu Eksperimental (\emph{Deep Tabular \& Sequence}):}
    \begin{itemize}
      \item \code{TabularEmbeddingMLP} (\code{src/csip/models/deep_tabular.py}): Jaringan saraf tiruan dengan arsitektur lapisan tersembunyi \code{[128, 64]}, \code{BatchNorm1d}, aktivasi \code{ReLU}, dan \code{Dropout(0.2)} yang dijalankan pada PyTorch CPU. Dipakai sebagai uji coba apakah struktur tabel CSIP mendapat untung dari representasi mendalam dibandingkan pohon gradien.
      \item \code{TextBiLSTM}: Pemroses teks ulasan berbahasa Portugis pasca-pengiriman memakai Bidirectional LSTM + Attention pada PyTorch CPU. Dibatasi ketat ke lapisan diagnostik pasca-kejadian karena komentar pelanggan belum eksis saat transaksi terjadi.
    \end{itemize}
\end{enumerate}

\section{Panduan Cara Kerja Seluruh Algoritma dalam Bahasa Manusia \& Sains Data}

Agar buku ini menjadi rujukan utuh baik bagi direktur maupun rekayasawan, subbab ini membedah delapan algoritma inti CSIP tanpa rumus matematika di luar ruang lingkup Sains Data.

\subsection{1. Regresi Logistik, Regresi Ridge \& Mengapa Kesederhanaan Sering Menang}
Secara intuitif, \textbf{Regresi Logistik} bekerja seperti papan skor timbangan berbobot. Setiap fitur---misalnya keterlambatan pengiriman atau rasio ongkos kirim---mendapat bobot angka positif atau negatif. Semua nilai lalu dijumlahkan menjadi skor log-odds dan dilewatkan ke kurva S (\emph{sigmoid}) agar berubah rapi menjadi peluang di rentang $0$ sampai $1$ (misalnya $0{,}74$ atau $74\%$).

Sementara itu, \textbf{Regresi Ridge} bekerja mirip untuk target berupa angka kontinu (seperti menebak keterlambatan sekian hari), tetapi dengan tambahan rem penalti L2. Mengapa butuh rem tersebut? Karena jarak pengiriman dan tarif ongkir saling berkorelasi; tanpa rem, komputer bisa memberi bobot $+500$ pada jarak dan $-498$ pada ongkir yang membuat prediksi meledak saat ada data baru yang sedikit menyimpang. Penalti L2 menjaga bobot fitur tetap kalem dan stabil saat menghadapi data masa depan.

\subsection{2. Random Forest vs. LightGBM (Ensembel Pohon)}
\textbf{Random Forest} ibarat komite 100 juri independen yang masing-masing hanya diberi akses ke sebagian catatan dan sebagian kolom. Karena mereka tidak saling mencontek, kesimpulan rata-ratanya jauh lebih stabil daripada satu pohon keputusan tunggal.

Sebaliknya, \textbf{LightGBM} ibarat tim estafet peneliti. Pohon ke-1 mencoba memprediksi risiko keluhan; lalu dicatat pesanan mana yang tebakannya meleset paling jauh (residu kesalahan). Pohon ke-2 dilatih khusus untuk menebak sisa kesalahan tersebut, pohon ke-3 menambal sisa pohon ke-2, dan seterusnya hingga 150 tahap dengan langkah belajar lembut (\code{learning_rate=0.05}). Keunggulan LightGBM pada data tabel e-commerce adalah kemampuannya menangkap aturan non-linear kompleks---misalnya \emph{``terlambat 2 hari pada barang murah hanya menaikkan risiko sedikit, tetapi terlambat 2 hari pada barang di atas \$200 BRL memicu lonjakan keluhan tajam''}---secara otomatis tanpa perlu puluhan aturan manual.

\subsection{3. K-Means Clustering \& Validasi Jumlah Klaster Optimal}
Untuk mengelompokkan puluhan ribu transaksi tanpa label target (\emph{unsupervised learning}), CSIP menggunakan \textbf{K-Means} pada fitur numerik yang telah distandardisasi agar skala rupiah/real Brasil tidak mengalahkan skala kilometer jarak. Algoritma ini memilih sejumlah titik pusat (\emph{centroid}), memasukkan setiap transaksi ke pusat terdekat, menghitung ulang rata-rata posisi tiap kelompok, lalu mengulang proses sampai keanggotaan stabil.

Berapa jumlah kelompok ($K$) yang paling masuk akal? CSIP tidak menebak angka $K$ berdasarkan selera visual, melainkan mengadu kandidat $K=2, 3, 4, 5, \dots$ dengan tiga indeks objektivitas klaster:
\begin{itemize}
  \item \textbf{Silhouette Score} (rentang $-1$ sampai $+1$, semakin tinggi semakin baik): Mengukur seberapa rapat anggota dalam satu kelompok dibandingkan jaraknya ke kelompok tetangga terdekat.
  \item \textbf{Davies-Bouldin Index} (semakin rendah mendekati $0$ semakin baik): Mengukur rasio kekompakan internal terhadap jarak antarpusat klaster.
  \item \textbf{Calinski-Harabasz Index} (semakin tinggi semakin baik): Rasio keragaman antarkelompok terhadap keragaman di dalam kelompok.
\end{itemize}

\subsection{4. Analisis Ketahanan Hidup Kaplan-Meier (\emph{Survival Analysis})}
Untuk pertanyaan bertema waktu seperti \emph{``berapa persen pembeli yang masih berpotensi datang belanja ulang setelah hari ke-30, ke-60, hingga ke-90 pascapenerimaan paket?''}, metode klasifikasi biasa sering keliru karena data \emph{tersensor} (\emph{right-censored}): pembeli yang baru berbelanja pada bulan Agustus 2018 belum punya kesempatan hidup penuh selama 90 hari sebelum pengumpulan data berakhir pada \code{2018-10-17}.

Estimator \textbf{Kaplan-Meier} secara elegan menghitung probabilitas bertahan secara bertahap pada setiap titik waktu tanpa membuang informasi dari pelanggan yang periode pantauannya belum selesai. Untuk menguji apakah perbedaan dua kelompok (misalnya pelanggan yang puas vs pelanggan yang mengalami keterlambatan) bermakna secara statistik, digunakan uji \emph{log-rank}.

\subsection{5. VAR Ridge untuk Simulator Kebijakan Makro Mingguan}
Kebutuhan simulasi tidak selalu berhenti pada level satu pesanan. Direksi sering bertanya: \emph{``Jika rata-rata keterlambatan seluruh kurir turun bulan depan, bagaimana dampaknya terhadap proporsi keluhan mingguan selama 6 sampai 12 minggu ke depan?''}

Di sini CSIP membangun agregasi seri waktu mingguan (87 minggu historis) lalu melatih model \textbf{Vector Autoregression (VAR) dengan regularisasi Ridge}. Secara konsep Data Science, VAR memprediksi nilai indikator minggu ini (tingkat keluhan, rata-rata keterlambatan, rasio ongkir, volume pesanan) sebagai fungsi linear dari nilai indikator-indikator tersebut pada 1--2 minggu sebelumnya (\emph{lags}), dengan regularisasi Ridge untuk mencegah kelebihan parameter karena panjang deret waktu mingguan hanya puluhan baris.

\begin{kejujuranbox}[Batas Wajar Penggunaan Model VAR]
Simulator makro menggunakan model autoregresif berbasis data observasional historis, bukan uji coba acak terkendali di tingkat negara. Oleh karena itu, seluruh keluaran skenario diklasifikasikan sebagai \textbf{proyeksi simulasi terkalibrasi secara historis} (Level Klaim L2), bukan jaminan kausal absolut saat terjadi perubahan struktur pasar ekstrem di dunia nyata.
\end{kejujuranbox}

\chapter{Evaluasi Empiris, Papan Peringkat Model \& Rekonsiliasi Audit Forensik}
\chaptermark{Leaderboard Empiris \& Audit Forensik}

Bab ini memaparkan angka empiris nyata yang tercatat di \code{artifacts/reports/ml_leaderboard.csv}, \code{artifacts/reports/locked_test_reevaluation.json}, serta \code{artifacts/metadata/MODEL_RESULTS_REGISTRY.csv}. Tidak ada satu pun angka yang dipoles atau diperhalus; bahkan temuan anomali dan model yang gagal ikut dipublikasikan secara terbuka.

\section{Papan Peringkat Empiris Seluruh Model (Leaderboard Lengkap)}

Mengapa CSIP menekankan metrik \textbf{PR-AUC} (Precision-Recall Area Under Curve) di samping ROC-AUC pada prediksi ulasan negatif? Karena hanya sekitar 11\%--12\% pesanan Olist yang berujung ulasan buruk (1 atau 2 bintang). Pada data tidak seimbang (\emph{imbalanced data}), metrik akurasi biasa maupun ROC-AUC mudah tampak tinggi meski model sering menebak positif palsu. PR-AUC berfokus murni pada ketepatan menangkap kelompok keluhan minoritas tersebut.

\begin{longtable}{@{} >{\raggedright\arraybackslash}p{2.5cm} >{\raggedright\arraybackslash}p{3.2cm} r r r r >{\raggedright\arraybackslash}p{2.2cm} @{}}
\caption{Papan Peringkat Empiris Klasifikasi Ulasan Negatif \& Retensi pada Set Uji Terkunci} \\
\toprule
\textbf{Tugas / Pos Fitur} & \textbf{Model} & \textbf{Val PR} & \textbf{Test PR} & \textbf{Test ROC} & \textbf{Test Brier} & \textbf{Status} \\
\midrule
\endfirsthead
\multicolumn{1}{@{}l}{\small\emph{(Tabel lanjutan)}} \\
\toprule
\textbf{Tugas / Pos Fitur} & \textbf{Model} & \textbf{Val PR} & \textbf{Test PR} & \textbf{Test ROC} & \textbf{Test Brier} & \textbf{Status} \\
\midrule
\endhead
\bottomrule
\endfoot
\multicolumn{7}{l}{\textbf{A. Prediksi Ulasan Negatif --- Pasca-Pengiriman (Pos $T_2$ / \code{E3})}} \\
Pos $T_2$ (E3) & \code{B1_logistic} & 0{,}3134 & 0{,}2436 & 0{,}6803 & --- & Baseline Linear \\
Pos $T_2$ (E3) & \code{B2_decision_tree} & --- & 0{,}2704 & 0{,}6642 & --- & Pembanding Pohon \\
Pos $T_2$ (E3) & \code{B2_knn} & --- & 0{,}2532 & 0{,}6383 & --- & Pembanding Jarak \\
Pos $T_2$ (E3) & \code{B2_naive_bayes} & --- & 0{,}1996 & 0{,}6523 & --- & Pembanding Prob. \\
Pos $T_2$ (E3) & \code{B3_random_forest} & --- & 0{,}3016 & 0{,}6727 & --- & Runner-Up \\
Pos $T_2$ (E3) & \code{B3_lightgbm} & --- & \textbf{0{,}3061} & \textbf{0{,}6745} & \textbf{0{,}1373*} & \textbf{Champion Utama} \\
\midrule
\multicolumn{7}{l}{\textbf{B. Peringatan Dini saat Checkout (Pos $T_0$ / \code{E1}, Tanpa Info Kurir Aktual)}} \\
Pos $T_0$ (E1 EW) & \code{B3_lightgbm} & --- & 0{,}1335 & 0{,}5468 & 0{,}2370 & Sinyal Lemah \\
\midrule
\multicolumn{7}{l}{\textbf{C. Prediksi Pembelian Ulang 90 Hari (Kohort Retensi)}} \\
Retensi 90 Hari & \code{B3_lightgbm} & --- & 0{,}0202 & 0{,}5021 & --- & \textbf{DITOLAK (Acak)} \\
\end{longtable}

\noindent\emph{*Catatan Brier Score \code{B3_lightgbm} Pos $T_2$: Nilai Brier mentah \code{0,1373} (dan \code{0,137478} pada sub-sampel evaluasi \code{model_explainability.json}) membaik secara drastis setelah kalibrasi Platt Scaling menjadi \textbf{0,076177} dengan ECE turun dari \textbf{0,235486} menjadi \textbf{0,007015} (dibedah rinci di Bab 8).}

Untuk tugas \textbf{Regresi Estimasi Hari Keterlambatan} (memprediksi berapa hari paket akan meleset atau lebih cepat dari janji sebelum terjadi), papan peringkat pada set uji terkunci menampilkan hasil berikut:

\begin{longtable}{@{} >{\raggedright\arraybackslash}p{4.1cm} r r r >{\raggedright\arraybackslash}p{4.8cm} @{}}
\caption{Papan Peringkat Empiris Regresi Estimasi Hari Keterlambatan} \\
\toprule
\textbf{Model Regresi Delay} & \textbf{MAE} & \textbf{RMSE} & \textbf{$R^2$} & \textbf{Catatan Audit Teknis} \\
\midrule
\endfirsthead
\multicolumn{1}{@{}l}{\small\emph{(Tabel lanjutan)}} \\
\toprule
\textbf{Model Regresi Delay} & \textbf{MAE} & \textbf{RMSE} & \textbf{$R^2$} & \textbf{Catatan Audit Teknis} \\
\midrule
\endhead
\bottomrule
\endfoot
\code{Naive Mean} (B0 Validasi) & 7{,}4529 hari & --- & --- & Tebakan rata-rata historis \\
\code{B1_ridge} (Set Uji Terkunci) & \textbf{4{,}7509 hari} & \textbf{5{,}9685 hari} & \textbf{0{,}4454} & Terbaik di set uji terkunci \\
\code{B3_lightgbm} (Set Uji Terkunci) & 4{,}8256 hari & 6{,}2458 hari & 0{,}3926 & Artefak tersimpan di registry \\
\end{longtable}

\section{Penolakan Model Pembelian Ulang (\emph{Repeat Purchase}): Kemenangan Disiplin Ilmiah}

Banyak pengembang ML pemula menyembunyikan eksperimen yang gagal agar portofolionya terlihat sempurna. CSIP justru menampilkan penolakan model \code{repeat_purchase} sebagai contoh bagaimana seorang Data Scientist bertanggung jawab menjaga perusahaan agar tidak menghamburkan anggaran pemasaran atas dasar ilusi statistik.

Fakta empiris pada data Olist setelah penyaringan ketat terhadap 962 pembeli ulang sebelum pengiriman dan 2.275 pembeli dalam zona sensor 90 hari menunjukkan:
\begin{itemize}
  \item Tingkat pembelian ulang aktual setelah menerima paket hanyalah sangat kecil (sekitar $\sim 2\%$).
  \item Pada set uji terkunci di masa depan, model \code{B3_lightgbm} mencetak \textbf{ROC-AUC = 0,5021} dengan interval kepercayaan 95\% \textbf{[0,472, 0,531]} dan \textbf{PR-AUC = 0,0202}.
  \item Angka ROC-AUC $0{,}5021$ yang mencakup nilai $0{,}500$ secara statistik \textbf{tidak dapat dibedakan dari melempar koin acak}. Mengapa? Karena data historis Olist hanya merekam satu kali transaksi rata-rata per akun (identifier unik \code{customer_unique_id} sangat jarang melakukan transaksi kedua di periode 2016--2018), serta tidak memuat variabel perilaku penelusuran aplikasi, riwayat promo personal, atau kategori kebutuhan konsumsi berulang.
  \item \textbf{Keputusan Resmi Sistem:} Model prediksi retensi \code{repeat_purchase} diberi status \textbf{DITOLAK / TIDAK DILAYAKKAN PRODUKSI} dan dilarang dipakai menggerakkan kampanye kupon otomatis di Decision Engine. Ini dibuktikan juga pada uji hipotesis \code{H3} di Bab 7.
\end{itemize}

\section{Rekonsiliasi Audit Forensik \& Kejujuran Perbedaan Angka}

Saat proyek ini menjalani peninjauan ketat internal (audit forensik artefak), ditemukan empat ketidakkonsistenan historis antara log awal pengembangan dengan hasil evaluasi set uji terkunci terbaru. Semua diselesaikan dan dibuka apa adanya:

\begin{kejujuranbox}[Empat Temuan Audit Forensik yang Diluruskan Secara Terbuka]
\begin{enumerate}
  \item \textbf{Rekonsiliasi Peringatan Dini Checkout (\code{E1} \code{0,2845 -> 0,1335}):} Dokumen awal pernah mencatat angka PR-AUC sekitar \code{0,2845} akibat penggunaan kombinasi pembagian data atau fitur pasca-transaksi yang belum sepenuhnya terisolasi. Setelah seluruh pos \code{T_0} dikunci murni ke 10 fitur sebelum pengiriman dan dinilai ulang di set uji terkunci pasca-Juni 2018, performa sejati \code{B3_lightgbm} pada Pos \code{T_0} adalah \textbf{PR-AUC = 0,1335} dan \textbf{ROC-AUC = 0,5468}. Sinyal memang ada di atas tebakan acak ($0{,}50$), tetapi masih tergolong lemah untuk pemicu intervensi agresif tanpa info pengiriman. Hal ini didokumentasikan pula dalam \code{CHECKOUT_RISK_FEASIBILITY.md}.
  \item \textbf{Kalibrasi Brier LightGBM \code{E3} (\code{0,1373 -> 0,0762}):} Angka \code{0,1373} di leaderboard berasal dari probabilitas mentah \code{B3_lightgbm} dengan bobot kelas seimbang (\code{class_weight="balanced"}). Karena bobot kelas menaikkan prediksi probabilitas agar menangkap kelas minoritas, kalibrasinya bergeser ke atas. Setelah diterapkan \emph{Platt Scaling} yang dilatih murni di set validasi, Brier Score sejati membaik ke \textbf{0,076177}.
  \item \textbf{Kemenangan \code{B1_ridge} atas \code{B3_lightgbm} pada Regresi Delay:} Pada pemodelan awal di val set, LightGBM ditetapkan sebagai artefak tersimpan (\code{best_delivery_delay_lightgbm.joblib}). Namun pada uji masa depan terkunci, model linear \code{B1_ridge} justru mengungguli LightGBM (\textbf{MAE 4,7509 vs 4,8256 hari}; \textbf{$R^2$ 0,4454 vs 0,3926}). Mengapa? Karena faktor keterlambatan geografis dan perkiraan SLA antarnegara bagian di Brasil berpola sangat aditif (linear), sehingga model linear teregulasi L2 lebih tahan terhadap perubahan musim dibanding pohon gradien. Temuan ini dicatat transparan agar tidak membabi-buta memuja algoritma kompleks.
  \item \textbf{Perbedaan Populasi Uji Antarlaporan (\code{N = 12.571}, \code{12.448}, \code{38.942}):} Penjelasan selisih jumlah baris uji di berbagai berkas laporan: $N = 12.571$ adalah ukuran set uji terkunci \code{ml_leaderboard.csv} pasca-\code{2018-06-30} untuk evaluasi penuh artefak \code{.joblib}; $N = 12.630$ untuk regresi delay (karena pesanan tertentu tidak menuntut nilai ulasan); $N = 38.942$ pada \code{model_slice_analysis.json} mencakup jendela stabilitas lebih lebar sejak \code{2018-03-01}; sedangkan pada audit kohort terdahulu terdapat irisan temporal terpisah sekitar $12.448$ baris. Semua berasal dari sumber tunggal yang terdokumentasi di konfigurasi masing-masing skrip.
\end{enumerate}
\end{kejujuranbox}

\chapter{Uji Hipotesis Statistik, Segmentasi Klaster \& Ketahanan Wilayah}
\chaptermark{Hipotesis H1--H6, Klaster \& Slice Wilayah}

\section{Enam Uji Hipotesis Bisnis dengan Koreksi FDR Benjamini-Hochberg}

\begin{bisnisbox}[Mengapa Menggunakan Koreksi Multi-Uji FDR?]
Jika seorang analis mengetes 20 hal acak dengan batas signifikansi standar $5\%$ ($p < 0{,}05$), secara teori probabilitas paling tidak satu tes tampak ``signifikan secara kebetulan''. Untuk melindungi manajemen dari penemuan palsu, CSIP menuntut setiap p-value mentah pada enam hipotesis bisnisnya disesuaikan menggunakan metode \textbf{Benjamini-Hochberg False Discovery Rate (FDR)} di modul \code{src/csip/hypotheses/runner.py} sebelum menyatakan hipotesis \code{SUPPORTED} atau \code{REJECTED}.
\end{bisnisbox}

Seluruh hasil berikut bersumber langsung dari \code{artifacts/reports/hypothesis_evaluation.json}:

\begin{longtable}{@{} p{0.8cm} >{\raggedright\arraybackslash}p{3.3cm} >{\raggedright\arraybackslash}p{2.6cm} >{\raggedright\arraybackslash}p{1.9cm} >{\raggedright\arraybackslash}p{2.1cm} >{\raggedright\arraybackslash}p{2.5cm} @{}}
\caption{Ringkasan Uji Enam Hipotesis Bisnis H1--H6 dengan Koreksi FDR Benjamini-Hochberg} \\
\toprule
\textbf{Kode} & \textbf{Pertanyaan Bisnis} & \textbf{Uji Statistik} & \textbf{$p_{\text{adj}}$ (FDR)} & \textbf{Ukuran Efek} & \textbf{Keputusan} \\
\midrule
\endfirsthead
\multicolumn{1}{@{}l}{\small\emph{(Tabel lanjutan)}} \\
\toprule
\textbf{Kode} & \textbf{Pertanyaan Bisnis} & \textbf{Uji Statistik} & \textbf{$p_{\text{adj}}$ (FDR)} & \textbf{Ukuran Efek} & \textbf{Keputusan} \\
\midrule
\endhead
\bottomrule
\endfoot
\textbf{H1} & Apakah pesanan dengan rasio ongkir tinggi lebih sering mengeluh? & Welch's $t$-test ($t = -11{,}858$) & $6{,}17 \times 10^{-32}$ & $-0{,}2068$ & \textbf{SUPPORTED} \\
\textbf{H2} & Apakah perbedaan keterlambatan menurunkan skor ulasan secara drastis? & ANOVA $F$-test ($F = 1657{,}035$) & $0{,}0$ ($<10^{-300}$) & $-0{,}0661$ & \textbf{SUPPORTED} \\
\textbf{H3} & Apakah skor kepuasan tinggi mendorong belanja ulang 90 hari? & Two-prop $z$-test ($z = -0{,}175$) & $0{,}8609$ & Rasio $0{,}9964$ (CI $0{,}957$--$1{,}037$) & \textbf{REJECTED} \\
\textbf{H4} & Apakah rata-rata skor kepuasan berbeda bermakna antarwilayah? & ANOVA $F$-test ($F = 29{,}381$) & $3{,}51 \times 10^{-13}$ & $29{,}3812$ & \textbf{SUPPORTED} \\
\textbf{H5} & Apakah sinyal teks \& pengiriman (Pos $T_2$) mengalahkan model tanpa kurir ($T_0$)? & Bootstrap komparatif PR-AUC & $0{,}00012$ & $\Delta \text{PR} = +0{,}1401$ & \textbf{SUPPORTED} \\
\textbf{H6} & Apakah pesanan dengan pembayaran cicilan lebih banyak memiliki keterlambatan atau perbedaan profil nilai? & Welch's $t$-test ($t = -4{,}729$) & $3{,}39 \times 10^{-6}$ & $47{,}3$ & \textbf{SUPPORTED} \\
\end{longtable}

\textbf{Makna Strategis bagi Bisnis dari H1--H6:}
Fakta bahwa \textbf{H2} menghasilkan skor uji $F = 1657{,}035$ menegaskan bahwa kegagalan memenuhi janji estimasi kurir (\code{delivery_delay_days}) adalah pembunuh skor ulasan nomor satu---jauh melampaui faktor harga barang itu sendiri. Sebaliknya, penolakan \textbf{H3} ($p_{\text{adj}} = 0{,}8609$) mematahkan asumsi umum tim pemasaran bahwa \emph{``asal pembeli kasih bintang 5, mereka pasti belanja lagi bulan depan''}. Pada model bisnis marketplace barang tahan lama satu kali beli seperti Olist, menaikkan kepuasan mencegah keluhan/refund dan menjaga reputasi platform, tetapi tidak otomatis menciptakan retensi berulang dalam 90 hari tanpa strategi produk habis pakai atau program loyalitas terpisah.

\section{Segmentasi Pelanggan (\emph{K-Means Clustering}): 3 Profil Operasional Nyata}

Berdasarkan \code{artifacts/reports/clustering_evaluation.json}, algoritma K-Means dieksekusi di atas \textbf{96.470 transaksi} dengan menguji beberapa kandidat jumlah segmen. Komputer memilih \textbf{$K = 3$ segmen} karena menghasilkan kombinasi evaluasi tertinggi:
\begin{itemize}
  \item \textbf{Silhouette Score:} \code{0,3865} (lebih baik daripada kandidat $K=4$ yang bernilai \code{0,3485} pada konfigurasi pipeline final ini).
  \item \textbf{Davies-Bouldin Index:} \code{1,2168}.
  \item \textbf{Calinski-Harabasz Index:} \code{3124,5}.
\end{itemize}

\begin{kejujuranbox}[Sinkronisasi Angka Historis Registry \code{K=4} vs Artefak \code{K=3}]
Di berkas log historis lama \code{artifacts/metadata/MODEL_RESULTS_REGISTRY.csv}, pernah tercatat percobaan eksploratif klaster awal bernilai $K=4$ dengan Silhouette \code{0,4180} pada prapemrosesan sebelum standardisasi ulang fitur. Pada arsitektur pipeline deterministik final di \code{clustering_evaluation.json} ($N=96.470$), metrik valid memilih \textbf{$K=3$ (\code{selected_k=3}, Silhouette \code{0,3865})}. Di setiap presentasi maupun penulisan angka rujukan resmi, gunakan \textbf{3 Segmen Utama ($N=96.470$)} dan jelaskan transisi audit ini jika ditanya asesor.
\end{kejujuranbox}

Profil nyata ketiga segmen tersebut adalah sebagai berikut:

\begin{longtable}{@{} p{2.0cm} >{\raggedleft\arraybackslash}p{1.9cm} >{\raggedleft\arraybackslash}p{1.6cm} >{\raggedleft\arraybackslash}p{2.0cm} >{\raggedright\arraybackslash}p{6.8cm} @{}}
\caption{Profil Tiga Segmen Operasional Pelanggan Hasil K-Means Clustering} \\
\toprule
\textbf{Segmen} & \textbf{Jumlah Pesanan} & \textbf{Porsi (\%)} & \textbf{Ulasan Negatif} & \textbf{Karakteristik Operasional \& Penanganan} \\
\midrule
\endfirsthead
\multicolumn{1}{@{}l}{\small\emph{(Tabel lanjutan)}} \\
\toprule
\textbf{Segmen} & \textbf{Jumlah Pesanan} & \textbf{Porsi (\%)} & \textbf{Ulasan Negatif} & \textbf{Karakteristik Operasional \& Penanganan} \\
\midrule
\endhead
\bottomrule
\endfoot
\textbf{Segmen 0} & 18.987 & 19,68\% & \textbf{15,33\%} & Segmen berisiko gesekan lebih tinggi (ongkir relatif mahal, durasi pengiriman lebih rentan atau kategori sensitif). Risiko keluhan melampaui rata-rata global. \\
\textbf{Segmen 1} & 75.244 & 78,00\% & 11,63\% & Tulang punggung volume normal marketplace: pesanan standar dengan ekspektasi pengiriman relatif terjaga dan profil nilai harian. \\
\textbf{Segmen 2} & 2.239 & 2,32\% & Variatif Khusus & Kelompok ceruk bernilai tinggi / keranjang besar / pola pembayaran dan logistik berdimensi ekstrem. Memerlukan penanganan layanan khusus bernilai tinggi. \\
\end{longtable}

\section{Analisis Ketahanan Subkelompok (\emph{Subgroup Slice Robustness})}

Menguji model hanya pada nilai rata-rata global sering menipu; sebuah model bisa terlihat hebat di pusat kota tetapi runtuh di kota pelabuhan atau daerah pegunungan. Melalui berkas \code{artifacts/reports/model_slice_analysis.json} (evaluasi stabilitas pesanan sejak \code{2018-03-01}, $N = 38.942$ baris), CSIP memecah performa berdasarkan negara bagian, tier keranjang belanja, dan status ketepatan waktu:

\begin{itemize}
  \item \textbf{Metrik Global Irisan Evaluasi ($N = 38.942$):} Klasifikasi Ulasan Negatif meraih \textbf{ROC-AUC = 0,7517} dan \textbf{Brier = 0,1491}; Regresi Keterlambatan mencatat \textbf{MAE = 5,6011 hari}.
  \item \textbf{Kesenjangan Geografis (\code{SP} vs \code{RJ}, \code{MG}, \code{RS}, \code{PR}):} Di pusat ekonomi Sao Paulo (\code{SP}), kesalahan rata-rata tebakan hari keterlambatan amat rendah yaitu \textbf{MAE = 4,2532 hari}. Namun di Rio de Janeiro (\code{RJ}), kesalahan melebar hampir dua kali lipat menjadi \textbf{MAE = 7,724 hari}. Negara bagian penyangga lain seperti Minas Gerais (\code{MG}), Rio Grande do Sul (\code{RS}), dan Parana (\code{PR}) ikut dipantau berkesinambungan karena karakteristik kepadatan kurir dan hambatan logistik lintas daerah sangat berbeda.
  \item \textbf{Kesenjangan Kondisi Ketepatan Waktu (\emph{On-Time} vs \emph{Delayed}):} Ketika paket benar-benar terlambat (\emph{Delayed}), model memisahkan pemicu keluhan sangat tajam dengan \textbf{ROC-AUC = 0,8383}. Sebaliknya, jika paket tiba tepat waktu (\emph{On-Time}), keluhan yang tersisa umumnya disebabkan barang cacat, ukuran salah, atau penjual lambat membalas chat---faktor produk yang belum terekam di fitur kurir sehingga \textbf{ROC-AUC = 0,6522}.
  \item \textbf{Stratifikasi Nilai Keranjang (\emph{Low / Med / High Tiers}):} Ketiga tier belanja dinilai terpisah untuk menjamin model tidak bias mengabaikan pesanan kecil atau salah memperkirakan pembeli barang mewah.
\end{itemize}

\chapter{Kalibrasi Probabilitas, Keterjelasan SHAP, Tata Kelola Klaim \& Simulator Kebijakan}
\chaptermark{Kalibrasi Platt, SHAP \& Skenario S0--S5}

\section{Kalibrasi Probabilitas (Platt Scaling): Mengubah Skor Acak Menjadi Peluang Jujur}

\begin{bisnisbox}[Analogi Bisnis: Skor Mentah vs Persentase Nyata yang Dapat Dipercaya]
Misalkan tim Anda melatih model dengan memberi bobot penalti berat agar sensitif mendeteksi keluhan (\code{class_weight="balanced"}). Hasilnya, model mengeluarkan angka skor \code{0,70} untuk ribuan paket. Jika seorang manajer CS percaya begitu saja bahwa 70 dari 100 paket tersebut pasti protes lalu menyewa tim lembur besar-besaran, perusahaan akan tekor berat---karena kenyataannya dari kelompok berskor 0,70 itu mungkin hanya sekitar 28\% yang benar-benar marah!

\textbf{Kalibrasi Probabilitas} adalah proses menerjemahkan skor mentah yang \emph{``terlalu paranoid''} agar angka \code{0,30} di layar benar-benar bermakna ``sekitar 30 dari 100 transaksi serupa di dunia nyata mengalami keluhan''. Tanpa kalibrasi, perhitungan keuntungan bersih (Expected Net Benefit) dalam rupiah/real di Decision Engine akan rusak total!
\end{bisnisbox}

Di modul CSIP, proses ini menggunakan \textbf{Platt Scaling} (regresi logistik satu dimensi di atas skor log-odds keluaran LightGBM) yang \textbf{dilatih khusus pada Set Validasi (April--Juni 2018)}, bukan pada Set Latih dan bukan pada Set Uji. Hasil terukur pada \code{artifacts/reports/model_explainability.json} menunjukkan lompatan keandalan yang luar biasa:
\begin{itemize}
  \item \textbf{Parameter Kalibrasi Terpasang:} Koefisien pengali \code{coef = 0,9020} dan konstanta pergeseran \code{intercept = -1,8089}. Pergeseran negatif sebesar $-1{,}8089$ inilah yang mengoreksi optimisme/keparanoidan berlebih dari pembobotan kelas seimbang.
  \item \textbf{Expected Calibration Error (ECE --- Selisih antara Keyakinan Model \& Kejadian Nyata):} Turun drastis 97\%, dari \textbf{0,235486} (tidak terkalibrasi, meleset rata-rata 23,5 poin persentase) menjadi \textbf{0,007015} (hanya meleset 0,7 poin persentase!).
  \item \textbf{Brier Score (Kesalahan Kuadrat Probabilitas):} Membaik dari \textbf{0,137478} menjadi \textbf{0,076177}.
\end{itemize}

\section{Atribusi Fitur Global via SHAP (\emph{15 Kontributor Teratas})}

Manajemen tidak menerima jawaban ``kata AI paket ini berisiko''. Setiap prediksi CSIP dibekali penguraian kontribusi per fitur memakai metode \textbf{TreeSHAP} (\emph{SHapley Additive exPlanations}) yang menghitung sumbangan adil tiap fitur pada sampel uji (\code{N = 800} pada berkas \code{artifacts/reports/model_explainability.json}). Urutan kekuatan pendorong keluhan utama pada model pasca-pengiriman (\code{E3}) mencakup 15 fitur berikut:
\begin{enumerate}
  \item \code{delivery_delay_days} --- Selisih hari terlambat terhadap estimasi kurir (faktor dominan mutlak).
  \item \code{delivery_lead_time_days} --- Total hari tunggu aktual dari bayar sampai paket mendarat.
  \item \code{is_late} --- Status biner keterlambatan paket.
  \item \code{delay_x_freight_ratio} --- Interaksi kekecewaan ganda (ongkir mahal + paket telat).
  \item \code{freight_ratio} --- Rasio ongkos kirim terhadap harga barang.
  \item \code{estimated_delivery_days} --- Panjang janji kurir saat checkout.
  \item \code{price} --- Nilai rupiah/real produk.
  \item \code{freight_value} --- Nominal biaya kirim.
  \item \code{is_interstate} --- Status rute pengiriman lintas negara bagian.
  \item \code{order_item_count} --- Banyaknya item di dalam pesanan.
  \item \code{product_weight_g} --- Berat fisik produk yang memengaruhi penanganan logistik.
  \item \code{high_value_late_flag} --- Sinyal pesanan mahal yang mengalami keterlambatan.
  \item \code{customer_state_te} --- Pola risiko wilayah penerima (Target Encoded di data latih).
  \item \code{seller_state_te} --- Pola risiko wilayah asal gudang penjual.
  \item \code{payment_installments} --- Jumlah cicilan pembayaran pembeli.
\end{enumerate}

\section{Hierarki Tata Kelola Klaim L1, L2, dan L3 (\emph{Claim Governance})}

Salah satu ciri kedewasaan ilmiah CSIP adalah aturan keras yang membatasi bagaimana tim bisnis dan aplikasi UI boleh menceritakan temuan model (\code{KNOWN_LIMITATIONS.md}):
\begin{itemize}
  \item \textbf{Level L1 --- Bukti Prediktif \& Asosiatif Murni:} Seluruh skor ROC-AUC, PR-AUC, nilai Brier, SHAP ranking, dan hasil uji $t$/ANOVA/chi-square yang terukur pada data historis. Diizinkan penuh dinyatakan sebagai fakta prediktif.
  \item \textbf{Level L2 --- Simulasi Terkalibrasi \& Proyeksi Observasional:} Hasil perhitungan proyeksi keuntungan bersih di Decision Engine dan simulasi skenario makro VAR Ridge (S0--S5). Wajib menyertakan kalimat penjelas bahwa angka merupakan proyeksi model historis dengan asumsi tingkat efektivitas intervensi dan kepatuhan kuota, bukan bukti empiris pasca-peluncuran.
  \item \textbf{Level L3 --- Bukti Kausal Tervalidasi Lapangan (\emph{Production Causal Impact}):} Klaim definitif berbunyi \emph{``tindakan CSIP menurunkan komplain secara nyata sebesar X\% di dunia nyata''}. Dalam versi rilis \code{v1.1.0} saat ini, \textbf{status L3 belum tercapai} karena Olist merupakan arsip historis statis sehingga belum ada eksperimen pengacakan terkontrol (A/B Test 80/20) langsung di lingkungan produksi nyata. Celah L2-ke-L3 ini tertulis gamblang dan menjadi agenda fase eksperimen lapangan berikutnya.
\end{itemize}

\section{Simulator Kebijakan Makro \& Enam Skenario Strategis (S0--S5)}

Di modul \code{artifacts/reports/world_model_evaluation.json} dan \code{artifacts/reports/scenario_simulation_results.json}, model seri waktu VAR Ridge dilatih pada \textbf{65 minggu pertama} lalu diuji keluar-sampel (\emph{out-of-sample backtest}) pada \textbf{22 minggu terakhir} (total 87 minggu), serta lulus uji pemulihan parameter sintetis dengan \code{max_rel_error = 0,0160} ($1{,}6\%$).

Setelah terverifikasi, enam skenario simulasi kebijakan dijalankan untuk membantu manajemen memilih strategi terbaik sebelum membakar anggaran:

\begin{longtable}{@{} p{0.9cm} >{\raggedright\arraybackslash}p{3.9cm} >{\raggedleft\arraybackslash}p{1.9cm} >{\raggedleft\arraybackslash}p{2.0cm} >{\raggedright\arraybackslash}p{5.3cm} @{}}
\caption{Simulasi Enam Skenario Kebijakan Makro S0--S5 (VAR Ridge 87 Minggu)} \\
\toprule
\textbf{Kode} & \textbf{Nama Skenario Kebijakan} & \textbf{Prediksi Keluhan} & \textbf{Rata-Rata Delay} & \textbf{Interpretasi Bisnis \& Efisiensi} \\
\midrule
\endfirsthead
\multicolumn{1}{@{}l}{\small\emph{(Tabel lanjutan)}} \\
\toprule
\textbf{Kode} & \textbf{Nama Skenario Kebijakan} & \textbf{Prediksi Keluhan} & \textbf{Rata-Rata Delay} & \textbf{Interpretasi Bisnis \& Efisiensi} \\
\midrule
\endhead
\bottomrule
\endfoot
\textbf{S0} & \emph{Baseline} (Tanpa Intervensi Baru) & 9,96\% & $-9{,}08$ hari & Status kuo operasi normal (paket rata-rata tiba 9 hari sebelum tenggat maksimal). \\
\textbf{S1} & Subsidi Ongkir Universal 15\% & 7,46\% & $-8{,}35$ hari & Keluhan turun ke 7,46\%, tetapi sangat boros karena menyubsidi seluruh pembeli yang sebenarnya puas. \\
\textbf{S2} & Pembenahan SLA Kurir (\emph{Carrier Overhaul}) & 9,05\% & \textbf{$-19{,}96$ hari} & Memangkas keterlambatan paling ekstrem, sangat efektif menggeser janji pengiriman lebih konservatif \& aman. \\
\textbf{S3} & Subsidi Terarah Risiko Gesekan (\emph{Targeted Friction Subsidy}) & 7,87\% & $-12{,}31$ hari & Hanya menyubsidi paket ber-rasio ongkir tinggi \& jarak jauh. Hampir menyamai hasil S1 dengan pemborosan anggaran jauh lebih sedikit! \\
\textbf{S4} & Guncangan Makro / Mogok Kurir (\emph{Macro Shock}) & \textbf{14,95\%} & \textbf{$+4{,}54$ hari} & Simulasi stres sistem: saat kurir mogok atau musim puncak memburuk, rata-rata pengiriman meleset $+4{,}5$ hari dan keluhan meledak hampir 15\%. \\
\textbf{S5} & Blitz Gabungan (\emph{Targeted Subsidy + SLA}) & \textbf{6,08\%} & $-19{,}25$ hari & Kombinasi subsidi terarah + perbaikan SLA menurunkan potensi keluhan hingga 6,08\% (penurunan relatif hampir 39\% terhadap S0). \\
\end{longtable}


% ==============================================================================
% BAGIAN III: DECISION ENGINE, CSV STUDIO, MLOPS & PENGUJIAN
% ==============================================================================
\part{Decision Engine, Universal CSV Studio, MLOps \& Pengujian}

\chapter{Mesin Keputusan Bisnis (\emph{Decision Engine}) \& Desain Eksperimen Uji A/B}
\chaptermark{Decision Engine \& Uji A/B 80/20}

\begin{bisnisbox}[Analogi Bisnis: Mengubah Lampu Merah Menjadi Perintah Kerja Nyata]
Model Machine Learning biasa berhenti setelah menampilkan angka di layar, misalnya \emph{``Pesanan \#8821 memiliki probabilitas keluhan 0,68''}. Angka seperti itu sering dibiarkan manajer karena bingung harus bertindak apa: dikirim tim CS? Dikasih voucher? Di-eskalasi ke manajer gudang? Dan kalau ada 10.000 paket berisiko hari ini sementara tim CS hanya sanggup menelepon 500 orang, siapa yang didahulukan?

Di situlah peran \textbf{Business Decision Engine} (\code{src/csip/decision_engine/}). Modul ini bertindak seperti kepala operasi virtual yang membaca probabilitas terkalibrasi dari model ML, menimbangnya dengan batas kapasitas kerja tim lapangan (\emph{capacity budget}), memeriksa batasan klaim ilmiah (L1 vs L2), lalu menerbitkan surat perintah kerja terstruktur untuk setiap pesanan atau segmen.
\end{bisnisbox}

\section{Arsitektur \texttt{BusinessDecisionEngine} \& 15 Field Rekomendasi Terstruktur}

Kode pada \code{src/csip/decision_engine/engine.py}, \code{rules.py}, dan \code{schemas.py} melarang keras keluaran keputusan yang bersifat ``kotak hitam tanpa pertanggungjawaban''. Setiap rekomendasi yang keluar ke UI maupun endpoint API wajib memvalidasi kontrak skema Pydantic bernama \code{DecisionRecommendation} dengan \textbf{15 field wajib}:

\begin{longtable}{@{} >{\raggedright\arraybackslash}p{4.6cm} >{\raggedright\arraybackslash}p{2.8cm} >{\raggedright\arraybackslash}p{8.0cm} @{}}
\caption{Skema Lima Belas Field Wajib Rekomendasi pada Business Decision Engine} \\
\toprule
\textbf{Nama Field Skema} & \textbf{Kelompok} & \textbf{Fungsi Operasional \& Audit} \\
\midrule
\endfirsthead
\multicolumn{1}{@{}l}{\small\emph{(Tabel lanjutan)}} \\
\toprule
\textbf{Nama Field Skema} & \textbf{Kelompok} & \textbf{Fungsi Operasional \& Audit} \\
\midrule
\endhead
\bottomrule
\endfoot
\code{decision_id} & Identitas & ID pelacakan unik rekomendasi untuk audit mundur pasca-tindakan. \\
\code{business_problem} & Konteks Inti & Masalah bisnis spesifik yang sedang diselesaikan (mis. risiko ulasan 1--2 bintang). \\
\code{relevant_entity} & Sasaran & Identitas pesanan, wilayah negara bagian, penjual, atau segmen pelanggan terdampak. \\
\code{prediction_or_evidence} & Bukti Kuantitatif & Angka probabilitas terkalibrasi atau temuan uji statistik (mis. H2/H4). \\
\code{priority_level} & Penentuan Urutan & Tingkat urgensi (\code{CRITICAL}, \code{HIGH}, \code{MEDIUM}, \code{LOW}) berdasar nilai risiko vs kuota tim. \\
\code{contributing_factors} & Keterjelasan & Daftar pendorong utama dari SHAP / interaksi fitur (mis. \code{delay_x_freight_ratio}). \\
\code{recommended_action} & Tindakan Konkret & Langkah intervensi spesifik bagi staf (proaktif telepon, kupon tebus, atau revisi kurir). \\
\code{responsible_function} & Pemilik Tugas & Departemen eksekutor: Tim CS, Operasi Logistik, atau Manajemen Penjual. \\
\code{expected_operational_objective} & Target Hasil & Efek operasional yang diincar (mis. meredam eskalasi keluhan publik). \\
\code{measurement_criteria} & KPI Keberhasilan & Metrik terukur untuk menilai sukses-tidaknya eksekusi setelah berjalan. \\
\code{limitations_and_uncertainty} & Kejujuran Ilmiah & Pengungkapan terbuka rentang ketidakpastian, selang kepercayaan, serta sifat historis data. \\
\code{recommendation_status} & Status Kelayakan & Menandai apakah tindakan siap dijalankan, ditahan (\emph{holdout}), atau khusus riset. \\
\code{evidence_type} & Tata Kelola Klaim & Klasifikasi sumber bukti (\code{L1_PREDICTIVE} vs \code{L2_CALIBRATED_SIMULATION}). \\
\code{estimated_cost_tier} & Kendali Biaya & Kategori beban anggaran tindakan (tanpa biaya/CS internal vs kompensasi finansial). \\
\code{created_at} & Cap Waktu & Stempel waktu penerbitan rekomendasi dalam format UTC ISO-8601. \\
\end{longtable}

\section{Kalibrasi Ambang Batas Kapasitas (\emph{Capacity Threshold Calibration})}

Mengapa Decision Engine tidak asal mematok ambang batas klasik $0{,}50$? Karena di data Olist, pesanan berprobabilitas terkalibrasi $\ge 0{,}50$ jumlahnya sangat sedikit, sedangkan jika diturunkan ke $0{,}15$ jumlah kasus melebihi kapasitas staf CS. Oleh sebab itu:
\begin{itemize}
  \item Sistem mengatur \emph{threshold} secara dinamis sesuai anggaran kapasitas tim CS harian (misalnya target menangani $5\%$ atau $10\%$ pesanan dengan probabilitas terkalibrasi tertinggi yang lolos nilai keranjang minimal).
  \item Keputusan tetap menjaga fungsi prioritas biaya: intervensi berbiaya voucher hanya dibuka pada pesanan dengan interaksi risiko ongkir dan nilai pesanan memadai agar biaya subsidi tidak melebihi nilai pesanan itu sendiri.
  \item Untuk model yang ditolak secara empiris---seperti \code{repeat_purchase} (ROC-AUC \code{0,5021})---aturan di \code{rules.py} memblokir otomatis setiap pemicu rekomendasi pemasaran otomatis guna mencegah pemborosan anggaran retensi.
\end{itemize}

\section{Rancangan Uji A/B Terkontrol 80/20 Menuju Bukti Kausal Level L3}

Sebagaimana ditegaskan dalam \code{KNOWN_LIMITATIONS.md}, proyek CSIP v1.1.0 bersikap jujur bahwa bukti lapangan Level L3 belum dapat diklaim di atas arsip statis Olist 2016--2018. Sebagai jembatan eksekusi profesional saat diterapkan di toko produksi nyata, Decision Engine menyediakan blueprint pengacakan eksperimental \textbf{80/20 Holdout}:
\begin{enumerate}
  \item \textbf{Kelompok Perlakuan (80\% Teratas di Zona Risiko):} Pesanan berprioritas \code{CRITICAL}/\code{HIGH} diarahkan menerima tindakan proaktif dari tim CS sesuai daftar kerja \code{DecisionRecommendation}.
  \item \textbf{Kelompok Kontrol Acak (20\% Ditahan / \emph{Holdout}):} Dua puluh persen pesanan berprofil risiko yang sama persis \emph{sengaja tidak diintervensi ekstra} dan menjalani layanan standar.
  \item \textbf{Evaluasi Kausal Bersih:} Setelah minimal satu siklus 30--60 hari penuh, selisih nyata angka ulasan negatif dan biaya bersih antara kelompok 80\% dan 20\% diuji dengan uji proporsi statistik serta estimasi margin keuntungan bersih (\emph{Net Benefit}). Hanya jika kelompok perlakuan menang signifikan secara statistik barulah sistem berhak mengklaim status bukti kausal \textbf{L3}.
\end{enumerate}

\chapter{Universal CSV Studio (\texttt{ml\_engine.py} \& \texttt{semantic\_profiler.py})}
\chaptermark{Universal CSV Studio (11 Sektor)}

\begin{bisnisbox}[Dari Satu Dataset E-Commerce ke Segala Tabel Bisnis CSV/Excel]
Pertanyaan paling tajam yang sering muncul adalah: \emph{``Apakah CSIP hanya bisa membaca data Olist Brasil, lalu mati kutu kalau dikasih file CSV laporan pabrik, data HR karyawan, atau catatan medis pasien dari rumah sakit?''}

Jawabannya adalah \textbf{TIDAK}. Sejak rilis v1.1.0, CSIP dilengkapi otak kedua bernama \textbf{Universal CSV Studio} yang tertanam di \code{src/csip/models/ml_engine.py} serta modul \code{src/csip/data/semantic_profiler.py}. Pengunggah awam cukup menaruh berkas CSV atau Excel (.xlsx) apa pun, lalu dalam hitungan detik CSIP mengenali sektor industrinya secara semantik, menyapu kolom bocor target, melatih ulang model yang sesuai langsung pada dataset baru tersebut, dan menyajikan simulasi interaktif ``Bagaimana Jika?'' berbahasa Indonesia natural.
\end{bisnisbox}

\section{Leksikon Semantik 11 Sektor Industri \& Inferensi Dinamis Zero-Shot}

Modul \code{src/csip/data/semantic_profiler.py} bertugas menjembatani nama kolom teknis yang kaku (seperti \code{freight_ratio}, \code{s6}, atau \code{biaya_operasional}) menjadi bahasa manusia yang hangat, lengkap dengan kata kerja arah perubahan (naik/turun) dan satuannya (\code{Rp}, \texttt{\$}, \texttt{\%}, \code{hari}, \code{km}, \code{mg/dL}, \code{°C}).

Cakupan leksikon semantik bawaan dalam \code{NATURAL_LEXICON} meliputi \textbf{11 sektor industri utama}:
\begin{enumerate}
  \item \textbf{Pasar Keuangan, Bursa Saham \& Kripto (OHLCV):} Harga penutupan, return harian, volatilitas harian, jarak rata-rata 7/30 hari, dorongan beli vs jual, volume transaksi.
  \item \textbf{Keuangan Bisnis, Akuntansi \& Perbankan Umum:} Laba bersih, omzet pendapatan, OPEX, HPP/COGS, margin kotor, beban pajak, arus kas bersih, plafon kredit, suku bunga.
  \item \textbf{E-Commerce, Ritel \& Toko Online:} Skor kepuasan pembeli, nilai belanja, harga barang, ongkir, rasio ongkir, keterlambatan kirim, lama tunggu, tenor cicilan, diskon produk.
  \item \textbf{SDM, Manajemen Tenaga Kerja \& HR:} Gaji, masa kerja, jam kerja mingguan/lembur, kepuasan karyawan, skor kinerja KPI, hari absensi, jam pelatihan.
  \item \textbf{Pemasaran Digital, Periklanan \& Retensi:} Tingkat konversi (CVR), rasio klik (CTR), jumlah klik, tayangan, biaya akuisisi (CAC), tingkat churn, tingkat retensi.
  \item \textbf{Properti, Bangunan \& Real Estate:} Harga properti, luas bangunan, luas tanah, kamar tidur/mandi, usia bangunan, jarak ke pusat kota.
  \item \textbf{Logistik, Rantai Pasok \& Transportasi:} Jarak tempuh km, waktu tempuh jam, konsumsi BBM liter, bobot muatan kg.
  \item \textbf{Manufaktur, Mesin, Sensor \& IoT:} Suhu mesin (\code{°C}), tekanan hidrolik (\code{psi/bar}), getaran poros (\code{Hz}), daya listrik (\code{kWh}), tingkat cacat produksi (\texttt{\%}), waktu mogok mesin (\code{jam}).
  \item \textbf{Kesehatan, Medis \& Gaya Hidup:} Indeks progresi risiko penyakit, BMI, tekanan darah (\code{mmHg}), kadar gula darah (\code{mg/dL}), kolesterol total/LDL/HDL, usia pasien.
  \item \textbf{Mutu Produk, Bahan Kimia \& Komposisi:} Peringkat mutu/grade, kadar flavonoid, kadar alkohol, kepekatan warna, tingkat keasaman malat, kadar proline.
  \item \textbf{Edukasi, Akademik \& Kemahasiswaan:} Nilai ujian/IPK, jam belajar mandiri, persentase tingkat kehadiran kelas.
\end{enumerate}

Bagaimana jika pengguna mengunggah kolom baru dengan nama lokal perusahaan yang belum terdaftar di kamus di atas (misalnya \code{persen_retur_daerah} atau \code{modal_kerja_idr})? Fungsi \code{infer_dynamic_semantics} melakukan \textbf{inferensi dinamis zero-shot}: membersihkan tanda pemisah (\code{_}, \code{-}), memetakan satuan otomatis lewat pola kata kunci (\code{pct}/\code{ratio} $\rightarrow$ \texttt{\%}, \code{idr}/\code{gaji} $\rightarrow$ \code{Rp}, \code{km} $\rightarrow$ \code{km}, \code{celsius}/\code{temp} $\rightarrow$ \code{°C}, hingga satuan cacah \code{kali}/\code{unit}), dan merangkai kalimat narasi yang tetap terdengar natural untuk tim non-teknis.

\section{Mekanisme Pengaman Otomatis di \texttt{ml\_engine.py}}

Saat pengguna mengunggah file CSV sembarang, risiko terbesar bukan sekadar akurasi rendah, melainkan \textbf{model menipu pengguna karena fitur bocor (\emph{data leakage}) atau error karena format berantakan}. Modul \code{src/csip/models/ml_engine.py} menanamkan enam lapis sabuk pengaman otomatis:

\begin{teknisbox}[Enam Lapis Sabuk Pengaman Universal CSV Engine (\texttt{ml\_engine.py})]
\begin{enumerate}
  \item \textbf{Batas Sampel Ramah RAM 8 GB (\code{sample_cap = 3500}):} Jika file CSV berisi puluhan atau ratusan ribu baris, mesin otomatis mengambil sampel acak terstratifikasi/representatif maksimal \code{3.500 baris} agar proses pelatihan + validasi silang selesai dalam hitungan detik di laptop standar tanpa membekukan RAM.
  \item \textbf{Deteksi Otomatis Tipe Masalah (\emph{Binary Classification vs Regression}):} Sistem memeriksa kolom target yang dipilih pengguna. Jika hanya memuat dua nilai diskrit bersih $\{0, 1\}$, jalur klasifikasi biner aktif; jika kontinu, jalur regresi aktif otomatis.
  \item \textbf{Split-Before-Transform (\code{test_size = 0.25}):} Data dipecah menjadi 75\% latih dan 25\% uji \emph{sebelum} transformasi statistik apa pun dihitung. Khusus dataset mikro ekstrem (\code{N < 8}), sistem memicu penanda jujur \code{IN_SAMPLE_FALLBACK_SMALL_N} alih-alih mengklaim skor generalisasi palsu.
  \item \textbf{Radar Anti-Kebocoran \& Multikolinearitas Khusus Data Latih:}
    \begin{itemize}
      \item Imputasi nilai kosong memakai \code{SimpleImputer(strategy="median")} yang di-\code{fit} \emph{hanya pada data latih}.
      \item Fitur yang memiliki korelasi terlalu dekat dengan target (\textbf{$> 0{,}992$}) otomatis dibuang karena hampir pasti merupakan salinan bocor dari target (misalnya kolom \code{total_bayar} saat menebak \code{harga_final} tanpa pajak).
      \item Fitur yang saling kembar antarsesama fitur (\textbf{inter-feature correlation $> 0{,}94$}) disaring di sisi latih agar koefisien dan peringkat faktor tidak terdistorsi.
    \end{itemize}
  \item \textbf{Validasi Silang Berlapis (\code{StratifiedKFold} / \code{KFold} $\le 5$ Lipatan):} Model tidak hanya diuji satu kali pecah, melainkan divalidasi lintas lipatan untuk mengukur kestabilan error di berbagai irisan sampel.
  \item \textbf{Ketidakpastian \code{prediction_interval_90} \code{[p5, p95]} \& Simulasi Kurva Terlabel Jelas:} Setiap prediksi regresi menyertakan pita batas 90\% (\code{p5} hingga \code{p95}) agar pembaca memahami margin melesetnya. Saat pengguna menggeser tuas simulasi, kurva sensitivitas 70 titik dengan penyebaran Gaussian ($\sigma = 2{,}0$) diberi penanda status jujur \code{HYPOTHETICAL_SENSITIVITY}---menandakan efek asosiatif ``jika fitur ini bergerak'', bukan klaim kausal eksperimental.
\end{enumerate}
\end{teknisbox}

\chapter{MLOps: Detektor Drift, Siklus Champion-Challenger, FastAPI \& Kokpit Streamlit}
\chaptermark{MLOps, Drift, FastAPI \& CSIP Studio}

\section{Mesin Pemantau Drift (\texttt{DriftEngine}) \& Pelatihan Ulang Berjendela}

Model yang hebat pada bulan Januari perlahan bisa membusuk di bulan Desember karena perilaku pasar berubah (\emph{Covariate Shift}) atau hubungan antara keterlambatan dan kemarahan bergeser (\emph{Concept Drift}). Modul \code{src/csip/monitoring/drift_engine.py} memantau tiga dimensi kesehatan model secara ringan di CPU:
\begin{enumerate}
  \item \textbf{Pergeseran Fitur (\emph{Covariate Shift}):} Memakai \textbf{Population Stability Index (PSI)} (ambang stabil ditetapkan pada \code{PSI_THRESHOLD_STABLE = 0,10}), uji non-parametrik \textbf{Kolmogorov-Smirnov (KS)}, dan \textbf{Wasserstein Distance} (\emph{Earth Mover's Distance}) untuk mengukur seberapa jauh distribusi data masuk bergeser dari distribusi data latih dasar.
  \item \textbf{Pergeseran Prediksi (\emph{Prediction Drift}):} Memakai \textbf{Jensen-Shannon Divergence (JSD)} untuk mendeteksi bila model tiba-tiba menyemburkan distribusi probabilitas risiko yang janggal dibanding periode referensi.
  \item \textbf{Peluruhan Konsep (\emph{Concept Drift}):} Setelah label ulasan aktual baru masuk, sistem memantau penurunan PR-AUC, F1-Score, serta naiknya kesalahan kalibrasi Brier Score.
\end{enumerate}

Ketika drift terdeteksi, modul \code{src/csip/monitoring/retrainer.py} (\code{SlidingWindowRetrainer}) dan \code{src/csip/models/registry.py} (\code{ModelLifecycleManager}) menjalankan protokol pergantian model \textbf{Champion-Challenger} secara ketat. Model penantang (\emph{Challenger}) hasil pelatihan ulang pada jendela waktu geser terbaru tidak boleh langsung menggantikan model lama (\emph{Champion}) tanpa syarat. Ia wajib membuktikan:
\begin{itemize}
  \item Kenaikan minimal \textbf{$+2\%$ PR-AUC lift} (memicu kandidat retraining) dan \textbf{$+1\%$ lift} pada evaluasi kelayakan promosi di gerbang lifecycle.
  \item Keterjagaan kalibrasi dengan selisih kesalahan Brier tidak memburuk melebihi toleransi ($\le 5\%$).
  \item Penguncian identitas artefak melalui hash kriptografis \textbf{SHA-256} sehingga jika versi baru mengalami kejanggalan pasca-peluncuran, operator dapat memutar balik (\emph{rollback}) ke artefak Champion sebelumnya secara deterministik.
\end{itemize}

\section{Arsitektur Layanan API (\texttt{FastAPI} --- 7 Endpoint Produksi)}

Untuk integrasi ke sistem backend toko, ERP gudang, maupun bot layanan pelanggan, berkas \code{src/csip/api/server.py} menyediakan REST API FastAPI dengan middleware observabilitas HTTP dan \textbf{7 endpoint aktif}:

\begin{longtable}{@{} p{1.5cm} >{\raggedright\arraybackslash}p{4.7cm} >{\raggedright\arraybackslash}p{9.2cm} @{}}
\caption{Daftar Tujuh Rute HTTP Aktif Layanan Mikro FastAPI} \\
\toprule
\textbf{Metode} & \textbf{Rute Endpoint} & \textbf{Fungsi \& Kegunaan Integrasi Sistem} \\
\midrule
\endfirsthead
\multicolumn{1}{@{}l}{\small\emph{(Tabel lanjutan)}} \\
\toprule
\textbf{Metode} & \textbf{Rute Endpoint} & \textbf{Fungsi \& Kegunaan Integrasi Sistem} \\
\midrule
\endhead
\bottomrule
\endfoot
\code{GET} & \code{/} & Halaman sambutan \& ringkasan identitas versi layanan CSIP. \\
\code{GET} & \code{/health} & Pemeriksaan detak jantung proses server (\emph{liveness probe}). \\
\code{GET} & \code{/ready} & Pemeriksaan kesiapan model \& artefak terkalibrasi (\emph{readiness probe}). \\
\code{POST} & \code{/predict/batch} & Menilai sekumpulan transaksi sekaligus \& mengembalikan probabilitas terkalibrasi + interval/risiko. \\
\code{POST} & \code{/decisions/recommend} & Mengubah skor risiko menjadi objek \code{DecisionRecommendation} berisi 15 field kerja lengkap. \\
\code{POST} & \code{/monitor/drift} & Menerima cuplikan data produksi terkini \& menghitung PSI, KS, Wasserstein, dan status JSD. \\
\code{GET} & \code{/models/registry} & Menampilkan katalog model terdaftar, versi aktif, \& jejak hash SHA-256 Champion/Challenger. \\
\end{longtable}

\section{Anatomi Kokpit \emph{CSIP Studio} (\texttt{app/streamlit\_app.py})}

Antarmuka pengguna pada \code{app/streamlit_app.py} dirancang dengan prinsip \emph{single-screen cognitive clarity} agar pimpinan non-teknis dapat memahami inti cerita dalam 10 detik tanpa tersesat di dalam belasan tab yang membingungkan:

\begin{enumerate}
  \item \textbf{Sumber Data Fleksibel di Bilah Samping / Atas:} Pengguna dapat langsung mengunggah berkas CSV/Excel miliknya sendiri atau memilih satu dari \textbf{3 Kasus Bawaan}:
    \begin{itemize}
      \item \emph{Kasus 1: Toko Online \& Kepuasan Pembeli (Data Olist)}
      \item \emph{Kasus 2: Risiko Kesehatan \& Gaya Hidup (Data Medis)}
      \item \emph{Kasus 3: Kualitas Produk \& Komposisi Bahan (Data Mutu)}
    \end{itemize}
  \item \textbf{Langkah 1 --- Pemilihan Target \& Terjemahan \emph{``Artinya begini''}:} Pengguna memilih indikator yang ingin dioptimalkan (dropdown \emph{``Yang ingin dioptimalkan''}). Sistem langsung merangkum arti kolom tersebut dalam bahasa sehari-hari di bawah label \textbf{``Artinya begini''}.
  \item \textbf{Langkah 2 --- Panel Tuas Simulasi (\emph{Top-3 Sliders + Setel ulang}):} Tiga faktor paling berpengaruh berdasarkan model disajikan sebagai penggeser interaktif (\emph{Tuas simulasi}) disertai tombol cepat \textbf{``Setel ulang''} untuk kembali ke posisi baseline awal.
  \item \textbf{Tiga Kartu KPI Bento:} Menampilkan ringkasan metrik baseline vs hasil simulasi serta indikator reliabilitas/ketidakpastian tepat di puncak layar utama.
  \item \textbf{Empat Lensa Analisis (\emph{Lensa tampilan} Radio Selector):}
    \begin{itemize}
      \item \code{1. Kurva Simulasi dan Titik Kritis}: Grafik respons non-linear ketika tuas digeser.
      \item \code{2. Peringkat Faktor Penentu}: Ranking kekuatan pendorong faktor dari model.
      \item \code{3. Sebaran Kelompok Data}: Peta distribusi segmen/populasi data aktual.
      \item \texttt{4. Observatorium Drift \& Stabilitas}: Kondisi kestabilan distribusi dan kesehatan evaluasi data/model.
    \end{itemize}
  \item \textbf{Mesin Narasi \emph{Nara}:} Di bagian bacaan eksekutif, komponen \textbf{Nara} merangkai temuan angka dari kartu Bento dan pergeseran tuas menjadi paragraf rekomendasi berbahasa Indonesia yang mengalir wajar dan jujur terhadap batasan ketidakpastian.
\end{enumerate}

\section{Katalog 14 Notebook Riset \& 22 Skrip Reproduksi Otomatis}

Seluruh alur eksplorasi hingga rilis produk tersusun rapi menjadi dua jalur yang saling menyinkronkan: jalur interaktif \textbf{14 Jupyter Notebook} di direktori \code{notebooks/} (untuk pembelajaran, visualisasi, dan peninjauan langkah demi langkah) serta jalur otomatis \textbf{22 skrip Python} di \code{scripts/} + \textbf{4 alat verifikasi} di \code{tools/} (untuk eksekusi batch yang dapat diaudit hash-nya).

{\small
\begin{longtable}{@{} >{\raggedright\arraybackslash}p{7.8cm} >{\raggedright\arraybackslash}p{8.0cm} @{}}
\caption{Katalog Empat Belas Jupyter Notebook Eksperimen dan Reproduksi Riset} \\
\toprule
\textbf{Berkas Notebook (\texttt{notebooks/})} & \textbf{Cakupan Eksperimen \& Artefak Dihasilkan} \\
\midrule
\endfirsthead
\multicolumn{1}{@{}l}{\small\emph{(Tabel lanjutan)}} \\
\toprule
\textbf{Berkas Notebook (\texttt{notebooks/})} & \textbf{Cakupan Eksperimen \& Artefak Dihasilkan} \\
\midrule
\endhead
\bottomrule
\endfoot
\code{01_business_and_data_understanding.ipynb} & Perumusan masalah bisnis, estimand, dan profil awal 9 tabel Olist. \\
\code{02_data_validation_and_cleaning.ipynb} & Pengujian kontrak data 4 lapis, sanitasi tipe, dan penguncian grain pesanan. \\
\code{03_exploratory_data_analysis.ipynb} & Eksplorasi distribusi keterlambatan, ongkir, wilayah Brasil, dan anomali ulasan. \\
\code{04_feature_engineering_and_construction.ipynb} & Konstruksi fitur berbasis Pos $T_0, T_1, T_2$ dan penjagaan target leakage. \\
\code{05_behavioral_discovery_unsupervised.ipynb} & Penemuan segmen via K-Means (Silhouette/DB/CH) dan kurva Kaplan-Meier. \\
\code{06_hypothesis_investigation.ipynb} & Eksekusi 6 uji hipotesis statistik H1--H6 dengan koreksi FDR BH. \\
\code{07_supervised_ml_and_competency_lab.ipynb} & Pelatihan \& pembandingan tangga baseline B0--B3 di split temporal. \\
\code{08_neural_network_and_deep_learning.ipynb} & Eksperimen CPU-friendly \code{TabularEmbeddingMLP} dan \code{TextBiLSTM}. \\
\code{09_model_evaluation_and_explainability.ipynb} & Kalibrasi Platt Scaling (ECE/Brier), TreeSHAP 15 fitur, dan irisan wilayah. \\
\code{10_system_and_world_modeling.ipynb} & Pembangunan model runtun waktu mingguan VAR Ridge (87 minggu) \& uji sintetik. \\
\code{11_alternative_world_simulation.ipynb} & Eksekusi enam skenario kebijakan strategis S0--S5 untuk rapat pimpinan. \\
\code{12_uncertainty_quantification.ipynb} & Pengukuran interval kepercayaan bootstrap 500x, Manski bounds, dan pita 90\%. \\
\code{13_model_review_and_failure_analysis.ipynb} & Analisis kegagalan model \code{repeat_purchase}, kesalahan \code{RJ}, dan kasus tepi. \\
\code{14_final_research_report.ipynb} & Ringkasan temuan penelitian akhir, batas klaim L1--L3, dan peta keputusan. \\
\end{longtable}
}

\chapter{Pengujian Otomatis, 12 Gerbang Kualitas, Pemetaan SKKNI \& Panduan Reproduksi}
\chaptermark{Pengujian, 12 Gerbang Rilis \& SKKNI}

\section{Rangkaian 216 Pengujian Otomatis di 39 Berkas \texttt{pytest}}

Kepercayaan terhadap CSIP tidak bersumber dari klaim verbal pengembang, melainkan ditegakkan oleh \textbf{39 berkas modul pengujian} (\code{tests/test_*.py}) dengan total terverifikasi \textbf{216 kasus uji (\code{216 passed})}. Pengujian tersebut mencakup uji kontrak data (\code{test_data_contracts.py}), pencegahan \emph{fan-out} penggabungan tabel (\code{test_cleaning_and_grain.py}), isolasi fitur pos waktu dan ketiadaan kebocoran masa depan, kesesuaian kalibrasi (\code{test_calibration_audit.py}), logika Decision Engine, detektor drift, 7 rute FastAPI (\code{test_api_v11.py}), simulasi layar hidup Streamlit menggunakan \code{AppTest} (\code{test_app_pages.py}), konsistensi SKKNI (\code{test_skkni_consistency.py}), serta pendeteksi klaim berlebihan (\code{test_claims_and_overclaims.py}).

\section{Dua Belas Gerbang Kualitas Rilis (\emph{Release Gates G1--G12})}

Sebelum rilis dinyatakan layak pakai atau dikemas oleh \code{scripts/package_release.py}, skrip penjaga \code{tools/verify_release.py} mewajibkan keduabelas gerbang berikut bernilai \textbf{PASS} tanpa pengecualian dan mencetak log ke \code{reports/verification/gates.json}:

\begin{longtable}{@{} p{1.0cm} >{\raggedright\arraybackslash}p{4.0cm} >{\raggedright\arraybackslash}p{9.8cm} @{}}
\caption{Dua Belas Gerbang Kualitas Rilis CSIP (G1--G12)} \\
\toprule
\textbf{Kode} & \textbf{Nama Gerbang Kualitas} & \textbf{Bukti Kepatuhan \& Cek Eksekutabel (\texttt{tools/verify\_release.py})} \\
\midrule
\endfirsthead
\multicolumn{1}{@{}l}{\small\emph{(Tabel lanjutan)}} \\
\toprule
\textbf{Kode} & \textbf{Nama Gerbang Kualitas} & \textbf{Bukti Kepatuhan \& Cek Eksekutabel (\texttt{tools/verify\_release.py})} \\
\midrule
\endhead
\bottomrule
\endfoot
\textbf{G1} & Single Version Truth & Keseragaman versi paket (\code{1.1.0}) antara \code{pyproject.toml}, \code{csip._version}, dan seluruh file pendamping model (\emph{sidecar JSON}). \\
\textbf{G2} & Zero Failed Tests & Seluruh kasus uji dalam \code{reports/testing/pytest_actual.txt} (216 pengujian) wajib lulus dengan 0 kegagalan. \\
\textbf{G3} & Clean Claims \& Anti-Overclaim & Pemindaian \code{tools/check_claims.py} memastikan nol frasa sesumbar terlarang di seluruh repositori. \\
\textbf{G4} & Preregistration \& Frozen Plan & Berkas \code{config/analysis_plan.yaml} utuh terverifikasi hash SHA-256 (\code{c340dab182cc...}). \\
\textbf{G5} & Valid Rolling-Origin Temporal CV & Validasi silang 4 lipatan jendela mengembang berjalan valid dengan embargo ketat $\ge 90$ hari. \\
\textbf{G6} & Single Evaluation Locked Test & Set uji terkunci dievaluasi satu kali tanpa kebocoran optimasi (audit slice \code{N=12.448, PR-AUC=0,3151 vs base 0,0961}; leaderboard sinkron \code{N=12.571, PR-AUC=0,3061}). \\
\textbf{G7} & Estimand Population & Pengungkapan eksplisit sifat retrospektif estimand M1--M4 serta perhitungan sensitivitas \emph{Manski bounds} (\texttt{12,10\%} vs \texttt{11,44\%}). \\
\textbf{G8} & Tested Dashboard with Provenance & Seluruh komponen UI Streamlit CSIP Studio dan 4 lensa tampilan teruji dengan \code{AppTest} terhubung ke artefak riil. \\
\textbf{G9} & Rerunable Competency Evidence & 100\% berkas bukti 63 KUK SKKNI eksis fisik di penyimpanan dan cocok hash SHA-256-nya. \\
\textbf{G10} & One-Command Complete Reproduction & Eksekusi satu perintah (\code{scripts/reproduce_all.py}) dan audit ke-14 notebook berjalan konsisten. \\
\textbf{G11} & CI Multi-Version Workflow & Konfigurasi alur integrasi berkelanjutan (\code{.github/workflows/ci.yml}) aktif memeriksa gerbang rilis. \\
\textbf{G12} & Honest Competency Provenance & Pengungkapan jujur keterlacakan SKKNI memakai taksonomi kategori status apa adanya. \\
\end{longtable}

\section{Pemetaan Standar Kompetensi SKKNI No. 299 Tahun 2020 (KKNI Level 7)}

Bagi praktisi maupun peserta uji kompetensi skema Ilmuwan Data (\emph{Data Scientist}) LSP TRI (\code{SKM/2905/LSP-TRI/VII/2025}), setiap komponen di CSIP telah dipetakan ke \textbf{11 Unit Kompetensi, 27 Elemen, dan 63 KUK} secara terbuka dan jujur di \code{competency/SKKNI_TRACEABILITY_MATRIX.csv} serta \code{docs/SKKNI_TRACEABILITY_GUIDE.md}:

{\small
\begin{longtable}{@{} >{\raggedright\arraybackslash}p{2.7cm} >{\raggedright\arraybackslash}p{3.9cm} c c >{\raggedright\arraybackslash}p{6.8cm} @{}}
\caption{Pemetaan Sebelas Unit Kompetensi SKKNI No. 299 Tahun 2020 (KKNI Level 7)} \\
\toprule
\textbf{Kode Unit} & \textbf{Judul Unit Kompetensi} & \textbf{Elem.} & \textbf{KUK} & \textbf{Bukti Teknis Utama di CSIP} \\
\midrule
\endfirsthead
\multicolumn{1}{@{}l}{\small\emph{(Tabel lanjutan)}} \\
\toprule
\textbf{Kode Unit} & \textbf{Judul Unit Kompetensi} & \textbf{Elem.} & \textbf{KUK} & \textbf{Bukti Teknis Utama di CSIP} \\
\midrule
\endhead
\bottomrule
\endfoot
\code{J.62DMI00.001.1} & Menentukan Objektif Bisnis & 4 & 9 & \code{docs/methodology/BUSINESS_UNDERSTANDING.md} \\
\code{J.62DMI00.002.1} & Menentukan Tujuan Teknis DS & 2 & 5 & \code{docs/methodology/PREDICTION_CONTRACT.md} \\
\code{J.62DMI00.005.1} & Menelaah Data & 2 & 5 & \code{reports/data/DATA_PROFILING_REPORT.md} \\
\code{J.62DMI00.006.1} & Memvalidasi Data & 3 & 7 & \code{src/csip/data/contracts.py} \\
\code{J.62DMI00.007.1} & Menentukan Objek Data & 2 & 5 & \code{docs/methodology/DATA_DICTIONARY.md} \\
\code{J.62DMI00.008.1} & Membersihkan Data & 3 & 7 & \code{src/csip/cleaning/cleaner.py} \\
\code{J.62DMI00.009.1} & Mengkonstruksi Data (FE) & 2 & 5 & \code{src/csip/features/engineer.py} \\
\code{J.62DMI00.012.1} & Membangun Skenario Model & 2 & 5 & \code{docs/EXPERIMENT_DOCUMENTATION.md} \\
\code{J.62DMI00.013.1} & Membangun Model & 3 & 7 & \code{src/csip/models/trainer.py} \\
\code{J.62DMI00.014.1} & Mengevaluasi Hasil Pemodelan & 2 & 5 & \code{reports/models/MODEL_ROBUSTNESS_AND_SLICE_ANALYSIS.md} \\
\code{J.62DMI00.015.1} & Melakukan Review Pemodelan & 3 & 6 & \code{reports/final/} (Dokumen Audit Akhir) \\
\midrule
\textbf{TOTAL} & \textbf{11 Unit Kompetensi Resmi} & \textbf{27} & \textbf{63} & \textbf{100\% Dapat Ditelusuri Fisik} \\
\end{longtable}
}

Kejujuran status verifikasi 63 KUK di atas dibagi tanpa rekayasa: \textbf{35 KUK (\texttt{55,6\%})} berstatus \code{VERIFIED} (terbukti di kode + tes otomatis + kecocokan SHA-256), \textbf{12 KUK (\texttt{19,0\%})} berstatus \code{ARTIFACT_ONLY} (templat/skema tersedia), \textbf{13 KUK (\texttt{20,6\%})} berstatus \code{CANDIDATE_DEMONSTRATION_REQUIRED} (menuntut penjelasan lisan \& demo langsung di hadapan asesor), dan \textbf{3 KUK (\texttt{4,8\%})} berstatus \code{EXTERNAL_CONFIRMATION_REQUIRED} (KUK 2.2, 3.3, 4.3 Unit 1 yang secara hukum memang memerlukan persetujuan formal dari klien/manajemen eksternal).

\section{Panduan Reproduksi Lengkap Lewat Terminal (CLI)}

Seluruh isi buku ini dapat diverifikasi ulang oleh siapa saja di laptop Windows/Linux/macOS dengan urutan perintah terminal berikut:

\begin{lstlisting}[language=bash, caption={Perintah Reproduksi, Verifikasi, Layanan API, dan Paket Rilis CSIP v1.1.0}]
# 1. Jalankan reproduksi pipeline penuh (pembersihan -> fitur -> pelatihan -> laporan)
python scripts/reproduce_all.py

# 2. Jalankan seluruh rangkaian 216 unit & integrasi test
python -m pytest -q

# 3. Audit integritas klaim bahasa (harus keluar kode 0) & perbarui manifest SHA-256
python tools/check_claims.py
python scripts/generate_sha256sums.py

# 4. Evaluasi seluruh 12 gerbang kualitas rilis G1-G12
python tools/verify_release.py

# 5. Luncurkan antarmuka interaktif CSIP Studio (Streamlit)
streamlit run app/streamlit_app.py

# 6. Luncurkan API produksi FastAPI beserta Swagger UI (/docs)
uvicorn csip.api.server:app --host 127.0.0.1 --port 8000

# 7. Kemas bundel distribusi ZIP siap rilis (dibatasi aman < 100 MB)
python scripts/package_release.py
\end{lstlisting}


% ==============================================================================
% BAGIAN IV: PAKET PRESENTASI MASTER & TANYA-JAWAB PENGUJI
% ==============================================================================
\part{Paket Presentasi Master, Live Demo \& Tanya-Jawab Penguji}

\chapter{Strategi Presentasi Multi-Audiens, Peta Waktu \& Aturan Emas Panggung}
\chaptermark{Strategi Presentasi \& Peta Waktu}

Banyak proyek Data Science yang bagus secara kode justru gagal meyakinkan direksi atau penguji hanya karena pembawaannya berputar-putar di rumus teknis atau sebaliknya, terdengar terlalu mengada-ada tanpa bukti data. Bab ini memetakan strategi presentasi agar satu proyek CSIP v1.1.0 dapat memukau empat tipe audiens yang sangat berbeda dalam ruangan yang sama.

\section{Memahami Siapa yang Duduk di Hadapan Anda}

\begin{longtable}{@{} >{\raggedright\arraybackslash}p{3.3cm} >{\raggedright\arraybackslash}p{4.7cm} >{\raggedright\arraybackslash}p{7.4cm} @{}}
\caption{Peta Kebutuhan Multi-Audiens Presentasi CSIP v1.1.0} \\
\toprule
\textbf{Profil Audiens} & \textbf{Pertanyaan Utama di Kepala Mereka} & \textbf{Senjata Pembuktian CSIP yang Harus Ditonjolkan} \\
\midrule
\endfirsthead
\multicolumn{1}{@{}l}{\small\emph{(Tabel lanjutan)}} \\
\toprule
\textbf{Profil Audiens} & \textbf{Pertanyaan Utama di Kepala Mereka} & \textbf{Senjata Pembuktian CSIP yang Harus Ditonjolkan} \\
\midrule
\endhead
\bottomrule
\endfoot
\textbf{Direksi / Pemilik Bisnis (Non-Teknis)} & ``Apa masalah uang yang selesai, seberapa hemat, dan kapan tim saya bisa memakai aplikasinya?'' & Fokus pada perbandingan skenario S0 (9,96\% keluhan) vs S3 Subsidi Terarah (7,87\%) vs S5 (6,08\%), kokpit \emph{CSIP Studio} satu layar, dan \emph{Decision Engine} prioritas kerja tim CS. \\
\textbf{Manajer Operasional / CS / Logistik} & ``Bagaimana staf saya tahu paket mana yang harus ditangani hari ini jika kuota terbatas?'' & Jelaskan objek kerja \code{DecisionRecommendation} 15 field, kesenjangan kurir Sao Paulo (MAE 4,25 hari) vs Rio de Janeiro (MAE 7,72 hari), dan ambang kapasitas kuota harian. \\
\textbf{Lead Data Scientist / Reviewer Teknis} & ``Apakah skornya hasil kebocoran data masa depan, bagaimana kalibrasinya, dan tahan kah lintas domain?'' & Beberkan isolasi waktu $T_0/T_1/T_2$, embargo 90 hari, penolakan \code{repeat_purchase} (ROC 0,5021), penurunan ECE via Platt Scaling (\code{0,235 -> 0,007}), dan pengaman Universal CSV Engine (\code{> 0,992} drop). \\
\textbf{Asesor Kompetensi SKKNI / KKNI Level 7} & ``Apakah 11 unit kompetensi dan seluruh artefak benar-benar konsisten, jujur, dan dapat diulang?'' & Tunjukkan matriks 11 Unit / 27 Elemen / 63 KUK dengan taksonomi status jujur, 12 gerbang \code{verify_release.py}, dan 216 tes \code{pytest} terverifikasi. \\
\end{longtable}

\section{Peta Waktu Presisi: Format 5 Menit, 15 Menit, dan 45 Menit}

Sesuaikan kedalaman presentasi dengan durasi panggung yang tersedia memakai tabel rencana berikut:

\begin{longtable}{@{} >{\raggedright\arraybackslash}p{2.4cm} >{\raggedright\arraybackslash}p{4.4cm} >{\raggedright\arraybackslash}p{4.4cm} >{\raggedright\arraybackslash}p{4.2cm} @{}}
\caption{Peta Manajemen Waktu Presentasi (5 Menit, 15 Menit, dan 45 Menit)} \\
\toprule
\textbf{Mode Durasi} & \textbf{Babak yang Dipakai} & \textbf{Alokasi Waktu Detail} & \textbf{Fokus Utama} \\
\midrule
\endfirsthead
\multicolumn{1}{@{}l}{\small\emph{(Tabel lanjutan)}} \\
\toprule
\textbf{Mode Durasi} & \textbf{Babak yang Dipakai} & \textbf{Alokasi Waktu Detail} & \textbf{Fokus Utama} \\
\midrule
\endhead
\bottomrule
\endfoot
\textbf{Kilat 5 Menit} (Elevator Pitch) & Babak 1, 3 (ringkas), 6 (skenario), 9--10 (demo singkat), Penutup & Menit 0--1: Masalah; 1--2: Anti-Leakage \& Hasil; 2--4: Demo Tuas CSIP Studio; 4--5: Universal CSV \& Penutup & Mengunci perhatian pada nilai bisnis, kejujuran bukti, dan demo langsung. \\
\textbf{Rapat 15 Menit} (Sidang Proyek) & Babak 1, 2, 3, 5, 6, 8, serta Demo Babak 9, 10, 12 + Penutup & Menit 0--6: Masalah s.d. Hasil \& Penolakan Model; 6--11: Demo Olist + Upload CSV Baru; 11--15: Decision Engine \& Q\&A & Menjawab tuntas sisi bisnis dan disiplin Sains Data sekaligus membuktikan portabilitas CSV. \\
\textbf{Master 45 Menit} (Presentasi Penuh + Asesmen) & \textbf{Seluruh Babak 1--14} + Sesi Tanya-Jawab Bab 16 & Menit 0--20: Babak 1--8 (Pemaparan Induk); 20--33: Babak 9--13 (Live Demo + API + CLI Gates); 33--45: Babak 14 + Q\&A & Pemaparan paripurna tanpa celah untuk ujian akhir, wawancara portofolio senior, atau asesmen SKKNI Level 7. \\
\end{longtable}

\section{Daftar Periksa 5 Menit Sebelum Manggung \& Tujuh Kalimat Haram}

Sebelum mikrofon menyala, pastikan 5 hal teknis sudah berdiri aktif di laptop Anda:
\begin{enumerate}
  \item \textbf{Terminal Verifikasi Siap:} Jalankan \code{python tools/verify_release.py} terlebih dahulu agar status 12 gerbang \code{PASS} dan \code{216 passed} segar di layar terminal terpisah.
  \item \textbf{Tab Browser 1 --- Streamlit CSIP Studio:} Jalankan \code{streamlit run app/streamlit_app.py} dan pastikan halaman awal terbuka mulus pada \emph{Kasus 1: Toko Online \& Kepuasan Pembeli (Data Olist)}.
  \item \textbf{Tab Browser 2 --- FastAPI Swagger UI:} Jalankan \code{uvicorn csip.api.server:app --host 127.0.0.1 --port 8000} dan buka \code{http://127.0.0.1:8000/docs} di tab kedua.
  \item \textbf{Berkas CSV Uji Coba Siap Unggah:} Siapkan satu file CSV bisnis nyata (misalnya sampel HR, penjualan, atau kesehatan) di folder mudah dijangkau untuk didemonstrasikan unggah langsung.
  \item \textbf{Koneksi Internet Tidak Wajib:} Tegaskan di depan hadirin bahwa seluruh aplikasi berjalan 100\% lokal tanpa panggilan API AI cloud apa pun.
\end{enumerate}

\begin{tipbox}[Tujuh Kalimat Terlarang yang Jangan Pernah Diucapkan]
\begin{enumerate}
  \item Jangan pernah berkata \emph{``Model kami pasti menurunkan komplain di toko nyata''} $\rightarrow$ Ganti dengan \emph{``Berdasarkan simulasi terkalibrasi historis L2, skenario target menunjukkan penurunan ke 7,87\%, dan pembuktian dampak lapangan L3 kami rancang lewat uji A/B holdout 80/20.''}
  \item Jangan pernah mengklaim \emph{``Dataset kami sepenuhnya tanpa kebocoran atau sempurna seratus persen''} $\rightarrow$ Katakan \emph{``Sistem menerapkan isolasi pos waktu T0/T1/T2 dan pemblokiran korelasi > 0,992 sehingga tidak ada kebocoran fitur pasca-kejadian yang terdeteksi.''}
  \item Jangan pernah berkata \emph{``Akurasi AI kami 88\%''} $\rightarrow$ Sebut metrik yang benar untuk kelas tak seimbang: \textbf{PR-AUC 0,3061 vs lantai dasar $\sim$0,11--0,12}, \textbf{ROC-AUC 0,6745} pada set uji terkunci, dan kalibrasi \textbf{ECE 0,007}.
  \item Jangan pernah malu atau menyembunyikan \code{repeat_purchase} $\rightarrow$ Justru tonjolkan penolakan model ROC-AUC 0,5021 tersebut sebagai bukti integritas Data Science Anda!
  \item Jangan menyebut \emph{``CSIP butuh GPU server mahal''} $\rightarrow$ Banggakan efisiensi laptop RAM $\le 8$ GB dan CPU-only.
  \item Jangan bilang \emph{``Ini proyek khusus Olist saja''} $\rightarrow$ Jelaskan dua kaki arsitektur: studi mendalam Olist + mesin adaptif Universal CSV Studio untuk 11 sektor industri.
  \item Jangan bilang \emph{``Lampu biru AI sudah 100\% siap menggantikan manajer''} $\rightarrow$ Tekankan konsep \emph{Decision Intelligence}: mesin memberi rekomendasi terurut berserta batasan ketidakpastian, keputusan eksekusi tetap dikendalikan manusia dan anggaran kuota.
\end{enumerate}
\end{tipbox}

\chapter{Skrip Presentasi Induk Kata-demi-Kata (Babak 1 s.d. Babak 8)}
\chaptermark{Skrip Presentasi Induk (Babak 1--8)}

Naskah berikut disusun agar dapat dibaca langsung apa adanya maupun disesuaikan intonasinya. Setiap babak memuat panduan visual layar (\emph{stagebox}), teks ucapan panjang lebar untuk jembatan non-teknis ke teknis (\emph{scriptbox}), dan catatan kontrol panggung (\emph{tipbox}).

% ------------------------------------------------------------------------------
\section{Babak 1 --- Pembukaan, Hook Bisnis \& Tujuan Utama Proyek}

\begin{stagebox}[Panduan Layar \& Durasi --- Babak 1]
\textbf{Durasi Rencana:} 2,5 menit. \quad \textbf{Visual Layar:} Slide Judul CSIP v1.1.0 + Infografis Masalah Biaya E-Commerce (Ongkir Mahal + Keterlambatan $\rightarrow$ Ulasan 1--2 Bintang) + Diagram Ringkas 5 Lapisan Arsitektur.
\end{stagebox}

\begin{scriptbox}[Naskah Bicara Lengkap --- Babak 1]
``Selamat pagi dan terima kasih atas waktu Bapak, Ibu pimpinan, serta rekan-rekan tim teknis maupun asesor sekalian. Hari ini saya ingin membagikan perjalanan rekayasa dan Sains Data dari sebuah platform bernama \textbf{CSIP --- Core Strategy \& Intelligence Platform versi 1.1.0}.

Izinkan saya memulai dengan sebuah masalah sederhana yang sangat mahal di dunia ritel dan e-commerce. Bayangkan Anda menjalankan sebuah marketplace nasional dengan hampir seratus ribu pesanan. Ketika seorang pelanggan membuka kotak paketnya lalu meninggalkan ulasan \textbf{bintang satu atau bintang dua}, sering kali kita baru tahu kejadian itu \emph{setelah nasi menjadi bubur}: reputasi penjual sudah tercoreng di halaman produk, tim CS dibanjiri tiket marah secara reaktif, dan biaya penanganan keluhan membengkak. Di sisi lain, kalau tim keuangan sekadar mengambil langkah panik seperti \emph{membagi subsidi ongkos kirim untuk semua pesanan}, anggaran promosi langsung habis terbakar untuk membiayai 88 persen pembeli yang sejatinya sudah puas tanpa diberi subsidi pun!

Di sinilah letak kesenjangan terbesar antara banyak proyek Machine Learning akademis dan kebutuhan operasional perusahaan nyata. Banyak tim data berhenti setelah membuat grafik korelasi atau mencetak skor akurasi di atas file Jupyter Notebook, tanpa pernah menjawab tiga pertanyaan kritis di meja direksi: \textbf{Pertama}, informasi apa saja yang benar-benar legal dan sudah tersedia secara waktu nyata sebelum prediksi dibuat? \textbf{Kedua}, ketika kapasitas tim CS kita terbatas hanya mampu menelepon sebagian kecil pembeli, paket mana yang didahulukan agar setiap rupiah kompensasi memberi penyelamatan reputasi tertinggi? Dan \textbf{ketiga}, apakah platform ini cukup cerdas untuk menerima data bisnis lain dalam format CSV secara langsung, ataukah ia mati kutu di luar satu dataset latihan?

Platform \textbf{CSIP v1.1.0} dibangun dari nol untuk menjawab ketiganya di bawah satu batasan praktis yang sangat ketat: sistem harus berjalan penuh \textbf{100 persen lokal di laptop standar ber-RAM maksimal 8 GB tanpa memerlukan GPU mahal dan tanpa mengirim satu baris data pelanggan pun ke API AI eksternal di cloud}. Platform ini berdiri di atas dua pilar yang saling mengunci: pertama, \textbf{cetak biru riset mendalam dan auditable pada hampir seratus ribu transaksi riil Olist Brazil} lengkap dengan mesin keputusan serta simulator kebijakan makro; dan kedua, \textbf{Universal CSV Studio}, mesin adaptif yang mampu memprofilkan, memfilter kebocoran data, memodelkan, dan mensimulasikan tabel CSV/Excel dari 11 sektor industri langsung dari satu layar antarmuka.''
\end{scriptbox}

\begin{tipbox}[Kontrol Intonasi \& Transisi]
Berhenti sejenak 1,5 detik setelah kalimat \emph{``membakar anggaran untuk 88 persen pembeli yang sudah puas''}. Tekankan frasa \textbf{100\% lokal, RAM 8 GB, tanpa cloud AI} dengan suara mantap karena poin privasi dan biaya komputasi selalu memikat pimpinan IT dan keuangan. Lalu alihkan slide ke peta skema 9 tabel Olist.
\end{tipbox}

% ------------------------------------------------------------------------------
\section{Babak 2 --- Mengurai 9 Tabel Olist, Ledakan Fan-Out \& Kebersihan Data}

\begin{stagebox}[Panduan Layar \& Durasi --- Babak 2]
\textbf{Durasi Rencana:} 2,5 menit. \quad \textbf{Visual Layar:} Skema Relasional 9 Tabel Olist + Ilustrasi Bahaya \emph{Join Fan-Out} (1 Order $\times$ 3 Item $\times$ 2 Pembayaran = 6 Duplikat Palsu) + Corong Reduksi \code{99.441 -> 95.824}.
\end{stagebox}

\begin{scriptbox}[Naskah Bicara Lengkap --- Babak 2]
``Untuk membuktikan rekayasa pada kondisi dunia nyata yang berantakan---bukan dataset mainan yang sudah rapi dari buku teks---CSIP memulai dari ekosistem transaksi aktual \textbf{Olist Brasil periode 2016 hingga 2018} yang tersebar di \textbf{sembilan tabel terpisah}: mulai dari tabel pesanan utama, daftar pelanggan, item keranjang belanja, rincian pembayaran cicilan atau voucher, ulasan bintang dan komentar berbahasa Portugis, katalog kategori produk, profil penjual, hingga peta koordinat geolokasi kode pos di seluruh Brasil.

Bagi rekan non-teknis, tantangan terbesar menggabungkan sembilan buku catatan terpisah terletak pada satu jebakan klasik yang kami sebut \textbf{ledakan duplikasi atau Cartesian Fan-Out}. Misalkan ada satu nomor pesanan yang berisi tiga buah barang dan dibayar memakai dua metode---sebagian pakai voucher, sisanya cicilan kartu kredit. Jika seorang analis menggabungkan tabel secara mentah tanpa memadatkan butirannya (\emph{grain}) lebih dulu ke tingkat satu-baris-per-order, komputer akan menggandakan pesanan tunggal tersebut menjadi tiga kali dua alias \textbf{enam baris kembar}! Akibatnya nanti, saat model dilatih dan diuji, satu pesanan yang sama bisa muncul separuh di buku latihan dan separuh lagi di buku ujian---menciptakan ilusi akurasi palsu yang menghancurkan validitas proyek.

Oleh karena itu, di lapisan pembersihan CSIP (\code{cleaner.py} dan kontrak validasi 4 lapis \code{contracts.py}), kami menegakkan satu aturan besi: setiap tabel turunan wajib diagregasi ke butiran tunggal \code{1 baris = 1 order_id unik} sebelum digabungkan. Dari total awal \textbf{99.441 pesanan mentah}, sistem melakukan penyaringan secara transparan dan tercatat: pesanan dengan status dibatalkan (\textbf{625 baris}), pesanan berstatus tidak tersedia yang ternyata 99 persennya memang tidak memiliki catatan item produk (\textbf{609 baris}), dan transaksi tanpa ulasan (\textbf{646 baris}) disisihkan dengan jejak audit yang jelas. Tidak hanya disisihkan, untuk 646 baris tanpa ulasan tersebut, tim kami juga menghitung batas sensitivitas ekstrem bernama \textbf{Manski Bounds}: bahkan seandainya seluruh pembeli yang tidak mengisi ulasan itu sebenarnya marah atau sebenarnya puas, tingkat keluhan populasi hanya bergerak dalam koridor sempit \textbf{11,44 persen hingga 12,10 persen}---selisih kurang dari 0,59 poin persentase. Pada akhirnya, kita memperoleh fondasi bersih berjumlah \textbf{95.824 baris pesanan tersampaikan yang utuh} tanpa satu pun duplikat ID transaksi.''
\end{scriptbox}

\begin{tipbox}[Kontrol Intonasi \& Transisi]
Gunakan gerakan tangan saat menjelaskan contoh $3 \times 2 = 6$ baris kembar agar hadirin non-teknis langsung mengangguk paham mengapa pembersihan data itu menyita kerja keras rekayasa yang serius. Penyebutan angka \textbf{Manski Bounds (11,44\%--12,10\%)} akan langsung memberi sinyal kepada reviewer Data Science senior bahwa Anda memikirkan bias data hilang secara ilmiah.
\end{tipbox}

% ------------------------------------------------------------------------------
\section{Babak 3 --- Benteng Anti-Kebocoran Waktu ($T_0, T_1, T_2$) \& Embargo 90 Hari}

\begin{stagebox}[Panduan Layar \& Durasi --- Babak 3]
\textbf{Durasi Rencana:} 2,5 menit. \quad \textbf{Visual Layar:} Garis Waktu Kronologis Transaksi dengan 3 Pos Pemeriksaan ($T_0$ Checkout, $T_1$ Dispatched, $T_2$ Delivered) + Diagram Partisi Train (`<= 2018-03-31`), Val (`Apr-Jun 2018`), Locked Test (`Jul-Ags 2018`) dengan Pita Embargo 90 Hari.
\end{stagebox}

\begin{scriptbox}[Naskah Bicara Lengkap --- Babak 3]
``Sekarang kita masuk ke bagian yang menurut saya membedakan Sains Data yang jujur dari sekadar pemodelan statistik biasa: \textbf{menghormati garis waktu realitas bisnis}.

Pikirkan hal ini secara logika sederhana. Misalkan ada seorang pembeli baru saja menekan tombol `Bayar' detik ini juga di aplikasi. Apakah sistem kita pada detik itu boleh menebak kemarahannya memakai data \emph{``paket ternyata terlambat 7 hari''}? Tentu tidak! Paketnya saja baru dibayar, bahkan belum masuk kardus di gudang penjual! Jika sebuah model memakai informasi dari masa depan untuk menebak kejadian di masa lalu, model itu tampak super cerdas di presentasi, tetapi langsung lumpuh saat dipasang di toko nyata. Di bidang kami, dosa ini disebut \textbf{kebocoran target atau feature leakage}.

Untuk menutup celah tersebut secara struktural di dalam kode (\code{src/csip/models/trainer.py} dan \code{temporal_cv.py}), CSIP membagi seluruh fitur menjadi tiga pos waktu yang tersegel rapat:
\begin{enumerate}
  \item \textbf{Pos $T_0$ --- Saat Checkout (10 Fitur):} Hanya memuat fakta ekonomi \& jarak yang sudah sah ada detik itu: harga, ongkos kirim, rasio ongkir terhadap harga, panjang janji estimasi kurir, apakah pengiriman lintas negara bagian, serta kategori atau wilayah penjual yang disandikan murni di dalam data latih.
  \item \textbf{Pos $T_1$ --- Saat Penyerahan ke Kurir (11 Fitur):} Menambahkan sinyal jeda keterlambatan penjual mengemas barang ke gudang ekspedisi.
  \item \textbf{Pos $T_2$ --- Pasca-Pengiriman (\code{E3}, 15 Fitur):} Barulah sesudah paket tiba di tangan penerima, model diizinkan melihat berapa hari tunggu aktualnya, seberapa jauh meleset dari janji (\code{delivery_delay_days}), dan fitur interaksi kekecewaan ganda (\code{delay_x_freight_ratio}). Pos $T_2$ bukan untuk meramal sebelum beli, melainkan untuk \textbf{diagnosis pemicu akar dan pemeringkatan prioritas intervensi CS} sebelum ulasan publik meledak.
\end{enumerate}

Lebih jauh lagi, CSIP melarang keras pengacakan data (\emph{random K-fold}) untuk data transaksi berurutan waktu. Kami menerapkan pemecahan kronologis ketat sesuai protokol pra-registrasi terkunci (\code{config/analysis_plan.yaml}): data latih hingga Maret 2018 ($N = 63.666$), data validasi untuk memilih model dan kalibrasi pada April sampai Juni 2018 ($N = 19.587$), dan \textbf{set uji terkunci pasca-Juni 2018} ($N = 12.571$) yang hanya dibuka satu kali di akhir. Untuk validasi silang temporal, empat jendela mengembang kami jaga dengan \textbf{jeda embargo minimal 90 hari} guna memastikan siklus tunggu pengiriman dan pengisian ulasan dari masa latih tidak pernah tumpah ke jendela validasi.''
\end{scriptbox}

\begin{tipbox}[Kontrol Intonasi \& Transisi]
Pastikan kata \textbf{Pos $T_0$, $T_1$, $T_2$} dan \textbf{embargo 90 hari} diucapkan perlahan. Ini adalah tameng terkuat Anda bila di sesi tanya jawab ada yang menanyakan validitas metodologi pemodelan waktu.
\end{tipbox}

% ------------------------------------------------------------------------------
\section{Babak 4 --- Tangga Baseline B0--B3 \& Pertandingan Algoritma yang Adil}

\begin{stagebox}[Panduan Layar \& Durasi --- Babak 4]
\textbf{Durasi Rencana:} 2 menit. \quad \textbf{Visual Layar:} Grafik Tangga Baseline B0 $\rightarrow$ B1 $\rightarrow$ B2 $\rightarrow$ B3 + Tabel Leaderboard Ringkas (\code{ml_leaderboard.csv}).
\end{stagebox}

\begin{scriptbox}[Naskah Bicara Lengkap --- Babak 4]
``Sesudah datanya bersih dan pos waktunya terkunci, pertanyaan berikutnya: algoritma mana yang harus kita pilih? Kami tidak memilih algoritma berdasarkan tren kata kunci, melainkan memakai disiplin \textbf{Tangga Baseline dari B0 sampai B4}.

Karena hanya sekitar 11 hingga 12 persen transaksi Olist berakhir sebagai ulasan negatif, memakai metrik \emph{akurasi persen biasa} akan sangat menipu---sebab menebak semua aman saja sudah benar 88 persen! Oleh sebab itu, papan peringkat kami berpusat pada \textbf{PR-AUC} (\emph{Precision-Recall Area Under Curve}), yang mengukur seberapa tajam model menemukan kelompok keluhan minoritas tanpa mengobral alarm palsu.

Pada Pos pasca-pengiriman $T_2$ di set uji masa depan terkunci ($N = 12.571$):
\begin{itemize}
  \item Model linear klasik \code{B1_logistic} mencetak \textbf{PR-AUC 0,2436} (ROC-AUC 0,6803).
  \item Pohon keputusan tunggal \code{B2_decision_tree} naik ke \textbf{PR-AUC 0,2704}, menyalip algoritma berbasis jarak KNN (0,2532) dan Naive Bayes (0,1996).
  \item Ensembel 100 pohon acak \code{B3_random_forest} melompat ke \textbf{0,3016}.
  \item Akhirnya, \textbf{LightGBM} (\code{B3_lightgbm}) yang membangun 150 pohon histogram secara berurutan di CPU laptop meraih posisi puncak dengan \textbf{PR-AUC 0,3061} dan \textbf{ROC-AUC 0,6745}---mendekati tiga kali lipat tingkat keluhan acak dasar!
\end{itemize}
Sebaliknya, mari kita lihat juga bagaimana performa peringatan dini di \textbf{Pos $T_0$ saat checkout} (tanpa info kurir): skor PR-AUC berada di \textbf{0,1335} dan ROC-AUC \textbf{0,5468}. Angka ini memang mengalahkan tebak acak, tetapi dengan lapang dada membuktikan bahwa tanpa mengetahui keterlambatan aktual kurir, daya prediksi di detik checkout secara alami jauh lebih terbatas dibanding pasca-pengiriman.''
\end{scriptbox}

\begin{tipbox}[Kontrol Intonasi \& Transisi]
Jelaskan alasan mengapa PR-AUC dipakai menggantikan Accuracy dengan kalimat: \emph{``Tebak semua orang sehat di rumah sakit akurasinya tinggi, tapi kita gagal menyelamatkan pasien sakit---itulah gunanya PR-AUC.''} Kalimat ini sangat membekas di telinga non-teknis.
\end{tipbox}

% ------------------------------------------------------------------------------
\section{Babak 5 --- Kejujuran Empiris: Penolakan Repeat Purchase \& Audit Forensik}

\begin{stagebox}[Panduan Layar \& Durasi --- Babak 5]
\textbf{Durasi Rencana:} 2,5 menit. \quad \textbf{Visual Layar:} Kartu Merah Status `DITOLAK' untuk \code{repeat_purchase} (ROC-AUC 0,5021, CI [0,472--0,531]) + Tabel Perbandingan \code{B1_ridge} vs \code{B3_lightgbm} pada Regresi Delay.
\end{stagebox}

\begin{scriptbox}[Naskah Bicara Lengkap --- Babak 5]
``Hadirin sekalian, ada satu bagian paling penting dalam presentasi ini yang justru jarang berani ditampilkan di proyek portofolio: \textbf{laporan tentang apa yang GAGAL dan temuan koreksi audit internal kami}.

Pertama, mengenai prediksi pembelian ulang dalam 90 hari atau \textbf{Repeat Purchase}. Semula kami berharap data riwayat pesanan bisa mendeteksi pembeli mana yang akan berbelanja lagi untuk dikirimi kampanye loyalitas. Setelah membersihkan 962 pelanggan yang sudah membeli ulang sebelum paket pertamanya tiba serta menyisihkan 2.275 pelanggan di batas jendela sensor 90 hari akhir kalender, tingkat belanja ulang riil ternyata sangat kecil ($\sim 2\%$). Saat LightGBM diuji pada set masa depan terkunci, hasilnya mencatatkan \textbf{ROC-AUC = 0,5021} dengan selang kepercayaan 95 persen dari \textbf{0,472 sampai 0,531}, serta PR-AUC hanya \textbf{0,0202}. Angka $0{,}502$ yang memeluk garis $0{,}500$ secara ilmiah tidak lebih baik daripada melempar koin! Artinya, dataset transaksi Olist tidak memuat perilaku kunjungan harian atau minat promo personal untuk meramalkan belanja ulang. Alih-alih memolesnya agar tampak bagus, kami mengetok palu tegas: \textbf{model retensi dinyatakan DITOLAK} dan Decision Engine memblokir otomatis pemicu kupon berdasarkan model ini demi menjaga anggaran perusahaan.

Kedua, tiga rekonsiliasi hasil audit forensik yang kami catat terbuka:
\begin{enumerate}
  \item Pada proyeksi \textbf{Regresi Estimasi Keterlambatan Hari}, di mana tebakan rata-rata naif B0 meleset sebesar MAE 7,45 hari di validasi, ternyata di set uji masa depan model linear sederhana \textbf{B1 Ridge mencetak kesalahan lebih kecil (MAE 4,7509 hari; $R^2$ 0,4454)} mengungguli kompleksitas \textbf{LightGBM (MAE 4,8256 hari; $R^2$ 0,3926)}! Ini membuktikan prinsip ilmiah bahwa model sederhana kerap lebih tangguh saat menghadapi perubahan musim dibandingkan pohon keputusan dalam.
  \item Peringatan dini Pos $T_0$ (E1) telah kami sinkronkan turun dari angka awal eksperimen \code{0,2845} ke angka murni bebas info kurir sebesar \code{0,1335}.
  \item Beda ukuran baris antar-artefak ($N = 12.571$ leaderboard penuh, $12.448$ irisan audit gerbang G6, $38.942$ jendela stabilitas irisan wilayah sejak Maret 2018) terdokumentasi eksplisit lengkap dengan parameter penyaringannya.''
\end{enumerate}
\end{scriptbox}

\begin{tipbox}[Kontrol Intonasi \& Transisi]
Ketika mengucapkan \emph{``model retensi dinyatakan DITOLAK dan kami larang masuk produksi''}, suara Anda harus terdengar justru paling bangga! Penguji berpengalaman tahu bahwa berani menolak model gagal senilai emas di dunia industri.
\end{tipbox}

% ------------------------------------------------------------------------------
\section{Babak 6 --- Uji Hipotesis H1--H6, 3 Segmen Pelanggan \& Simulator Kebijakan S0--S5}

\begin{stagebox}[Panduan Layar \& Durasi --- Babak 6]
\textbf{Durasi Rencana:} 2,5 menit. \quad \textbf{Visual Layar:} Kartu Checklist H1--H6 FDR BH + Donat 3 Segmen (\code{N=96.470}) + Grafik Batang Simulasi Kebijakan S0--S5.
\end{stagebox}

\begin{scriptbox}[Naskah Bicara Lengkap --- Babak 6]
``Setelah mengetahui kekuatan model, kami menerjemahkan temuan menjadi jawaban atas \textbf{enam hipotesis bisnis inti (H1 hingga H6)} memakai koreksi multi-uji \emph{Benjamini-Hochberg FDR} agar tidak ada signifikansi kebetulan:
\begin{itemize}
  \item \textbf{H2 terbukti sangat dominan ($F = 1657{,}035$, $p_{\text{adj}} \approx 0$):} Pelanggaran janji estimasi kurir (\code{delivery_delay_days}) adalah penghancur skor kepuasan terkuat di seluruh ekosistem.
  \item \textbf{H1 ($t = -11{,}858$), H4 lintas wilayah ($F = 29{,}381$), H5 keunggulan pasca-kirim atas checkout ($\Delta\text{PR-AUC} = +0{,}1401$), serta H6 keterkaitan cicilan ($t = -4{,}729$) semuanya terkonfirmasi secara statistik.}
  \item \textbf{Sebaliknya, H3 kembali DITOLAK ($z = -0{,}175$, $p_{\text{adj}} = 0{,}8609$):} Kepuasan skor bintang 5 di transaksi pertama tidak terbukti otomatis mendorong pembelian kedua dalam 90 hari pada barang tahan lama Olist. Mengelola ekspektasi ini menghemat jutaan rupiah kesalahan strategi pemasaran!
\end{itemize}

Selanjutnya, lewat algoritma \textbf{K-Means Clustering} pada \textbf{96.470 transaksi}, evaluasi objektif Silhouette (\code{0,3865}), Davies-Bouldin (\code{1,2168}), dan Calinski-Harabasz (\code{3124,5}) menetapkan \textbf{3 Segmen Operasional} sebagai konfigurasi final (menggantikan percobaan awal $K=4$ di catatan historis lama): \textbf{Segmen 1 (78,0\% atau 75.244 pesanan)} sebagai arus normal berisiko rendah; \textbf{Segmen 0 (19,68\% atau 18.987 pesanan)} sebagai zona berisiko gesekan dengan tingkat keluhan melonjak ke \textbf{15,33 persen}; serta \textbf{Segmen 2 (2,32\% atau 2.239 pesanan)} sebagai ceruk keranjang khusus bernilai tinggi.

Lalu pertanyaannya untuk meja direksi: \emph{kebijakan makro mana yang paling hemat biaya?} Melalui model seri waktu \textbf{VAR Ridge 87 minggu} (dilatih 65 minggu dan diuji keluar-sampel 22 minggu dengan error pemulihan sintetik 1,6 persen), simulator kebijakan memperbandingkan enam skenario. Dari baseline \textbf{S0 dengan keluhan 9,96 persen}, opsi pemborosan \textbf{S1 Subsidi Ongkir Universal 15\%} memang menekan keluhan ke \textbf{7,46\%} tetapi menghambur-hamburkan subsidi untuk seluruh pasar. Sementara itu, opsi cerdas \textbf{S3 Subsidi Terarah Risiko Gesekan} mencapai hampir hasil yang sama yaitu \textbf{7,87\%} hanya dengan menyasar segmen berongkir berat dan rute rawan! Jika dipadukan dengan perbaikan SLA kurir pada skenario \textbf{S5 Blitz Gabungan}, tingkat keluhan diproyeksikan turun hingga \textbf{6,08\%}---sementara di ujung ekstrem lain, skenario stres \textbf{S4 Guncangan Makro} memperingatkan bahaya lonjakan keluhan ke \textbf{14,95\%} jika terjadi gangguan ekspedisi nasional.''
\end{scriptbox}

\begin{tipbox}[Kontrol Intonasi \& Transisi]
Bandingkan secara kontras angka \textbf{S1 (7,46\% boros seluruh toko)} vs \textbf{S3 (7,87\% tepat sasaran hemat anggaran)} vs \textbf{S5 (6,08\% kombinasi optimal)}. Pimpinan bisnis langsung menangkap nilai ROI dari Data Science melalui perbandingan ini!
\end{tipbox}

% ------------------------------------------------------------------------------
\section{Babak 7 --- Kalibrasi Platt Scaling, 15 Faktor SHAP \& Ketahanan Wilayah}

\begin{stagebox}[Panduan Layar \& Durasi --- Babak 7]
\textbf{Durasi Rencana:} 2 menit. \quad \textbf{Visual Layar:} Grafik Kurva Kalibrasi (ECE `0,235 -> 0,007`) + Bar Chart 15 Fitur TreeSHAP + Peta Wilayah \code{SP (4,25 hari)} vs \code{RJ (7,72 hari)}.
\end{stagebox}

\begin{scriptbox}[Naskah Bicara Lengkap --- Babak 7]
``Supaya direksi maupun manajer CS berani bertaruh dengan anggaran nyata di atas angka model, dua syarat wajib harus terpenuhi: \textbf{angkanya harus terkalibrasi jujur} dan \textbf{alasannya harus dapat dibaca manusia}.

Secara bawaan, karena kita melatih LightGBM dengan penalti pembobotan kelas seimbang untuk mengejar kelas keluhan minoritas, skor probabilitas mentahnya terlalu tinggi (\emph{over-confident}), menyisakan kesalahan kalibrasi \emph{Expected Calibration Error (ECE)} sebesar \textbf{0,2355} alias meleset rata-rata 23,5 poin persentase! Dengan melatih \textbf{Platt Scaling} murni di atas Set Validasi April--Juni 2018 (\code{coef = 0,9020}, \code{intercept = -1,8089}), kami menekan kesalahan kalibrasi ECE turun 97 persen menjadi \textbf{0,007015} (hanya meleset 0,7 poin persentase!) dan memperhalus \emph{Brier Score} dari \textbf{0,1375} menjadi \textbf{0,0762}. Kini jika sistem berbunyi \emph{``risiko 30 persen''}, di dunia nyata memang berkisar 30 dari 100 paket serupa yang mengeluh.

Dari sisi transparansi via \textbf{TreeSHAP} (\code{N=800}), 15 pendorong teratas membuktikan akal sehat logika bisnis: peringkat 1--4 didominasi oleh \code{delivery_delay_days}, lama tunggu total, status terlambat, dan interaksi kekecewaan ganda (\code{delay_x_freight_ratio}). Namun kami juga membedah ketahanan per subkelompok (\code{N=38.942} sejak Maret 2018): kami menemukan ketimpangan geografis tajam di mana prediksi hari keterlambatan di pusat kota Sao Paulo sangat akurat dengan kesalahan \textbf{MAE 4,25 hari}, sedangkan di wilayah Rio de Janeiro melebar ke \textbf{7,72 hari}. Demikian pula model sangat jitu pada paket terlambat (\textbf{ROC-AUC 0,8383}), tetapi wajar turun menjadi \textbf{0,6522} pada paket tepat waktu karena alasan komplain di paket tepat waktu umumnya berasal dari cacat barang fisik atau respon toko yang belum tercatat di data logistik.''
\end{scriptbox}

\begin{tipbox}[Kontrol Intonasi \& Transisi]
Pengakuan tentang perbedaan wilayah \code{SP (4,25 hari)} vs \code{RJ (7,72 hari)} dan celah \emph{On-Time} vs \emph{Delayed} membuat audiens merasa Anda sudah mengaudit kelemahan produk Anda sendiri sebelum mereka sempat bertanya.
\end{tipbox}

% ------------------------------------------------------------------------------
\section{Babak 8 --- Jembatan Eksekusi: Decision Engine, Uji A/B 80/20 \& Universal CSV}

\begin{stagebox}[Panduan Layar \& Durasi --- Babak 8]
\textbf{Durasi Rencana:} 2 menit. \quad \textbf{Visual Layar:} Diagram Alur JSON 15 Field \code{DecisionRecommendation} + Desain Eksperimen Uji A/B 80/20 (L2 $\rightarrow$ L3) + Diagram Ekspansi 11 Sektor Universal CSV Studio.
\end{stagebox}

\begin{scriptbox}[Naskah Bicara Lengkap --- Babak 8]
``Bagian terakhir sebelum kita masuk ke \emph{Live Demo} adalah menjawab: \textbf{lalu setelah skor terkalibrasi keluar, aksi nyata apa yang terjadi dan bagaimana CSIP bekerja di luar data Olist?}

Pertama, modul \textbf{Business Decision Engine} mengambil probabilitas terkalibrasi, menimbang batas kuota harian staf CS, lalu mengeluarkan objek keputusan \code{DecisionRecommendation} tervalidasi skema Pydantic dengan \textbf{15 field kerja lengkap}---mulai dari target entitas, tingkat prioritas (\code{CRITICAL/HIGH/MEDIUM/LOW}), faktor pemicu SHAP, tindakan spesifik staf, departemen penanggung jawab, hingga field wajib bernama \code{limitations_and_uncertainty} dan \code{evidence_type}. Sesuai hierarki klaim kami, proyeksi historis berstatus \textbf{Level L2}; untuk melangkah secara sah menuju bukti kausal lapangan \textbf{Level L3}, CSIP membekali rencana eksekusi \textbf{Uji A/B Holdout 80/20}: 80 persen pesanan berisiko mendapat penanganan proaktif, sementara 20 persen ditahan acak sebagai pembanding bersih selama satu siklus operasi penuh.

Kedua, menjawab kebutuhan dunia usaha yang tidak pernah terbatas pada satu dataset e-commerce semata, CSIP v1.1.0 memisahkan studi kasus Olist dengan mesin \textbf{Universal CSV Studio} (\code{ml_engine.py} dan \code{semantic_profiler.py}). Apa pun file CSV atau Excel Anda---entah data sensor mesin pabrik, data gaji HR, progresi pasien medis, hingga rapor mutu produk kimia di \textbf{11 sektor industri}---CSIP otomatis memprofilkan arti kolomnya, memecah latih-uji sebelum transformasi (\emph{Split-Before-Transform}), menyaring kolom bocor target berkorelasi \textbf{$> 0{,}992$}, menyertakan pita ketidakpastian 90 persen \code{[p5, p95]}, serta memvisualisasikan simulasi tuasnya secara real time.

Sekarang, mari kita buktikan semuanya langsung di layar aplikasi \emph{CSIP Studio}, API FastAPI, dan penguji terminal tanpa ada yang ditutupi.''
\end{scriptbox}

\begin{tipbox}[Kontrol Intonasi \& Transisi]
Tutup kalimat terakhir sambil beralih penuh dari slide presentasi ke jendela peramban yang sudah membuka kokpit \emph{CSIP Studio}.
\end{tipbox}


\chapter{Skrip Live Demo Aplikasi, API \& Terminal Langkah-demi-Langkah (Babak 9 s.d. Babak 14)}
\chaptermark{Skrip Live Demo \& Terminal (Babak 9--14)}

Bab ini memandu setiap gerakan kursor dan ketukan keyboard selama demonstrasi langsung. Bacalah naskah bicara sambil mengarahkan mata audiens ke komponen layar yang ditunjuk.

% ------------------------------------------------------------------------------
\section{Babak 9 --- Membuka \emph{CSIP Studio}, Kasus 1 Olist \& Langkah 1 Target}

\begin{stagebox}[Aksi Layar \& Durasi --- Babak 9]
\textbf{Durasi Rencana:} 2 menit. \quad \textbf{Aksi Layar:} Buka tab \emph{CSIP Studio} (\code{http://localhost:8501}). Tunjukkan struktur satu layar bersih (tanpa belasan tab rumit). Pastikan preset terpilih adalah \emph{Kasus 1: Toko Online \& Kepuasan Pembeli (Data Olist)}. Sorot panel pilihan \textbf{Yang ingin dioptimalkan} dan kotak penjelasan \textbf{Artinya begini}.
\end{stagebox}

\begin{scriptbox}[Naskah Bicara Live Demo --- Babak 9]
``Sekarang mari kita lihat wajah platform CSIP yang berhadapan langsung dengan pengguna akhir. Alih-alih memaksa eksekutif atau manajer operasional tersesat di dalam sepuluh tab teknis, \textbf{CSIP Studio} di Streamlit dirancang di atas satu layar terpadu yang jernih.

Di sisi sumber data, kita memiliki pengunggah berkas CSV atau Excel bebas serta tiga kasus bawaan untuk pengujian cepat: \emph{Kasus 1: Toko Online \& Kepuasan Pembeli (Data Olist)}, \emph{Kasus 2: Risiko Kesehatan \& Gaya Hidup (Data Medis)}, dan \emph{Kasus 3: Kualitas Produk \& Komposisi Bahan (Data Mutu)}.

Mari kita aktifkan dulu \textbf{Kasus 1: Data Olist}. Di Langkah 1, pada menu \textbf{`Yang ingin dioptimalkan'}, pengguna memilih sasaran bisnisnya---misalnya kepuasan pembeli. Begitu dipilih, perhatikan bagian \textbf{`Artinya begini'} tepat di bawahnya: mesin profil semantik kami otomatis menerjemahkan kolom teknis dataset menjadi kalimat manusia yang langsung dipahami siapa pun di ruangan ini, tanpa perlu menghafal nama kolom database.''
\end{scriptbox}

\begin{tipbox}[Tips Kontrol Layar]
Jangan terburu-buru menggeser halaman ke bawah. Biarkan audiens melihat 3 detik bahwa seluruh kendali berada di tempat yang intuitif dan responsif.
\end{tipbox}

% ------------------------------------------------------------------------------
\section{Babak 10 --- Menggerakkan Tuas Simulasi, Kartu Bento \& Narasi \emph{Nara}}

\begin{stagebox}[Aksi Layar \& Durasi --- Babak 10]
\textbf{Durasi Rencana:} 2,5 menit. \quad \textbf{Aksi Layar:} Arahkan ke panel \textbf{Tuas simulasi} (3 slider faktor paling berpengaruh). Geser penggeser keterlambatan/ongkir ke arah lebih baik, tunjuk perubahan angka di ketiga kartu KPI Bento, baca kotak narasi \textbf{Nara}, lalu klik tombol \textbf{Setel ulang} untuk membuktikan pengembalian ke baseline.
\end{stagebox}

\begin{scriptbox}[Naskah Bicara Live Demo --- Babak 10]
``Di sebelah sini kita melihat \textbf{Panel Tuas Simulasi}. Dari puluhan variabel dalam data, algoritma pemeringkat secara otomatis hanya memunculkan \textbf{tiga tuas pengungkit paling kuat} agar pengambil keputusan fokus pada prioritas nyata, bukan tergoda mengutak-atik variabel sepele.

Lihat tiga kartu KPI Bento di bagian atas layar: saat ini angka menampilkan posisi dasar (\emph{baseline}) beserta indikasi rentang keandalan model. Sekarang bayangkan seorang manajer logistik bertanya di rapat: \emph{`Bagaimana kira-kira nasib metrik target kita kalau rata-rata keterlambatan kirim atau beban rasio ongkir berhasil ditekan dari titik ini?'} Mari saya geser tuasnya perlahan ke kiri...

Perhatikan dua hal yang langsung bereaksi secara waktu nyata: pertama, ketiga kartu KPI Bento langsung membandingkan angka baseline terhadap hasil skenario baru beserta selisih perubahannya; kedua, mari kita baca panel \textbf{`Nara'}---mesin narasi otomatis kami. Nara tidak sekadar melempar angka mentah, melainkan menyusun paragraf ringkasan eksekutif berbahasa Indonesia lengkap dengan pengingat jujur bahwa ini adalah proyeksi sensitivitas berbasis korelasi model terkalibrasi (\code{HYPOTHETICAL_SENSITIVITY}), bukan klaim kausal mutlak. Dan bila kita ingin kembali menguji skenario lain dari kondisi awal, cukup klik satu tombol \textbf{`Setel ulang'}.''
\end{scriptbox}

\begin{tipbox}[Tips Kontrol Layar]
Gerakkan slider satu per satu secara perlahan agar audiens menyaksikan angka Bento dan teks Nara ikut berubah. Klik \textbf{Setel ulang} sebelum pindah ke lensa visualisasi!
\end{tipbox}

% ------------------------------------------------------------------------------
\section{Babak 11 --- Menjelajah Keempat \emph{Lensa Tampilan} Analitik}

\begin{stagebox}[Aksi Layar \& Durasi --- Babak 11]
\textbf{Durasi Rencana:} 2,5 menit. \quad \textbf{Aksi Layar:} Klik satu demi satu keempat pilihan radio pada \textbf{Lensa tampilan}: (1) \code{1. Kurva Simulasi dan Titik Kritis}, (2) \code{2. Peringkat Faktor Penentu}, (3) \code{3. Sebaran Kelompok Data}, (4) \code{4. Observatorium Drift & Stabilitas}.
\end{stagebox}

\begin{scriptbox}[Naskah Bicara Live Demo --- Babak 11]
``Supaya pengambil keputusan bisnis maupun Data Scientist bisa membedah pola tanpa berpindah-pindah aplikasi, kokpit menyediakan empat sudut pandang radio di bawah menu \textbf{Lensa tampilan}:

\begin{enumerate}
  \item \textbf{Lensa 1 --- Kurva Simulasi dan Titik Kritis:} Saat diklik, grafik ini memperlihatkan kurva respons 70 titik dengan penyebaran Gaussian ($\sigma = 2{,}0$). Kita dapat melihat dengan mata telanjang di mana letak \emph{titik kritis atau tekukan}---yakni momen di mana penurunan risiko mulai melandai sehingga menyuntikkan subsidi lebih lanjut tidak lagi memberi hasil sepadan.
  \item \textbf{Lensa 2 --- Peringkat Faktor Penentu:} Lensa ini membuktikan transparansi model. Ia menampilkan urutan kontribusi pendorong faktor secara jujur sehingga staf bisnis mengerti \emph{mengapa} sebuah tuas diposisikan sebagai prioritas utama dan mana faktor yang kontribusinya kecil.
  \item \textbf{Lensa 3 --- Sebaran Kelompok Data:} Di sini kita melihat peta distribusi aktual dari kelompok transaksi atau segmen data. Ini menjaga pengambil keputusan agar tidak membuat kebijakan hanya dari nilai rata-rata semu sementara sebaran populasinya berkelompok tajam.
  \item \textbf{Lensa 4 --- Observatorium Drift \& Stabilitas:} Terakhir, ini adalah ruang kendali keandalan Data Science. Di sini dipantau stabilitas performa antarlipatan validasi serta pemantauan pergeseran distribusi data (\emph{drift})---mengingatkan tim kapan model masih layak dipercaya dan kapan waktunya dilatih ulang.''
\end{enumerate}
\end{scriptbox}

% ------------------------------------------------------------------------------
\section{Babak 12 --- Demonstrasi Multi-Domain: Beralih Preset \& Unggah CSV Baru}

\begin{stagebox}[Aksi Layar \& Durasi --- Babak 12]
\textbf{Durasi Rencana:} 2,5 menit. \quad \textbf{Aksi Layar:} Pilih \emph{Kasus 2: Risiko Kesehatan \& Gaya Hidup (Data Medis)} atau unggah file \code{.csv} uji coba eksternal. Tunjukkan bagaimana \emph{Yang ingin dioptimalkan}, \emph{Artinya begini}, 3 tuas baru, dan kartu KPI beradaptasi otomatis tanpa satu baris pun kode diubah.
\end{stagebox}

\begin{scriptbox}[Naskah Bicara Live Demo --- Babak 12]
``Sekarang mari kita jawab pertanyaan terpenting bagi kegunaan dunia nyata: \emph{Bagaimana jika besok divisi HR, manufaktur, keuangan, atau rumah sakit datang membawa berkas CSV mereka sendiri yang sama sekali tidak ada hubungannya dengan Olist Brasil?}

Perhatikan saat saya beralih ke \textbf{Kasus 2: Risiko Kesehatan \& Gaya Hidup (Data Medis)}---atau saat kita mengunggah file CSV acak melalui kotak unggahan di atas. Dalam hitungan detik, mesin \textbf{Universal CSV Studio} bekerja penuh di CPU lokal laptop ini:
\begin{enumerate}
  \item Modul \code{semantic_profiler.py} mencocokkan nama kolom dengan kamus \textbf{11 sektor industri} sekaligus menjalankan inferensi zero-shot untuk kolom asing; lihat bagian \emph{Artinya begini} langsung berubah mulus berbicara tentang indikator medis dan satuannya (\code{mg/dL}, \code{mmHg}, atau \code{poin})!
  \item Modul \code{ml_engine.py} membatasi sampel aman di maksimal \textbf{3.500 baris} agar RAM 8 GB tetap ringan, melakukan \emph{Split-Before-Transform} 75/25 sebelum imputasi median, serta menyalakan \textbf{radar anti-kebocoran}: bila ada kolom berkolerasi $> 0{,}992$ dengan target atau $> 0{,}94$ antarsesama fitur, sistem otomatis membersihkannya di sisi data latih sebelum melakukan validasi silang K-Fold dan menghitung pita interval prediksi 90 persen \code{[p5, p95]}.
\end{enumerate}
Jadi satu arsitektur CSIP melayani dua dunia sekaligus: ketelitian riset historis mendalam pada Olist, sekaligus kelincahan analitik serbaguna untuk tabel CSV/Excel bisnis apa pun.''
\end{scriptbox}

% ------------------------------------------------------------------------------
\section{Babak 13 --- Verifikasi Produksi: FastAPI Swagger \& Gerbang Terminal}

\begin{stagebox}[Aksi Layar \& Durasi --- Babak 13]
\textbf{Durasi Rencana:} 2 menit. \quad \textbf{Aksi Layar:} Pindah ke tab Swagger UI \code{http://127.0.0.1:8000/docs}. Tunjukkan ke-7 rute endpoint terutama \code{/ready}, \code{/predict/batch}, \code{/decisions/recommend}, dan \code{/monitor/drift}. Kemudian perlihatkan terminal menjalankan \code{python tools/verify_release.py} yang mencetak 12 gerbang \code{PASS} + output \code{pytest} (\code{216 passed}).
\end{stagebox}

\begin{scriptbox}[Naskah Bicara Live Demo --- Babak 13]
``Untuk audiens rekayasa perangkat lunak dan MLOps, CSIP Studio hanyalah ujung depan yang dapat disaksikan manusia. Di belakangnya berdiri mesin siap-integrasi yang kita lihat di tab Swagger UI FastAPI ini.

Di sini tersedia tujuh rute terstruktur: \code{/health} dan \code{/ready} untuk pengawasan kesiapan peladen; \code{/predict/batch} untuk menembak ratusan pesanan sekaligus; \code{/decisions/recommend} yang langsung mengeluarkan struktur JSON \code{DecisionRecommendation} berisi 15 field kerja siap santap sistem tiket perusahaan; \code{/monitor/drift} untuk memeriksa indeks PSI, KS, Wasserstein, dan JSD secara terprogram; serta \code{/models/registry} untuk jejak versi dan hash rollback.

Dan bagaimana cara membuktikan semua yang saya ucapkan sejak menit pertama bukan sekadar janji presentasi? Mari tengok jendela terminal ini. Ketika saya menjalankan \code{python tools/verify_release.py}, sistem memeriksa otomatis kedua belas gerbang rilis dari G1 sampai G12---mulai dari keseragaman versi, hasil \textbf{216 kasus uji \code{pytest} yang lulus 100\%} tanpa kegagalan, pemindaian frasa anti-sesumbar, keutuhan hash SHA-256 protokol pra-registrasi, verifikasi embargo 90 hari, hingga kecocokan hash fisik 63 KUK SKKNI. Semuanya menyala hijau: \textbf{12 dari 12 gerbang berstatus PASS}.''
\end{scriptbox}

% ------------------------------------------------------------------------------
\section{Babak 14 --- Rencana Darurat (\emph{Contingency Plan}) \& Pidato Penutup Panggung}

\subsection{Protokol Jika Demo Macet di Hadapan Penguji}
\begin{itemize}
  \item \textbf{Jika Streamlit Port Sibuk/Layar Lambat:} Jangan panik atau mengetik berulang-ulang di depan proyektor. Segera beralih ke tab dokumentasi artefak lokal (\code{artifacts/reports/ml_leaderboard.csv}, \code{reports/verification/gates.json}, dan halaman \code{docs/BUKU_INDUK_CSIP_LENGKAP.pdf}) sambil berkata: \emph{``Sambil layanan lokal memuat ulang cache, seluruh angka tervalidasi dan laporan pengujinya sudah terkunci deterministik di berkas artefak resmi kita di sini...''}
  \item \textbf{Jika CSV Eksternal User Memiliki Format Aneh/Rusak:} Tunjukkan betapa bergunanya pesan pengaman atau fallback jujur sistem (\code{IN_SAMPLE_FALLBACK_SMALL_N} bila baris $< 8$, atau pembersihan NaN median) sebagai bukti bahwa aplikasi gagal secara terkontrol (\emph{graceful degradation}) tanpa memuntahkan error mentah berbahaya.
\end{itemize}

\begin{scriptbox}[Naskah Bicara Penutup Master (60 Detik Terakhir)]
``Sebagai penutup, izinkan saya merangkum filosofi utama proyek \textbf{CSIP v1.1.0} dalam tiga kalimat.

\textbf{Pertama}, kami memilih \textbf{kejujuran ilmiah di atas angka kosmetik}: mengunci pos waktu $T_0/T_1/T_2$, menghormati embargo 90 hari, menurunkan error kalibrasi hingga 0,7 poin persentase, berani menolak model retensi yang tidak lebih baik dari tebak acak, dan jujur saat model linear sederhana Ridge ternyata mengungguli pohon kompleks di set uji masa depan.

\textbf{Kedua}, kami menjembatani jarak antara Sains Data dan meja keputusan lewat \textbf{Decision Engine 15 field} serta komparasi kebijakan makro terarah (seperti skenario S3 dan S5) yang melindungi anggaran pemasaran perusahaan dari subsidi membabi buta.

\textbf{Ketiga}, seluruh platform berjalan ringan di laptop RAM 8 GB, menjaga privasi 100 persen lokal tanpa ketergantungan API awan, teruji oleh 216 tes otomatis serta 12 gerbang rilis, dan mampu langsung membantu keputusan lintas 11 sektor industri melalui Universal CSV Studio.

Terima kasih banyak atas perhatian Bapak dan Ibu. Sekarang saya dengan senang hati membuka sesi tanya jawab dan bedah teknis.''
\end{scriptbox}

% ==============================================================================
% BAB 16: ENSIKLOPEDIA Q&A 25 PERTANYAAN
% ==============================================================================
\chapter{Ensiklopedia 25 Pertanyaan \& Jawaban Penguji (\emph{Q\&A Defense Playbook})}
\chaptermark{Ensiklopedia Q\&A Penguji (25 Kasus)}

Bab penutup ini mempersiapkan Anda menghadapi hampir segala jenis pertanyaan tajam dari direksi non-teknis, Lead Data Scientist, rekayasawan MLOps, maupun asesor kompetensi SKKNI. Setiap pertanyaan memuat dua lapis jawaban: \textbf{Jawaban Singkat Eksekutif (30 detik)} dan \textbf{Jawaban Mendalam Data Science \& Rekayasa}.

\section{Kelompok A --- Strategi Data, Kecukupan Dataset \& Universal CSV Studio (Q1--Q6)}

\subsection*{Q1. Apakah 1 dataset Olist ($\sim$100 ribu pesanan) cukup bagus atau terlalu kecil untuk proyek Data Science serius?}
\begin{itemize}
  \item \textbf{Jawaban Singkat Eksekutif:} Cukup memadai secara volume sekaligus realistis untuk studi kasus e-commerce terstruktur---dengan syarat kita mengolah ke-9 tabel relasionalnya secara benar dan memisahkan batasan generalisasinya. Kekurangan keragaman industri sudah kami selesaikan lewat arsitektur \emph{Universal CSV Studio}.
  \item \textbf{Jawaban Mendalam DS:} Dengan \code{99.441} pesanan dipecah di 9 tabel (\code{order_items}, \code{payments}, \code{reviews}, \code{geolocation}, dll.) selama hampir dua tahun kalender Brasil (\code{2016-09} s.d. \code{2018-10}), Olist tergolong dataset tabular industri menengah-kuat untuk uji coba pencegahan \emph{fan-out}, kalibrasi ketimpangan kelas 11\%, dan validasi silang temporal dengan embargo 90 hari ($N=12.571$ locked test). Namun, kelemahan intrinsik Olist memang terungkap oleh sains kita sendiri: minimnya transaksi berulang per \code{customer_unique_id} menyebabkan model retensi mentok di ROC-AUC \code{0,5021}. Itu sebabnya CSIP tidak membanggakan diri hanya pada Olist, melainkan membangun \code{ml_engine.py} + \code{semantic_profiler.py} agar mampu memodelkan dataset eksternal lintas 11 sektor.
\end{itemize}

\subsection*{Q2. Kalau penguji atau klien memasukkan file CSV dari bisnis lain di luar e-commerce, apa strategi CSIP agar tetap akurat dan aman?}
\begin{itemize}
  \item \textbf{Jawaban Singkat Eksekutif:} CSIP tidak mewarisi model Olist secara membabi buta. Saat file CSV baru masuk, sistem langsung mengenali bahasa kolomnya lintas 11 sektor, menyaring kolom yang berpotensi curang/bocor, melatih model khusus di atas file tersebut secara lokal, lalu menampilkan simulasinya dengan pita batas ketidakpastian.
  \item \textbf{Jawaban Mendalam DS:} Strateginya terpusat pada jalur adaptif \code{src/csip/models/ml_engine.py}: (1) pembatasan beban memori RAM 8 GB pada \code{sample_cap = 3500}; (2) deteksi otomatis apakah target bertipe biner $\{0,1\}$ (klasifikasi) atau kontinu (regresi); (3) penerapan \emph{Split-Before-Transform} (\code{test_size = 0,25}) sebelum imputasi median di \code{fit} pada train saja; (4) filter anti-kebocoran korelasi target $> 0{,}992$ dan multikolinearitas $> 0{,}94$ pada sisi latih; (5) validasi silang K-Fold/StratifiedKFold maksimal 5 lipatan; (6) penyertaan interval 90\% \code{[p5, p95]} dan leksikon semantik \code{semantic_profiler.py} berleksikon 11 industri plus inferensi zero-shot.
\end{itemize}

\subsection*{Q3. Bagaimana jika kolom target di file CSV baru pengguna adalah kolom bocor yang sebenarnya merupakan salinan dari kolom input lain?}
\begin{itemize}
  \item \textbf{Jawaban Singkat Eksekutif:} Radar anti-kebocoran di CSIP akan mencium kejanggalan korelasi yang hampir seratus persen mirip, langsung membuang kolom input yang mencurigakan itu sebelum melatih AI agar hasil simulasi tidak menipu pemilik bisnis.
  \item \textbf{Jawaban Mendalam DS:} Dalam fase prapemrosesan train-only di \code{ml_engine.py}, matriks korelasi antara fitur kandidat dan target dihitung ketat. Setiap fitur yang korelasinya melampaui ambang \textbf{$> 0{,}992$} otomatis dieliminasi dari matriks fitur latih (dan tidak ikut diikutsertakan saat evaluasi tes), sementara fitur yang berkorelasi antarfitur $> 0{,}94$ juga disederhanakan agar koefisien dan penentuan 3 tuas teratas tidak didominasi kolom kembar.
\end{itemize}

\subsection*{Q4. Bagaimana CSIP Studio tetap bekerja bila pengguna mengunggah dataset sangat kecil, misalnya kurang dari 8 baris?}
\begin{itemize}
  \item \textbf{Jawaban Singkat Eksekutif:} Sistem tidak akan memaksa membelah data secara acak dan mengaku-ngaku punya skor ujian valid; sebaliknya ia memunculkan bendera peringatan jujur bahwa data terlalu sedikit untuk dipecah sebagai evaluasi generalisasi.
  \item \textbf{Jawaban Mendalam DS:} Ketika $N < 8$, \code{ml_engine.py} memicu jalur \code{IN_SAMPLE_FALLBACK_SMALL_N}. Flag status ini dikirim eksplisit ke antarmuka agar \emph{Nara} dan komponen UI melaporkan keterbatasan sampel secara transparan alih-alih menipu pemakai dengan skor validasi palsu.
\end{itemize}

\subsection*{Q5. Mengapa tidak mengandalkan LLM Cloud besar (seperti GPT/Gemini/Claude) saja untuk menganalisis CSV yang diunggah?}
\begin{itemize}
  \item \textbf{Jawaban Singkat Eksekutif:} Pertama karena privasi: data gaji karyawan, medis pasien, dan keuangan perusahaan tidak boleh bocor keluar laptop. Kedua karena ketepatan hitung: LLM pandai menulis kalimat, tetapi sering berhalusinasi saat diminta menghitung regresi, K-Fold, dan interval 90\% secara akurat.
  \item \textbf{Jawaban Mendalam DS:} Sesuai prinsip tata kelola, komputasi deterministik (LightGBM, Ridge, kalibrasi, uji statistik, filter korelasi) harus dipisahkan dari lapisan narasi. CSIP melakukan seluruh komputasi statistik di scikit-learn/LightGBM secara lokal pada RAM $\le 8$ GB dan memakai mesin template semantik deterministik (\emph{Nara}) tanpa satu pun pemanggilan eksternal---menjamin nol kebocoran data sensitif dan hasil deterministik.
\end{itemize}

\subsection*{Q6. Apakah grafik simulasi tuas di CSIP Studio (\emph{Kurva Simulasi}) membuktikan hubungan sebab-akibat?}
\begin{itemize}
  \item \textbf{Jawaban Singkat Eksekutif:} Tidak! Grafik itu adalah alat baca sensitivitas ``seandainya nilai fitur bergerak berdasar pola riwayat data lama'', dan aplikasinya selalu memberi label peringatan terbuka agar manajer tidak salah menganggapnya jaminan kausal mutlak.
  \item \textbf{Jawaban Mendalam DS:} Kurva 70 titik pada \code{ml_engine.py} (sebaran Gaussian $\sigma = 2{,}0$) ditandai dengan penanda eksplisit \code{HYPOTHETICAL_SENSITIVITY}. Perubahan satu tuas sembari menahan fitur lain mencerminkan respons kondisional/asosiatif model terlatih (\emph{partial-dependence style}) di dalam wilayah klaim Level L2, bukan hasil manipulasi eksperimental teracak.
\end{itemize}

\section{Kelompok B --- Isolasi Kebocoran Waktu, Metodologi \& Algoritma ML (Q7--Q16)}

\subsection*{Q7. Jelaskan paling sederhana apa bahayanya \emph{Join Fan-Out} saat menggabungkan 9 tabel Olist!}
\begin{itemize}
  \item \textbf{Jawaban Singkat Eksekutif:} Satu pesanan berisi 3 barang dan 2 cara bayar kalau digabung langsung tanpa dipadatkan akan meledak jadi 6 baris kembar. Akibatnya komputer menghitung 1 keluhan seolah terjadi 6 kali, dan satu transaksi sama masuk bersamaan ke data belajar dan data ujian!
  \item \textbf{Jawaban Mendalam DS:} Penggabungan langsung relasi $1\text{-ke-}N$ (\code{orders} $\rightarrow$ \code{order_items} $\rightarrow$ \code{order_payments}) menciptakan perkalian Kartesius. Sebelum penggabungan di \code{cleaner.py}, setiap tabel anak wajib diagregasi di atas kunci primer \code{order_id} (total nilai item, jumlah cicilan maksimum, rata-rata skor ulasan bila ganda) sehingga integritas \emph{grain} \code{1 row = 1 order_id} terverifikasi oleh uji otomatis di \code{test_cleaning_and_grain.py}.
\end{itemize}

\subsection*{Q8. Apa beda Pos $T_0$, $T_1$, dan $T_2$? Mengapa model terbaik justru memakai $T_2$ yang datang setelah pengiriman?}
\begin{itemize}
  \item \textbf{Jawaban Singkat Eksekutif:} $T_0$ menilai risiko di detik pelanggan membayar, sedangkan $T_2$ menilai paket persis setelah tiba di rumah pembeli untuk menyortir siapa yang harus segera dihubungi tim CS sebelum mereka menulis komentar bintang satu di aplikasi.
  \item \textbf{Jawaban Mendalam DS:} Memakai fitur $T_2$ seperti \code{delivery_delay_days} untuk mengklaim prediksi \emph{sebelum pengiriman} adalah kebocoran waktu fatal. CSIP jujur melaporkan kedua pos secara terpisah: Pos $T_0$ menghasilkan PR-AUC terbatas \code{0,1335} (ROC \code{0,5468}), membuktikan kelayakan prediksi dini di detik checkout masih lemah. Sebaliknya, Pos $T_2$ (\code{E3}, 15 fitur) mencetak PR-AUC \code{0,3061} yang sah dipakai untuk \textbf{triase CS pasca-pengiriman sebelum survei ulasan ditutup/terisi} maupun root-cause audit.
\end{itemize}

\subsection*{Q9. Mengapa harus ada jeda \emph{Embargo 90 Hari} pada validasi silang temporal? Bukankah 7 hari sudah cukup?}
\begin{itemize}
  \item \textbf{Jawaban Singkat Eksekutif:} Di Brasil paket antarwilayah memakan waktu berminggu-minggu dan pelanggan baru mengisi ulasan setelah barang tiba. Jika jedanya cuma 7 hari, pesanan di ujung masa latihan masih belum selesai urusannya dan tumpang tindih dengan masa ujian.
  \item \textbf{Jawaban Mendalam DS:} Dalam \code{RollingOriginTemporalCV} (4 lipatan jendela mengembang), embargo $\ge 90$ hari melindungi tidak hanya batas atas pengiriman jarak jauh plus tenggat isi ulasan, tetapi juga menyelaraskan jendela kohort pengamatan 90 hari agar tidak ada kebocoran status target ke depan (\emph{label overlap}) antarfold. Untuk uji retensi di tepi kanan kalender \code{2018-10-17}, \code{2.275} pelanggan dalam 90 hari terakhir ikut disaring sebagai \emph{right-censoring buffer}.
\end{itemize}

\subsection*{Q10. Mengapa metrik utama memilih PR-AUC bukan Accuracy atau sekadar ROC-AUC?}
\begin{itemize}
  \item \textbf{Jawaban Singkat Eksekutif:} Karena hampir 89\% pembeli tidak menulis ulasan buruk. Kalau model menebak ``semua aman'', akurasinya otomatis hampir 89\% padahal tidak menangkap satu pun pelanggan yang marah! PR-AUC memaksa model membuktikan ketajamannya tepat di kelompok keluhan yang langka itu.
  \item \textbf{Jawaban Mendalam DS:} Pada kelas tak seimbang ($\sim 11\text{--}12\%$ positif ulasan negatif), True Negatives yang masif membuat False Positive Rate (FPR) tetap kecil sehingga ROC-AUC cenderung tampak optimistis. Sebaliknya, \emph{Precision-Recall AUC} tidak melibatkan True Negatives, sehingga secara langsung mengukur trade-off ketepatan alarm keluhan terhadap daya tangkap penarikan (\emph{recall}) bagi kapasitas tim CS.
\end{itemize}

\subsection*{Q11. Mengapa tidak malu mempublikasikan model \texttt{repeat\_purchase} yang gagal dengan ROC-AUC 0,5021?}
\begin{itemize}
  \item \textbf{Jawaban Singkat Eksekutif:} Justru menyembunyikannya adalah bencana bagi bisnis! Jika kami meloloskan model berketepatan setara lempar koin itu, perusahaan akan membuang anggaran promo berbulan-bulan ke target pelanggan yang acak. Menolak model gagal menghemat uang nyata.
  \item \textbf{Jawaban Mendalam DS:} Pada kohort bersih retensi (setelah mengeluarkan 962 pre-delivery repeaters dan 2.275 right-censored), \code{B3_lightgbm} memperoleh ROC-AUC \code{0,5021} dengan 95\% CI \code{[0,472, 0,531]} dan PR-AUC \code{0,0202}, searah dengan penolakan hipotesis H3 ($p_{\text{adj}} = 0{,}8609$). Ini memvalidasi keterbatasan informasional tabel transaksi Olist dan mengeksekusi penolakan produksi di level aturan \code{rules.py}.
\end{itemize}

\subsection*{Q12. Di papan peringkat regresi delay, \texttt{B1\_ridge} (MAE 4,75 hari) lebih baik daripada \texttt{B3\_lightgbm} (MAE 4,83 hari). Mengapa model linear mengalahkan AI gradien pohon?}
\begin{itemize}
  \item \textbf{Jawaban Singkat Eksekutif:} Hubungan antara jarak lokasi dan janji estimasi pengiriman polanya cukup lurus dan teratur. Model linear sederhana dengan penalti stabil lebih tahan terhadap pergeseran musim dibanding pohon bercabang banyak yang terlalu sensitif pada detail masa lalu.
  \item \textbf{Jawaban Mendalam DS:} Di set uji terkunci ($N=12.630$), \code{B1_ridge} mencatat MAE \code{4,7509}, RMSE \code{5,9685}, dan $R^2 = 0{,}4454$, mengungguli \code{B3_lightgbm} (MAE \code{4,8256}, RMSE \code{6,2458}, $R^2 = 0{,}3926$). Karena komponen varians aditif dari jangka estimasi vs rute antarnegara bagian dominan linear sementara cabang pohon LightGBM bisa mengalami sedikit degradasi out-of-time terhadap pola musiman pertengahan 2018, temuan ini dicatat jujur sebagai bahan evaluasi penggantian champion di rilis mendatang.
\end{itemize}

\subsection*{Q13. Bagaimana persisnya Platt Scaling memperbaiki kalibrasi dari ECE 0,235 menjadi 0,007 tanpa menyentuh set uji?}
\begin{itemize}
  \item \textbf{Jawaban Singkat Eksekutif:} Model latih awal dibuat ``paranoid'' agar peka mendeteksi komplain. Setelah selesai belajar, kami menyetel ulang skala persentasenya memakai cermin terpisah (Set Validasi) supaya bila layar menyebut risiko 20\%, kenyataannya memang sekitar 20 dari 100 pesanan serupa mengeluh.
  \item \textbf{Jawaban Mendalam DS:} Pembobotan \code{class_weight="balanced"} mendistorsi log-odds keluaran mentah menuju kelas minoritas. CSIP fitting regresi logistik 1D (\emph{Platt Scaling}, \code{coef = 0,9020}, \code{intercept = -1,8089}) \textbf{eksklusif pada Set Validasi April--Juni 2018}. Evaluasi menunjukkan ECE terjun dari \code{0,235486} ke \code{0,007015} dan Brier Score dari \code{0,137478} ke \code{0,076177}. Set uji terkunci pasca-Juni tetap murni belum tersentuh hingga penilaian tunggal akhir.
\end{itemize}

\subsection*{Q14. Mengapa terjadi perbedaan angka klaster ($K=3$ vs $K=4$) dan populasi uji ($N=12.571$ vs $12.448$ vs $38.942$) antar dokumen?}
\begin{itemize}
  \item \textbf{Jawaban Singkat Eksekutif:} Dokumen riset mengalami pemutakhiran penyaringan yang makin ketat. Semua perbedaan itu kami audit dan catat eksplisit agar siapapun yang membuka berkas lama maupun terbaru bisa menelusuri asal-usul tiap angka tanpa bingung.
  \item \textbf{Jawaban Mendalam DS:} (1) Pada klastering, catatan eksploratif awal \code{MODEL_RESULTS_REGISTRY.csv} menyimpan $K=4$ (Silhouette \code{0,4180}), sedangkan pipeline deterministik terbaru \code{clustering_evaluation.json} ($N=96.470$) membandingkan $K$ pasca-standardisasi dan memilih \code{selected_k = 3} (Silhouette \code{0,3865} vs $K=4$ \code{0,3485}); (2) Pada evaluasi, $N=12.571$ ialah ukuran locked test leaderboard E3 (dan $12.630$ untuk delay tanpa wajib review), $N=12.448$ pada \code{gates.json} ialah irisan audit pra-rekonsiliasi G6 (\code{PR-AUC 0,3151}), dan $N=38.942$ pada \code{model_slice_analysis.json} mencakup evaluasi jendela ketahanan wilayah lebih lebar sejak Maret 2018 (\code{ROC 0,7517}).
\end{itemize}

\subsection*{Q15. Apa implikasi praktis bagi bisnis bahwa model di Rio de Janeiro (\texttt{RJ}) jauh lebih meleset (MAE 7,72 hari) dibanding Sao Paulo (\texttt{SP}, MAE 4,25 hari)?}
\begin{itemize}
  \item \textbf{Jawaban Singkat Eksekutif:} Artinya kita tidak boleh menyamakan standar janji kurir maupun batas waspada untuk seluruh negara. Paket tujuan Rio de Janeiro perlu penyangga waktu ekstra atau evaluasi mitra kurir lokal sebelum janji pengiriman ditampilkan ke pembeli.
  \item \textbf{Jawaban Mendalam DS:} Dari temuan slice geografis di \code{model_slice_analysis.json} dan terkonfirmasinya hipotesis \textbf{H4} ($F = 29{,}381$), keseragaman error tidak terpenuhi lintas negara bagian (\emph{heteroskedastisitas spasial}). Rekomendasi rekayasanya adalah mengalibrasi \emph{safety buffer} estimasi kedatangan spesifik per negara bagian dan menandai \code{RJ} sebagai zona sensitif dalam aturan Decision Engine.
\end{itemize}

\subsection*{Q16. Bagaimana cara CSIP memitigasi masalah 646 transaksi tanpa ulasan? Tidakkah membuangnya menciptakan bias seleksi?}
\begin{itemize}
  \item \textbf{Jawaban Singkat Eksekutif:} Benar, selalu ada kemungkinan orang yang diam itu sebenarnya marah atau sebenarnya puas. Karena itu kami menghitung simulasi skenario terburuk hingga terbaik: ternyata batas pergeserannya sangat tipis, maksimal kurang dari 0,6 poin persentase.
  \item \textbf{Jawaban Mendalam DS:} Alih-alih berasumsi bahwa data hilang bersikap \emph{Missing Completely at Random (MCAR)}, CSIP menghitung batas non-parametrik \textbf{Manski Bounds}: dengan asumsi terburuk seluruh 646 pesanan tanpa ulasan bernilai negatif maupun asumsi terbaik seluruhnya positif, laju keluhan populasi terkunci di selang \textbf{11,44\% hingga 12,10\%} ($\Delta \le 0{,}59\%$), memenuhi gerbang kualitas \textbf{G7}.
\end{itemize}

\section{Kelompok C --- Decision Engine, Uji A/B, MLOps, Keamanan \& SKKNI (Q17--Q25)}

\subsection*{Q17. Bagaimana Decision Engine menentukan siapa yang dibantu jika CS hanya sanggup menangani 5\% pesanan hari ini?}
\begin{itemize}
  \item \textbf{Jawaban Singkat Eksekutif:} Sistem tidak memakai batas kaku 50\%. Ia mengurutkan pesanan dari tingkat bahaya terkalibrasi paling tinggi dipadukan dengan nilai ekonomis pesanan, lalu memotong pas di batas kuota kerja tim CS hari itu.
  \item \textbf{Jawaban Mendalam DS:} \code{BusinessDecisionEngine} menggabungkan ambang kapasitas dinamis dengan bobot prioritas serta kontrol biaya (\code{estimated_cost_tier}). Pesanan berprioritas \code{CRITICAL/HIGH} memperoleh rekomendasi tervalidasi skema Pydantic 15 field agar anggaran subsidi tidak terbuang untuk intervensi pesanan bernilai kecil yang biaya kompensasinya melebihi margin.
\end{itemize}

\subsection*{Q18. Mengapa status proyek masih L2 dan bagaimana langkah persis agar lolos klaim kausal L3?}
\begin{itemize}
  \item \textbf{Jawaban Singkat Eksekutif:} Karena dataset Olist adalah catatan masa lalu yang statis---kita belum pernah benar-benar menelepon pembelinya secara langsung di aplikasi toko. Untuk membuktikan dampak nyata (L3), sudah disiapkan rencana eksperimen acak 80/20 saat terpasang di toko sungguhan.
  \item \textbf{Jawaban Mendalam DS:} Skenario S0--S5 (VAR Ridge 87 minggu) dan prioritas Decision Engine berada pada level \textbf{L2 (Simulasi Terkalibrasi Historis)}. Untuk meraih klaim kausal \textbf{L3}, CSIP mendefinisikan protokol eksperimen lapangan \textbf{80/20 Holdout}: 80\% pesanan berisiko di zona target menerima intervensi proaktif dan 20\% ditahan secara acak sebagai grup kontrol sejati, lalu dievaluasi beda tingkat ulasan negatif aktual serta margin bersih setelah satu siklus operasional.
\end{itemize}

\subsection*{Q19. Jelaskan mekanisme deteksi drift di \texttt{DriftEngine}---kapan sebuah model harus diganti?}
\begin{itemize}
  \item \textbf{Jawaban Singkat Eksekutif:} Seperti cek kesehatan rutin kendaraan: CSIP memantau apakah profil paket masuk mulai bergeser dari zaman belajar dulu, apakah tebakan risikonya terasa janggal, serta apakah tebakannya mulai meleset setelah jawaban pelanggan datang. Jika ya, model penantang diuji coba lebih dulu sebelum boleh menggantikan model juara.
  \item \textbf{Jawaban Mendalam DS:} \code{DriftEngine} memantau tiga lapis: \emph{Covariate Shift} via PSI (\code{PSI_THRESHOLD_STABLE = 0,10}), Kolmogorov-Smirnov (KS), dan Wasserstein Distance; \emph{Prediction Drift} via Jensen-Shannon Divergence (JSD); serta \emph{Concept Decay} via penurunan PR-AUC, F1, dan kenaikan Brier. Modul \code{SlidingWindowRetrainer} \& \code{ModelLifecycleManager} kemudian mensyaratkan Challenger mengalahkan Champion dengan kenaikan PR-AUC minimal $+2\%$ trigger / $+1\%$ lift promosi tanpa memperburuk kalibrasi Brier lebih dari $\le 5\%$, dikunci hash SHA-256 demi kapabilitas rollback instan.
\end{itemize}

\subsection*{Q20. Bagaimana platform menjaga keamanan terhadap unggahan CSV berbahaya, injeksi prompt, atau kerusakan jalur direktori?}
\begin{itemize}
  \item \textbf{Jawaban Singkat Eksekutif:} Seluruh file yang diunggah diperlakukan sebagai masukan yang tidak langsung dipercaya: dicek ukuran, ekstensi, dan tipe kolomnya, diproses lokal di memori tanpa mengeksekusi perintah sistem dari isi teks file, dan tanpa melempar teks ke model LLM eksternal.
  \item \textbf{Jawaban Mendalam DS:} Pemrosesan unggahan CSV di CSIP menerapkan validasi skema, batas ukuran sampel \code{sample_cap = 3500} terhadap risiko kehabisan RAM, sanitasi penamaan string di \code{semantic_profiler.py}, tidak pernah melewatkan isi sel CSV ke fungsi \code{eval()}/\code{exec()} ataupun antarmuka shell OS, serta melokalisasi narasi di mesin deterministik bebas prompt-injection eksternal.
\end{itemize}

\subsection*{Q21. Apa bukti konkret bahwa platform ini benar-benar jalan ringan di laptop RAM 8 GB tanpa GPU?}
\begin{itemize}
  \item \textbf{Jawaban Singkat Eksekutif:} Seluruh kode memilih algoritma CPU yang sangat hemat memori seperti LightGBM berbasis histogram dan Ridge, membatasi sampel CSV berlebih pada 3.500 baris, dan berjalan satu layanan bersih tanpa menumpuk model raksasa miliaran parameter di memori.
  \item \textbf{Jawaban Mendalam DS:} Pelatihan baseline B0--B3, jaringan saraf ringan \code{TabularEmbeddingMLP} (\code{[128,64]}), TreeSHAP ($N=800$), hingga layanan FastAPI \& Streamlit berjalan mulus di arsitektur CPU-only laptop 8 GB. Berkas rilis utuhnya pun terkendali di bawah batas paket aman ($< 100$ MB via \code{package_release.py}).
\end{itemize}

\subsection*{Q22. Mengapa arsitektur Neural Network (\texttt{TabularEmbeddingMLP} \& \texttt{TextBiLSTM}) tidak dipilih menjadi Champion utama?}
\begin{itemize}
  \item \textbf{Jawaban Singkat Eksekutif:} Pada tabel transaksi terstruktur seperti logistik dan ongkir, rakitan pohon keputusan modern seperti LightGBM jauh lebih akurat, jauh lebih hemat memori RAM, dan jauh lebih mudah dijelaskan alasannya ke manajemen dibanding jaringan saraf tiruan yang rumit.
  \item \textbf{Jawaban Mendalam DS:} Eksperimen B4 tetap dibangun di CPU PyTorch (\code{deep_tabular.py} dan notebook 08) sebagai bagian komparasi tangga ilmiah serta ekstraksi teks ulasan Portugis. Namun di data tabular berukuran menengah ($\sim$95 ribu baris) dengan 15 fitur numerik/rekayasa, ensembel pohon histogram LightGBM memberikan keseimbangan terbaik antara PR-AUC (\code{0,3061}), kecepatan inferensi CPU, dan keterjelasan deterministik TreeSHAP.
\end{itemize}

\subsection*{Q23. Bagaimana penguji eksternal yakin bahwa 216 tes dan semua laporan di PDF ini tidak dikarang manual?}
\begin{itemize}
  \item \textbf{Jawaban Singkat Eksekutif:} Penguji dapat mengetik tepat satu baris perintah di terminal (\code{verify_release.py}). Komputer sendiri yang akan mengecek ulang semua bukti, menjalankan tes, memindai kalimat terlarang, mencocokkan sidik jari kriptografis file, dan menjawab PASS atau FAIL.
  \item \textbf{Jawaban Mendalam DS:} Repositori menanamkan \code{tools/check_claims.py} (pemindai regex anti-sesumbar untuk G3), \code{scripts/generate_sha256sums.py}, \code{reports/testing/pytest_actual.txt} (\code{216 passed} di 39 file uji), dan \code{tools/verify_release.py} (pengevaluasi 12 gerbang G1--G12 yang menuliskan bukti ke \code{reports/verification/gates.json}). Ketidakcocokan sekecil apa pun antara klaim dan artefak langsung membuat gerbang berstatus FAIL.
\end{itemize}

\subsection*{Q24. Untuk uji kompetensi SKKNI No. 299/2020 (KKNI Level 7), mengapa tidak seluruh 63 KUK diklaim berstatus \texttt{VERIFIED}?}
\begin{itemize}
  \item \textbf{Jawaban Singkat Eksekutif:} Karena aturan sertifikasi jujur! Kode program hanya bisa membuktikan kemampuan teknis, tetapi wawancara langsung di hadapan penguji dan tanda tangan persetujuan anggaran dari klien tidak boleh dipalsukan oleh skrip komputer.
  \item \textbf{Jawaban Mendalam DS:} Penyelarasan gerbang \textbf{G12} di \code{SKKNI_TRACEABILITY_GUIDE.md} membagi jujur 63 KUK lintas 11 unit menjadi: \textbf{35 KUK (55,6\%)} \code{VERIFIED} terikat tes+SHA-256, \textbf{12 KUK (19,0\%)} \code{ARTIFACT_ONLY}, \textbf{13 KUK (20,6\%)} \code{CANDIDATE_DEMONSTRATION_REQUIRED} (membutuhkan demo langsung seperti di Bab 15 buku ini!), dan \textbf{3 KUK (4,8\%)} \code{EXTERNAL_CONFIRMATION_REQUIRED} pada Unit 1 yang mewajibkan otorisasi eksternal formal.
\end{itemize}

\subsection*{Q25. Jika besok CSIP dilanjutkan ke fase versi 2.0 produksi penuh di perusahaan nyata, apa tiga prioritas rekayasa paling bernilai tambah?}
\begin{itemize}
  \item \textbf{Jawaban Singkat Eksekutif:} Pertama, menjalankan uji acak A/B 80/20 selama satu bulan untuk mengunci pembuktian dampak uang nyata (Level L3). Kedua, menggabungkan model regresi keterlambatan \code{Ridge} yang terbukti unggul di uji akhir. Ketiga, menyambungkan sinyal perpesanan penjual-pembeli dan stok barang agar keluhan pada paket tepat waktu ikut tertangkap.
  \item \textbf{Jawaban Mendalam DS:} Prioritas v2.0 berbasis bukti audit adalah: (1) mengoperasionalkan A/B 80/20 lewat \code{/decisions/recommend} guna transisi kausal L2 $\rightarrow$ L3; (2) mempromosikan \code{B1_ridge} atau hibrida linear-boosting terkalibrasi geografis (\code{RJ} safety-buffer) untuk menekan MAE delay di luar Sao Paulo; dan (3) memperkaya fitur produk, riwayat rating SKU penjual, serta percakapan pra-kirim guna mengatasi rendahnya sinyal checkout $T_0$ (\code{0,1335}) maupun slice \emph{On-Time} (\code{0,6522}).
\end{itemize}

\backmatter
\chapter*{Catatan Penutup \& Keutuhan Dokumen Induk}
Seluruh deskripsi arsitektur, angka metrik leaderboard, parameter algoritma, hasil uji hipotesis, batas kalibrasi, pemetaan SKKNI, dan naskah presentasi di dalam buku ini diturunkan langsung dari kode aktif dan artefak teraudit pada repositori \textbf{CSIP v1.1.0}. Tidak ada angka metrik fiktif yang dibuat untuk memperindah laporan. Gunakan berkas sumber LaTeX \code{docs/BUKU_INDUK_CSIP_LENGKAP.tex} serta skrip kompilator resmi \code{scripts/build_csip_book.py} setiap kali terjadi pembaruan artefak proyek agar dokumentasi selalu selaras satu kebenaran dengan kenyataan kode.

\end{document}
